% L’agriculture traditionnelle : de grandes améliorations — Géographie Tle A — Cours signé OnBuch+
% Programme officiel MINESEC. Compiler avec : tectonic N04-agriculture-traditionnelle-grandes-ameliorations.tex
\newcommand{\DOCMATIERE}{Géographie}
\newcommand{\DOCNIVEAU}{Tle A}
\newcommand{\DOCMODULE}{Module 1 : Le Cameroun}
\newcommand{\DOCLECON}{N04}
\newcommand{\DOCTITRE}{L’agriculture traditionnelle : de grandes améliorations}
% OnBuch+ — préambule « cours signés » ENRICHI (compilé avec tectonic / XeLaTeX)
% Système visuel « Pop » d'OnBuch+ : fond crème (jamais blanc), bordures encre
% épaisses, ombres « dures » décalées, titres Archivo Black en capitales.
% Version MATHÉMATIQUES : graphes (pgfplots/tikz), boîtes pédagogiques
% (définition, propriété, méthode, attention, le savais-tu, exercice résolu,
% exercice, TP, bilan), page de garde et signature OnBuch+ ancrée en bas.
%
% Macros attendues dans main.tex AVANT \input{preamble} :
%   \DOCMATIERE  \DOCNIVEAU  \DOCMODULE  \DOCLECON (numéro)  \DOCTITRE
\documentclass[11pt]{article}

\usepackage{fontspec}
\usepackage[french]{babel}
\usepackage[a4paper,margin=2cm,top=3.4cm,bottom=2.8cm,headheight=1.4cm,footskip=1.3cm]{geometry}
\usepackage{amsmath,amssymb,amsfonts}
\usepackage{mathtools} % \xmapsto, \coloneqq... souvent utilisés par les modèles
\usepackage{cancel} % simplifications barrées (bilans de forces)
% \abs{z} (module, valeur absolue) et \norm{u} (norme), utilisés par les modèles
\providecommand{\abs}{}\let\abs\relax\DeclarePairedDelimiter{\abs}{\lvert}{\rvert}
\providecommand{\norm}{}\let\norm\relax\DeclarePairedDelimiter{\norm}{\lVert}{\rVert}
\usepackage[table]{xcolor} % [table] : fournit aussi \rowcolors (alternance de lignes)
\usepackage{pagecolor}
\usepackage{tikz}
\usetikzlibrary{arrows.meta,positioning,calc,shapes.geometric,shapes.misc,shapes.arrows,decorations.pathmorphing,decorations.pathreplacing,decorations.markings,patterns,shadows,mindmap,trees,fit,backgrounds,matrix,intersections,angles,quotes}
\usepackage{pgfplots}
\usepackage[version=4]{mhchem} % équations nucléaires \ce{^{235}_{92}U}
\usepackage[european,siunitx]{circuitikz} % schémas électriques
\pgfplotsset{compat=1.18}
\usepgfplotslibrary{fillbetween,groupplots} % aires entre courbes (calcul intégral)
\usepackage{enumitem}
\usepackage{fancyhdr}
% [most] charge toutes les bibliothèques tcolorbox (listings, minted, raster,
% documentation...) et gonfle inutilement la mémoire TeX sur un document long
% (~300 boîtes) : on ne charge que ce qui est réellement utilisé.
\usepackage[skins,breakable]{tcolorbox}
\usepackage{graphicx}
\graphicspath{{/home/user/t-t/geo-onbuch/images/}}
\usepackage{colortbl}
\usepackage{booktabs}
\usepackage{tabularx}
\usepackage{multirow}
\usepackage{diagbox} % case d'en-tête diagonale des tableaux à double entrée
% Colonnes extensibles centrée / gauche / droite (C, L, R), souvent utilisées par les modèles
\newcolumntype{C}{>{\centering\arraybackslash}X}
\newcolumntype{L}{>{\raggedright\arraybackslash}X}
\newcolumntype{R}{>{\raggedleft\arraybackslash}X}
\usepackage{array}
\usepackage{multicol}
\usepackage{siunitx}
% parse-numbers=false : accepte aussi \SI{2,5 \times 10^{-3}}{...}
\sisetup{locale=FR,output-decimal-marker={,},group-minimum-digits=4,parse-numbers=false}
\DeclareSIUnit\ohm{\ensuremath{\symup{\Omega}}} % U+2126 absent de la police
\DeclareSIUnit\division{div} % réglages d'oscilloscope (ms/div, V/div)
\DeclareSIUnit\tr{tr} % tours (vitesse de rotation, ex. \SI{1500}{\tr\per\minute})
\DeclareSIUnit\molar{M} % concentration molaire (ex. \SI{3}{\molar}, \SI{1}{\micro\molar})
\DeclareSIUnit\an{an} % année (le modèle confond parfois avec \year, un compteur TeX) — ex. \SI{5}{\kilo\meter\per\an}
\DeclareSIUnit\FCFA{FCFA} % franc CFA (ex. \SI{500}{\FCFA\per\kilo\gram})
\DeclareSIUnit\degreeBrix{\textdegree Brix} % teneur en sucre d'un sirop/jus (réfractométrie)
\DeclareSIUnit\month{mois} % le modèle utilise parfois \month, un compteur TeX, comme unité de durée
\DeclareSIUnit\microliter{\micro\liter} % le modèle écrit parfois \microliter en un seul token
\DeclareSIUnit\million{millions} % dénombrement (ex. \SI{15}{\million\per\mL} spermatozoïdes)
\usepackage{qrcode}
% \degree : commande fréquemment utilisée par les modèles pour un angle ou une
% température en dehors de \SI{}{} (non fournie par les paquets chargés).
\providecommand{\degree}{^{\circ}}
% \qed (fin de démonstration), utilisé par les modèles sans amsthm
\providecommand{\qed}{\hfill$\square$}
% \barre : double barre des valeurs interdites dans un tableau de variations (tkz-tab)
\providecommand{\barre}{\ensuremath{\|}}
% environnement proof (amsthm non chargé) utilisé par les modèles
\ifdefined\proof\else\newenvironment{proof}[1][Démonstration]{\par\noindent\textit{#1.}\ }{\qed\par}\fi
\usepackage{lastpage}
\usepackage{hyperref}
\hypersetup{hidelinks}

% ============================================================
% Palette « Pop » OnBuch+ — mêmes hex que la classe P (plus_ui.dart)
% ============================================================
\definecolor{poporange}{HTML}{F59321}
\definecolor{popdark}{HTML}{E07A0C}
\definecolor{popcream}{HTML}{FFF6E8}
\definecolor{popink}{HTML}{1C1714}
\definecolor{poppurple}{HTML}{7A5AE0}
\definecolor{popgreen}{HTML}{2E9E6A}
\definecolor{poppink}{HTML}{E0457B}
\definecolor{popblue}{HTML}{2F7DE1}
\definecolor{popgold}{HTML}{C98A12}
\definecolor{popmuted}{HTML}{8A7F76}
\definecolor{popline}{HTML}{EADFD2}
\definecolor{popgray}{HTML}{8A7F76}
\colorlet{popgrey}{popgray}
% Alias tolérants (le modèle nomme parfois les couleurs différemment)
\colorlet{popwhite}{white}
\colorlet{popblack}{popink}
\colorlet{popred}{poppink}
\colorlet{popyellow}{popgold}
% Teintes claires (fonds de tableaux/encadrés) — le modèle nomme parfois
% les couleurs avec un suffixe L (ex. popblueL) sans les avoir définies.
\colorlet{poporangeL}{poporange!20}
\colorlet{popdarkL}{popdark!20}
\colorlet{poppurpleL}{poppurple!20}
\colorlet{popgreenL}{popgreen!20}
\colorlet{poppinkL}{poppink!20}
\colorlet{popblueL}{popblue!20}
\colorlet{popgoldL}{popgold!20}
\colorlet{popredL}{popred!20}
\colorlet{popgrayL}{popgray!20}
% Teintes claires pour fonds de tableaux / zones
\colorlet{popblueL}{popblue!10!white}
\colorlet{poporangeL}{poporange!14!white}
\colorlet{popgreenL}{popgreen!12!white}
\colorlet{poppurpleL}{poppurple!12!white}
\colorlet{poppinkL}{poppink!12!white}
\colorlet{popgoldL}{popgold!14!white}

\pagecolor{popcream}

% ============================================================
% Typographie
% ============================================================
\defaultfontfeatures{Ligatures=TeX}
\setmainfont{PlusJakartaSans-Medium.ttf}[Path=../../../fonts/,BoldFont=PlusJakartaSans-Bold.ttf]
% Maths en sans-serif (Fira Math) avec chiffres et lettres droites de Plus
% Jakarta, pour que les formules se fondent dans le texte.
\usepackage[math-style=ISO,bold-style=ISO]{unicode-math}
\setmathfont{FiraMath-Regular.otf}[Path=../../../fonts/]
\setmathfont{PlusJakartaSans-Medium.ttf}[Path=../../../fonts/,range={up/num,up/{Latin,latin},\mathup}]
\usepackage{icomma} % virgule décimale sans espace en mode math
% Caractères mathématiques Unicode absents des polices de texte (Plus Jakarta,
% Archivo Black) : rendus par leur équivalent mathématique, en texte comme en
% formule — sinon ils s'affichent en carré vide (ex. « Arithmétique dans ℤ »).
\usepackage{newunicodechar}
\newunicodechar{ℝ}{\ensuremath{\mathbb{R}}}
\newunicodechar{ℤ}{\ensuremath{\mathbb{Z}}}
\newunicodechar{ℂ}{\ensuremath{\mathbb{C}}}
\newunicodechar{ℕ}{\ensuremath{\mathbb{N}}}
\newunicodechar{ℚ}{\ensuremath{\mathbb{Q}}}
\newunicodechar{𝔻}{\ensuremath{\mathbb{D}}}
\newunicodechar{α}{\ensuremath{\alpha}}
\newunicodechar{β}{\ensuremath{\beta}}
\newunicodechar{γ}{\ensuremath{\gamma}}
\newunicodechar{δ}{\ensuremath{\delta}}
\newunicodechar{ε}{\ensuremath{\varepsilon}}
\newunicodechar{θ}{\ensuremath{\theta}}
\newunicodechar{λ}{\ensuremath{\lambda}}
\newunicodechar{μ}{\ensuremath{\mu}}
\newunicodechar{σ}{\ensuremath{\sigma}}
\newunicodechar{τ}{\ensuremath{\tau}}
\newunicodechar{φ}{\ensuremath{\varphi}}
\newunicodechar{ψ}{\ensuremath{\psi}}
\newunicodechar{ω}{\ensuremath{\omega}}
\newunicodechar{Σ}{\ensuremath{\Sigma}}
% majuscules produites par \MakeUppercase dans les titres (α-aminés -> Α)
\newunicodechar{Α}{\ensuremath{\alpha}}
\newunicodechar{Β}{\ensuremath{\beta}}
\newunicodechar{Γ}{\ensuremath{\Gamma}}
\newunicodechar{Δ}{\ensuremath{\Delta}}
\newunicodechar{Ε}{\ensuremath{\varepsilon}}
\newunicodechar{Θ}{\ensuremath{\Theta}}
\newunicodechar{Λ}{\ensuremath{\Lambda}}
\newunicodechar{Μ}{\ensuremath{\mu}}
\newunicodechar{Τ}{\ensuremath{\tau}}
\newunicodechar{Φ}{\ensuremath{\Phi}}
\newunicodechar{Ψ}{\ensuremath{\Psi}}
\newunicodechar{Ω}{\ensuremath{\Omega}}
\newunicodechar{↦}{\ensuremath{\mapsto}}
\newunicodechar{∈}{\ensuremath{\in}}
\newunicodechar{∉}{\ensuremath{\notin}}
\newunicodechar{∧}{\ensuremath{\wedge}}
\newunicodechar{⇔}{\ensuremath{\Leftrightarrow}}
\newunicodechar{⟺}{\ensuremath{\Longleftrightarrow}}
\newunicodechar{⇒}{\ensuremath{\Rightarrow}}
\newunicodechar{⟹}{\ensuremath{\Longrightarrow}}
\newunicodechar{∘}{\ensuremath{\circ}}
\newunicodechar{∪}{\ensuremath{\cup}}
\newunicodechar{∩}{\ensuremath{\cap}}
\newunicodechar{≡}{\ensuremath{\equiv}}
\newunicodechar{⊂}{\ensuremath{\subset}}
\newunicodechar{⊥}{\ensuremath{\perp}}
\newunicodechar{∃}{\ensuremath{\exists}}
\newunicodechar{∀}{\ensuremath{\forall}}
\newunicodechar{∛}{\ensuremath{\sqrt[3]{\,}}}
\newunicodechar{∜}{\ensuremath{\sqrt[4]{\,}}}
\newunicodechar{⃗}{\ensuremath{{}^{\to}}}
\newunicodechar{ⁿ}{\ifmmode^{n}\else\textsuperscript{\ensuremath{n}}\fi}
\newunicodechar{ˣ}{\ifmmode^{x}\else\textsuperscript{\ensuremath{x}}\fi}
\newunicodechar{ᵇ}{\ifmmode^{b}\else\textsuperscript{\ensuremath{b}}\fi}
\newunicodechar{ʳ}{\ifmmode^{r}\else\textsuperscript{\ensuremath{r}}\fi}
\newunicodechar{ᵘ}{\ifmmode^{u}\else\textsuperscript{\ensuremath{u}}\fi}
\newunicodechar{ʸ}{\ifmmode^{y}\else\textsuperscript{\ensuremath{y}}\fi}
\newunicodechar{ᵃ}{\ifmmode^{a}\else\textsuperscript{\ensuremath{a}}\fi}
\newunicodechar{ᵉ}{\ifmmode^{e}\else\textsuperscript{\ensuremath{e}}\fi}
\newunicodechar{ᵏ}{\ifmmode^{k}\else\textsuperscript{\ensuremath{k}}\fi}
\newunicodechar{ᵖ}{\ifmmode^{p}\else\textsuperscript{\ensuremath{p}}\fi}
\newunicodechar{ᴺ}{\ifmmode^{N}\else\textsuperscript{\ensuremath{N}}\fi}
\newunicodechar{ᶜ}{\ifmmode^{c}\else\textsuperscript{\ensuremath{c}}\fi}
\newunicodechar{ʰ}{\ifmmode^{h}\else\textsuperscript{\ensuremath{h}}\fi}
\newunicodechar{ᵠ}{\ifmmode^{q}\else\textsuperscript{\ensuremath{q}}\fi}
\newunicodechar{ₙ}{\ifmmode_{n}\else\textsubscript{\ensuremath{n}}\fi}
\newunicodechar{ₐ}{\ifmmode_{a}\else\textsubscript{\ensuremath{a}}\fi}
\newunicodechar{ᵢ}{\ifmmode_{i}\else\textsubscript{\ensuremath{i}}\fi}
\newunicodechar{ᵣ}{\ifmmode_{r}\else\textsubscript{\ensuremath{r}}\fi}
\newunicodechar{ₖ}{\ifmmode_{k}\else\textsubscript{\ensuremath{k}}\fi}
\newunicodechar{ᵦ}{\ifmmode_{\beta}\else\textsubscript{\ensuremath{\beta}}\fi}
\newunicodechar{ₓ}{\ifmmode_{x}\else\textsubscript{\ensuremath{x}}\fi}
\newfontfamily\popdisplay{ArchivoBlack-Regular.ttf}[Path=../../../fonts/]
\newfontfamily\poplogo{SpaceGrotesk-Bold.ttf}[Path=../../../fonts/]
\newfontfamily\popsemi{PlusJakartaSans-SemiBold.ttf}[Path=../../../fonts/]
\newfontfamily\popxbold{PlusJakartaSans-ExtraBold.ttf}[Path=../../../fonts/]
\newfontfamily\popxboldls{PlusJakartaSans-ExtraBold.ttf}[Path=../../../fonts/,LetterSpace=3.0]

\setlist[enumerate]{leftmargin=1.7em,itemsep=5pt,topsep=4pt}
\setlist[itemize]{leftmargin=1.7em,itemsep=4pt,topsep=4pt,label={\color{poporange}\textbullet}}
\setlength{\parindent}{0pt}
\setlength{\parskip}{5pt}
\renewcommand{\arraystretch}{1.25}

% pgfplots : style Pop par défaut
\pgfplotsset{
  every axis/.append style={axis line style={popink,line width=1pt},tick style={popink},
    grid style={popline,line width=0.5pt},label style={font=\small},tick label style={font=\footnotesize}},
  popcurve/.style={poporange,line width=1.8pt,no markers,smooth},
}
% Même style disponible dans un \draw TikZ simple (hors pgfplots)
\tikzset{popcurve/.style={poporange,line width=1.8pt}}
\tikzset{popgrid/.style={popline,very thin}}
\colorlet{popcurve}{poporange} % le nom du style est parfois utilisé comme couleur

% ============================================================
% Pill / squiggle
% ============================================================
\newcommand{\poppill}[3]{%
  \tikz[baseline=-0.55ex]\node[draw=popink,line width=1.2pt,rounded corners=2.6pt,%
    fill=#2,inner xsep=6.5pt,inner ysep=3pt]{%
    \popxboldls\fontsize{7}{7}\selectfont\color{#3}\MakeUppercase{#1}};%
}
\newcommand{\popsquiggle}[1]{%
  \tikz[baseline=-0.4ex]{
    \draw[#1,line width=2.6pt,line cap=round]
      plot [smooth] coordinates {(0,0.09) (0.34,0.01) (0.68,0.21) (1.02,0.01) (1.36,0.10)};
  }%
}
% Mot-clé surligné (marqueur orange sous le texte)
\newcommand{\cle}[1]{{\popxbold\color{popdark}#1}}

% ============================================================
% En-tête / pied
% ============================================================
\pagestyle{fancy}
\fancyhf{}
\renewcommand{\headrulewidth}{0pt}
\renewcommand{\footrulewidth}{0pt}
\lhead{\raisebox{-0.15ex}{\poplogo\fontsize{13}{13}\selectfont\color{popink}OnBuch\color{poporange}+}}
\rhead{\poppill{\DOCMATIERE\ \textbullet\ \DOCNIVEAU\ \textbullet\ Leçon \DOCLECON}{white}{popink}}
\lfoot{\popxbold\fontsize{7.6}{7.6}\selectfont\color{popmuted}\MakeUppercase{Cours signé OnBuch+}\ \textbullet\ {\popsemi\DOCTITRE}}
\rfoot{\tikz[baseline=-0.6ex]\node[rounded corners=8.5pt,fill=popink,minimum height=17pt,minimum width=17pt,inner xsep=5pt,inner sep=0]{\popxbold\fontsize{8}{8}\selectfont\color{popcream}\thepage\,/\,\pageref*{LastPage}};}

% ============================================================
% Titres
% ============================================================
\newcommand{\chaptitre}[2]{%
  \begin{tcolorbox}[enhanced,colback=poporange,colframe=popink,boxrule=2.6pt,arc=11pt,
      drop shadow=popink,left=15pt,right=15pt,top=12pt,bottom=11pt,before skip=3pt,after skip=18pt]
    \poppill{#2}{popink}{popcream}\par\vspace{9pt}
    {\raggedright\hyphenpenalty=10000\exhyphenpenalty=10000\popdisplay\fontsize{22}{24}\selectfont\color{popcream}\MakeUppercase{#1}\par}
  \end{tcolorbox}
}
% Section numérotée : pastille orange avec le numéro + titre Archivo
\newcounter{popsec}
\newcommand{\coursec}[1]{%
  \stepcounter{popsec}\setcounter{popsub}{0}%
  \par\vspace{16pt}\noindent
  \begin{minipage}{\linewidth}
  \tikz[baseline=-0.7ex]\node[draw=popink,line width=1.6pt,fill=poporange,rounded corners=4pt,minimum size=22pt,inner sep=0]{\popdisplay\fontsize{12}{12}\selectfont\color{popcream}\thepopsec};%
  \hspace{8pt}{\popdisplay\fontsize{15}{17}\selectfont\color{popink}\MakeUppercase{#1}}\par
  \vspace{4pt}\noindent{\color{poporange}\rule{36pt}{3.6pt}}
  \end{minipage}\par\nopagebreak\vspace{8pt}
}
\newcounter{popsub}
\newcommand{\courssub}[1]{%
  \stepcounter{popsub}%
  \par\vspace{9pt}\noindent{\popxbold\fontsize{11.5}{13}\selectfont\color{popink}{\color{poporange}\thepopsec.\thepopsub}\hspace{6pt}#1}\par\nopagebreak\vspace{4pt}
}

% ============================================================
% Boîtes pédagogiques — même recette : fond blanc, bordure épaisse colorée,
% ombre dure encre, badge plein. Titre optionnel [#1].
% ============================================================
\tcbset{popbase/.style={breakable,enhanced,colback=white,boxrule=2.3pt,arc=9pt,
  shadow={3pt}{-3pt}{0pt}{popink},left=14pt,right=14pt,top=11pt,bottom=11pt,before skip=11pt,after skip=15pt,
  fontupper=\popsemi\fontsize{10.6}{14.5}\selectfont\color{popink},
  fontlower=\popsemi\fontsize{10.6}{14.5}\selectfont\color{popink}}}
% Coche dessinée (le glyphe ✓ n'existe pas dans les polices du document)
\newcommand{\popcheck}[1]{\tikz[baseline=-0.5ex]\draw[#1,line width=1.6pt,line cap=round,line join=round](-0.13,0.0)--(-0.04,-0.09)--(0.13,0.1);}
\providecommand{\coche}{\popcheck{popgreen}}
\providecommand{\resultat}[1]{\par\smallskip\noindent\fcolorbox{popgreen}{popgreenL}{\parbox{\dimexpr\linewidth-2\fboxsep-2\fboxrule\relax}{\textbf{Résultat.} #1}}\par\smallskip}
\newcommand{\popboxhead}[3]{\poppill{#1}{#2}{white}\ifx\relax#3\relax\else\hspace{6pt}{\popxbold\fontsize{10.6}{12}\selectfont\color{popink}#3}\fi\par\vspace{7pt}}

\newtcolorbox{aretenir}[1][]{popbase,colframe=popgreen,before upper={\popboxhead{À retenir}{popgreen}{#1}}}
\newtcolorbox{exemplebox}[1][]{popbase,colframe=poporange,before upper={\popboxhead{Exemple}{poporange}{#1}}}
\newtcolorbox{definition}[1][]{popbase,colframe=popblue,before upper={\popboxhead{Définition}{popblue}{#1}}}
\newtcolorbox{propriete}[1][]{popbase,colframe=poppurple,before upper={\popboxhead{Propriété}{poppurple}{#1}}}
\newtcolorbox{methode}[1][]{popbase,colframe=popgold,before upper={\popboxhead{Méthode}{popgold}{#1}}}
\newtcolorbox{attention}[1][]{popbase,colframe=poppink,before upper={\popboxhead{Attention}{poppink}{#1}}}
\newtcolorbox{experience}[1][]{popbase,colframe=popblue,colback=popblueL,before upper={\popboxhead{Expérience}{popblue}{#1}}}
% « Le savais-tu ? » : carte violette pleine (contexte, culture, Cameroun)
\newtcolorbox{savaistu}[1][]{popbase,colback=poppurple,colframe=popink,
  fontupper=\popsemi\fontsize{10.4}{14.2}\selectfont\color{white},
  before upper={\poppill{Le savais-tu\,?}{white}{poppurple}\ifx\relax#1\relax\else\hspace{6pt}{\popxbold\color{white}#1}\fi\par\vspace{7pt}}}
% Exercice résolu : énoncé en haut, \tcblower puis la solution sur fond crème
\newcounter{popexo}\newcounter{popexr}
\newtcolorbox{exoresolu}[1][]{popbase,colframe=poporange,colbacklower=popcream,
  segmentation style={popink,line width=1.2pt,dashed},
  before upper={\stepcounter{popexr}\popboxhead{Exercice résolu \thepopexr}{poporange}{#1}},
  before lower={\poppill{Solution}{popink}{popcream}\par\vspace{6pt}}}
% Exercice (non résolu en place) — la correction est dans la partie Corrigés
\newtcolorbox{exercice}[1][]{popbase,colframe=popink,
  before upper={\stepcounter{popexo}\popboxhead{Exercice \thepopexo}{popink}{#1}}}
% Corrigé
\newtcolorbox{corrige}[1][]{popbase,colframe=popgreen,colback=popgreenL,
  before upper={\popboxhead{Corrigé}{popgreen}{#1}}}
% Objectifs / prérequis / bilan
\newtcolorbox{objectifs}{popbase,colframe=popink,colback=poporangeL,
  before upper={\popboxhead{Objectifs de la leçon}{popink}{}}}
\newtcolorbox{prerequis}{popbase,colframe=popink,colback=popblueL,
  before upper={\popboxhead{Prérequis}{popblue}{}}}
\newtcolorbox{bilan}{popbase,colframe=popink,colback=white,boxrule=2.8pt,
  before upper={\popboxhead{Fiche bilan}{poporange}{L'essentiel à connaître pour l'examen}}}
% Figure encadrée avec légende
\newtcolorbox{popfigure}[1][]{enhanced,breakable=false,colback=white,colframe=popline,boxrule=1.4pt,arc=8pt,
  left=10pt,right=10pt,top=10pt,bottom=8pt,before skip=10pt,after skip=12pt,halign=center,
  lower separated=false,halign lower=center,
  fontlower=\popsemi\fontsize{9}{11.5}\selectfont\color{popmuted},
  #1}
\newcommand{\legende}[1]{\par\vspace{6pt}{\popsemi\fontsize{9}{11.5}\selectfont\color{popmuted}#1}}

% ============================================================
% Signature OnBuch+
% ============================================================
\newcommand{\poplogoinline}[1]{{\poplogo\fontsize{#1}{#1}\selectfont\color{popink}OnBuch\color{poporange}+}}
\newcommand{\signatureonbuch}{%
  \begin{tcolorbox}[enhanced,colback=popink,colframe=popink,boxrule=2.2pt,arc=10pt,
      drop shadow=poporange,left=14pt,right=14pt,top=10pt,bottom=10pt,before skip=0pt,after skip=0pt]
    \tikz[baseline=-0.7ex]\node[circle,fill=popgreen,draw=popcream,line width=1.3pt,minimum size=22pt,inner sep=0]{\popcheck{white}};%
    \hspace{9pt}\begin{minipage}[c]{0.62\linewidth}
      {\popxbold\fontsize{12}{13}\selectfont\color{popcream}Signé\ \textbullet\ {\poplogo OnBuch\color{poporange}+}}\par\vspace{2pt}
      {\raggedright\popsemi\fontsize{8}{10.5}\selectfont\color{popline}\MakeUppercase{Équipe pédagogique OnBuch+}\\\MakeUppercase{Conforme au programme officiel MINESEC}\par}
    \end{minipage}\hfill
    {\poplogo\fontsize{22}{22}\selectfont\color{popcream}OnBuch\color{poporange}+}
  \end{tcolorbox}%
}

% ============================================================
% Page de garde
% #1 titre · #2 matière · #3 niveau · #4 module · #5 n° leçon · #6 durée ·
% #7 description / objectifs (texte ou itemize)
% ============================================================
\newcommand{\pagegarde}[7]{%
  \thispagestyle{empty}
  \newgeometry{a4paper,margin=1.7cm,top=1.6cm,bottom=1.4cm}%
  \begin{tikzpicture}[remember picture,overlay]
    \fill[poporange!18] ([xshift=-3.2cm,yshift=-3.6cm]current page.north east) circle (4.6cm);
    \fill[poppurple!14] ([xshift=2.4cm,yshift=7.6cm]current page.south west) circle (3.8cm);
  \end{tikzpicture}%
  {\poplogo\fontsize{30}{30}\selectfont\color{popink}OnBuch\color{poporange}+}\hfill
  \raisebox{0.6ex}{\poppill{Cours signé \textbullet\ Premium}{popink}{popcream}}\par
  \vspace{1pt}\popsquiggle{poppurple}\par
  \vspace{8pt}
  \begin{tcolorbox}[enhanced,colback=poporange,colframe=popink,boxrule=2.8pt,arc=13pt,
      drop shadow=popink,left=18pt,right=18pt,top=12pt,bottom=13pt,after skip=10pt]
    \poppill{#2 \textbullet\ #3}{popink}{popcream}\hspace{5pt}\poppill{#4}{white}{popink}\par\vspace{8pt}
    {\popdisplay\fontsize{14}{14}\selectfont\color{popink}LEÇON #5}\par\vspace{4pt}
    {\raggedright\hyphenpenalty=10000\exhyphenpenalty=10000\popdisplay\fontsize{25}{27}\selectfont\color{popcream}\MakeUppercase{#1}\par}
    \vspace{8pt}
    {\popxbold\fontsize{9.5}{9.5}\selectfont\color{popink}\MakeUppercase{Durée conseillée : #6}}
  \end{tcolorbox}
  \begin{tcolorbox}[enhanced,colback=white,colframe=popink,boxrule=2.2pt,arc=10pt,
      drop shadow=popink,left=16pt,right=16pt,top=10pt,bottom=10pt,after skip=8pt]
    \poppill{Dans cette leçon}{poppurple}{white}\par\vspace{5pt}
    \popsemi\fontsize{9.3}{12.6}\selectfont\color{popink}#7
  \end{tcolorbox}
  \noindent\begin{minipage}[t]{0.487\linewidth}
    \begin{tcolorbox}[enhanced,colback=popgreen,colframe=popink,boxrule=2.2pt,arc=10pt,
        drop shadow=popink,left=11pt,right=11pt,top=10pt,bottom=10pt,height=3.1cm,valign=center]
      \tcbox[colback=white,colframe=popink,boxrule=1.3pt,arc=5pt,boxsep=0pt,nobeforeafter,
        left=3pt,right=3pt,top=3pt,bottom=3pt,valign=center]{\qrcode[height=1.9cm]{https://play.google.com/store/apps/details?id=com.luvvix.onbuch&hl=fr}}%
      \hspace{8pt}\begin{minipage}[c]{0.5\linewidth}
        {\popdisplay\fontsize{11}{12.5}\selectfont\color{white}\MakeUppercase{Télécharge OnBuch}}\par\vspace{4pt}
        {\raggedright\popsemi\fontsize{8.4}{11}\selectfont\color{white}Gratuit \textbullet\ Google Play\\Tuteur IA, annales, concours, résultats\par}
      \end{minipage}
    \end{tcolorbox}
  \end{minipage}\hfill
  \begin{minipage}[t]{0.487\linewidth}
    \begin{tcolorbox}[enhanced,colback=poppurple,colframe=popink,boxrule=2.2pt,arc=10pt,
        drop shadow=popink,left=11pt,right=11pt,top=10pt,bottom=10pt,height=3.1cm,valign=center]
      \tikz[baseline=-0.7ex]\node[circle,fill=white,draw=popink,line width=1.3pt,minimum size=1.6cm,inner sep=0]{\popdisplay\fontsize{24}{24}\selectfont\color{poppurple}+};%
      \hspace{8pt}\begin{minipage}[c]{0.6\linewidth}
        {\popdisplay\fontsize{11}{12.5}\selectfont\color{white}\MakeUppercase{Passe à OnBuch+}}\par\vspace{4pt}
        {\raggedright\popsemi\fontsize{8.4}{11}\selectfont\color{white}Tous les cours signés, Tuteur IA illimité, fiches PDF — onglet OnBuch+ de l'app\par}
      \end{minipage}
    \end{tcolorbox}
  \end{minipage}
  \vspace{8pt}
  \popcontenu
  \vfill
  \signatureonbuch
  \restoregeometry
  \newpage
}

% Rangée « Ce cours contient » (page de garde)
\newcommand{\poptile}[3]{% #1 couleur #2 gros texte #3 légende
  \begin{tcolorbox}[enhanced,colback=#1,colframe=popink,boxrule=2pt,arc=9pt,shadow={2.5pt}{-2.5pt}{0pt}{popink},
      left=6pt,right=6pt,top=9pt,bottom=9pt,halign=center,valign=center,height=1.75cm,width=\linewidth,nobeforeafter]
    {\popdisplay\fontsize{12.5}{13}\selectfont\color{white}\MakeUppercase{#2}}\par\vspace{3pt}
    {\centering\popsemi\fontsize{7.8}{9.5}\selectfont\color{white}#3\par}
  \end{tcolorbox}}
\newcommand{\popcontenu}{%
  \par\noindent{\popxboldls\fontsize{7.5}{7.5}\selectfont\color{popmuted}CE COURS SIGNÉ CONTIENT}\par\vspace{7pt}
  \noindent\begin{minipage}[t]{0.235\linewidth}\poptile{popblue}{Cours}{complet, illustré, pas à pas}\end{minipage}\hfill
  \begin{minipage}[t]{0.235\linewidth}\poptile{poporange}{Méthodes}{et exercices résolus}\end{minipage}\hfill
  \begin{minipage}[t]{0.235\linewidth}\poptile{poppink}{Exercices}{type examen + corrigés}\end{minipage}\hfill
  \begin{minipage}[t]{0.235\linewidth}\poptile{popgreen}{Fiche bilan}{l'essentiel à retenir}\end{minipage}\par}

% Sommaire
\newtcolorbox{sommaire}{enhanced,colback=white,colframe=popink,boxrule=2.2pt,arc=10pt,
  drop shadow=popink,left=18pt,right=18pt,top=13pt,bottom=13pt,before skip=6pt,after skip=20pt,
  before upper={\poppill{Sommaire}{poppurple}{white}\par\vspace{9pt}\popsemi\fontsize{10.6}{14.5}\selectfont\color{popink}}}

% Signature de fin de cours — poussée tout en bas de la dernière page
\newcommand{\findecours}{%
  \par\vfill\nopagebreak
  \begin{center}\popsquiggle{poporange}\end{center}\vspace{4pt}
  \signatureonbuch
}
\AtBeginDocument{\def\llbracket{\mathopen{[\mkern-3.5mu[}}\def\rrbracket{\mathclose{]\mkern-3.5mu]}}} % intervalles d'entiers (glyphe ⟦ absent de la police)
% Glyphes absents de Fira Math : redéfinis à partir de glyphes disponibles
\AtBeginDocument{%
  \def\setminus{\mathbin{\backslash}}\let\smallsetminus\setminus
  \def\vdots{\mathord{\vbox{\baselineskip3.2pt\lineskiplimit0pt\kern4pt\hbox{.}\hbox{.}\hbox{.}}}}%
  \def\ddots{\mathinner{\mkern1mu\raise6.5pt\hbox{.}\mkern2mu\raise3.6pt\hbox{.}\mkern2mu\raise0.7pt\hbox{.}\mkern1mu}}%
  \def\triangle{\mathord{\scalebox{0.9}{$\symup{\Delta}$}}}%
  \def\nabla{\mathord{\raisebox{\depth}{\scalebox{1}[-1]{$\symup{\Delta}$}}}}%
  \def\checkmark{\popcheck{popdark}}%
  \def\star{\mathbin{\ast}}\def\bigstar{\mathord{\ast}}%
  \def\diamond{\mathbin{\scalebox{0.6}{$\Diamond$}}}%
  \def\bigcup{\mathop{\mathchoice{\raisebox{-0.3em}{\scalebox{1.7}{$\cup$}}}{\scalebox{1.25}{$\cup$}}{\cup}{\cup}}\displaylimits}%
  \def\bigcap{\mathop{\mathchoice{\raisebox{-0.3em}{\scalebox{1.7}{$\cap$}}}{\scalebox{1.25}{$\cap$}}{\cap}{\cap}}\displaylimits}%
  \def\lesseqqgtr{\mathrel{\lessgtr}}\def\bigoplus{\mathop{\oplus}\displaylimits}%
}
\providecommand{\ssi}{si et seulement si} % abréviation fréquente
\AtBeginDocument{\providecommand{\wideparen}{\overparen}} % arc de cercle

%% FIGFIX-BEGIN v5 — correctifs globaux des figures (docs/onbuch-generation/tools/figfix)
\usepackage{adjustbox}\usepackage{etoolbox}\tcbuselibrary{raster}
% toute figure TikZ/pgfplots plus large que la ligne est réduite à la largeur disponible
\BeforeBeginEnvironment{tikzpicture}{\begin{adjustbox}{max width=\linewidth}}
\AfterEndEnvironment{tikzpicture}{\end{adjustbox}}
% légendes pgfplots lisibles par-dessus les courbes
\pgfplotsset{every axis/.append style={legend style={fill=white,fill opacity=0.92,text opacity=1,draw=black!15,inner sep=3pt}}}
% étiquettes d'arêtes (syntaxe edge["..."]) avec fond blanc
\tikzset{every edge quotes/.append style={fill=white,inner sep=1.2pt}}
%% FIGFIX-END
\begin{document}

\pagegarde{L’agriculture traditionnelle : de grandes améliorations}{Géographie}{Tle A}{Module 1 : Le Cameroun}{N04}{2 h}{Cette leçon étudie les transformations de l'agriculture traditionnelle camerounaise : amélioration de l'outillage, évolution des techniques culturales et adoption de nouvelles cultures. Elle analyse ensuite les formes d'agriculture (itinérante sur brûlis, vivrière, de seconde génération) à partir de cartes, photographies et données statistiques du Cameroun.
\vspace{4pt}
\begin{multicols}{2}\small
\begin{itemize}
\item À la fin de cette leçon, je sais définir l'agriculture traditionnelle et identifier ses caractéristiques au Cameroun
\item À la fin de cette leçon, je sais décrire les améliorations de l'outillage agricole (houe, machette, traction animale, motoculteur)
\item À la fin de cette leçon, je sais expliquer l'évolution des techniques culturales (jachère raccourcie, engrais, sélection des semences, irrigation)
\item À la fin de cette leçon, je sais citer et localiser les nouvelles cultures adoptées au Cameroun (maïs, manioc amélioré, palmier à huile, riz)
\item À la fin de cette leçon, je sais décrire le fonctionnement et les limites de l'agriculture itinérante sur brûlis
\item À la fin de cette leçon, je sais distinguer agriculture vivrière et agriculture de seconde génération à partir d'exemples camerounais
\end{itemize}\end{multicols}}

\begin{sommaire}
\begin{multicols}{2}
\begin{enumerate}
\item L'agriculture traditionnelle camerounaise : cadre et caractéristiques
\item Un outillage amélioré
\item L'évolution des techniques culturales
\item L'adoption de nouvelles cultures
\item L'agriculture itinérante sur brûlis et l'agriculture vivrière
\item L'agriculture de seconde génération et bilan des améliorations
\item Activité d'intégration
\item Méthodes et exercices résolus
\item Exercices
\item Corrigés des exercices
\item Fiche bilan
\end{enumerate}
\end{multicols}
\end{sommaire}

\begin{objectifs}
\begin{itemize}
\item À la fin de cette leçon, je sais définir l'agriculture traditionnelle et identifier ses caractéristiques au Cameroun
\item À la fin de cette leçon, je sais décrire les améliorations de l'outillage agricole (houe, machette, traction animale, motoculteur)
\item À la fin de cette leçon, je sais expliquer l'évolution des techniques culturales (jachère raccourcie, engrais, sélection des semences, irrigation)
\item À la fin de cette leçon, je sais citer et localiser les nouvelles cultures adoptées au Cameroun (maïs, manioc amélioré, palmier à huile, riz)
\item À la fin de cette leçon, je sais décrire le fonctionnement et les limites de l'agriculture itinérante sur brûlis
\item À la fin de cette leçon, je sais distinguer agriculture vivrière et agriculture de seconde génération à partir d'exemples camerounais
\item À la fin de cette leçon, je sais légender une carte de répartition des systèmes agricoles du Cameroun
\item À la fin de cette leçon, je sais construire un croquis et un diagramme (histogramme, courbe) à partir de données agricoles
\item À la fin de cette leçon, je sais inventorier, classer et interpréter des documents (cartes, photographies, tableaux statistiques) sur l'agriculture camerounaise
\end{itemize}
\end{objectifs}

\begin{prerequis}
\begin{itemize}
\item Les grandes zones agro-écologiques du Cameroun (soudanienne, soudano-sahélienne, forestière, hautes terres) vues en classe de Première
\item La notion de milieu naturel et de contraintes climatiques (pluviométrie, saisons)
\item La répartition de la population camerounaise et la pression démographique sur les terres
\item Les notions de base de lecture de carte (échelle, légende, orientation)
\item La place de l'agriculture dans l'économie camerounaise (part du PIB, population active agricole)
\end{itemize}
\end{prerequis}

\newpage

\coursec{L'agriculture traditionnelle camerounaise : cadre et caractéristiques}

\textbf{Situation d'accroche.} À Foumbot, dans les Hauts-Plateaux de l'Ouest, un jeune planteur utilise encore la houe de son grand-père, mais il seme désormais du \cle{maïs hybride}, achète des engrais à crédit et vend une partie de ses récoltes à la coopérative agricole du village. À quelques centaines de kilomètres de là, dans la forêt du Sud, une famille pratique toujours l'agriculture itinérante sur \cle{brûlis} comme il y a un siècle : elle défriche, brûle, sème, puis abandonne la parcelle après deux ou trois récoltes. Deux mondes agricoles, un seul pays.

\textbf{Problématique.} Comment expliquer que l'agriculture camerounaise, longtemps jugée « traditionnelle », ait connu de si grandes améliorations tout en restant dominée par les petits exploitants ? Cette leçon montre comment l'outillage, les techniques et les cultures ont évolué, et ce qui distingue l'agriculture vivrière de l'agriculture de seconde génération. Mais avant de mesurer ces progrès, il faut d'abord comprendre ce qu'est l'agriculture traditionnelle camerounaise, quel poids elle occupe dans la vie du pays, et sous quelles formes elle se présente d'une région à l'autre. C'est l'objet de cette première section.

\courssub{Qu'est-ce que l'agriculture traditionnelle ?}

\textbf{L'intuition d'abord.} Imagine le champ de ton grand-père ou de ton oncle au village : une parcelle d'un hectare ou deux, défrichée à la machette, travaillée à la houe, plantée de manioc, de macabo, de plantain et de maïs mélangés sur le même champ. La famille entière participe aux travaux ; l'essentiel de la récolte est consommé à la maison, et seul le surplus est vendu au marché hebdomadaire pour payer les frais scolaires ou le savon. Voilà, concrètement, ce qu'est l'agriculture traditionnelle.

\begin{definition}[Agriculture traditionnelle]
L'\cle{agriculture traditionnelle} est un système agricole fondé sur le \textbf{travail manuel familial}, utilisant un outillage simple (machette, houe, daba), caractérisé par de \textbf{faibles rendements} et une production destinée d'abord à l'\textbf{autoconsommation} (agriculture de subsistance), la vente ne concernant que les surplus. Elle se pratique sur de petites exploitations et repose sur des savoir-faire hérités et transmis de génération en génération.
\end{definition}

Trois traits structurels permettent de reconnaître ce système :
\begin{itemize}
\item la \cle{motorisation} est quasi absente : la force de travail est humaine, parfois animale (culture attelée au Nord) ;
\item les \cle{intrants} (engrais, pesticides, semences améliorées) sont peu ou pas utilisés, d'où des rendements faibles, souvent inférieurs à 1 tonne de maïs par hectare contre 5 à 6 tonnes dans une agriculture intensifiée ;
\item la \cle{main-d'œuvre} est essentiellement familiale : parents, enfants, parents élargis, parfois entraidés par les groupements de travail collectif (les \emph{djangui} de l'Ouest).
\end{itemize}

\begin{attention}[Ne confonds pas « traditionnelle » et « archaïque »]
L'erreur la plus fréquente au Bac consiste à écrire que l'agriculture traditionnelle est « figée », « primitive » ou « archaïque ». C'est faux. L'agriculture traditionnelle camerounaise \textbf{évolue} : elle adopte de nouvelles cultures, améliore son outillage, intègre des semences sélectionnées et vend de plus en plus sa production. Le titre même de la leçon — « de grandes améliorations » — doit te rappeler que le traditionnel n'est pas l'immobile. Dans une copie, parle toujours d'une agriculture traditionnelle \emph{en mutation}.
\end{attention}

\courssub{Le poids de l'agriculture dans la vie nationale}

Pourquoi consacrer une leçon entière à ce secteur ? Parce que l'agriculture est la colonne vertébrale de l'économie et de la société camerounaises.

\begin{exemplebox}[L'agriculture, premier employeur du Cameroun]
L'agriculture occupe environ \textbf{6 Camerounais actifs sur 10} (soit 60 à 70 \% de la population active, selon les sources et les années ; donne cet ordre de grandeur dans tes copies). Elle contribue pour environ \textbf{15 à 20 \% du PIB} et fournit l'essentiel des denrées alimentaires consommées dans le pays : le Cameroun assure ainsi une large autosuffisance alimentaire pour les produits de base (manioc, plantain, maïs, mil, arachide), même s'il importe encore du riz, du blé et du poisson. Elle fournit aussi des matières premières à l'industrie (coton pour la CICAM, sucre pour la SOSUCAM) et des produits d'exportation (cacao, café, coton, banane).
\end{exemplebox}

\begin{popfigure}
\begin{center}
\begin{tikzpicture}[scale=1.05]
% Secteurs : agriculture 62 %, industrie 10 %, services 28 %
\fill[popgreenL] (0,0) -- (90:2.2) arc (90:-133.2:2.2) -- cycle;
\fill[poporangeL] (0,0) -- (-133.2:2.2) arc (-133.2:-169.2:2.2) -- cycle;
\fill[popblueL] (0,0) -- (-169.2:2.2) arc (-169.2:-270:2.2) -- cycle;
\draw[popdark, thick] (0,0) circle (2.2);
\node[popdark, font=\small\bfseries] at (25:1.35) {Agriculture};
\node[popdark, font=\small] at (-15:1.35) {62 \%};
\node[popdark, font=\small\bfseries] at (-151:1.4) {Ind.};
\node[popdark, font=\small] at (-151:1.0) {10 \%};
\node[popdark, font=\small\bfseries] at (140:1.4) {Services};
\node[popdark, font=\small] at (140:1.0) {28 \%};
\end{tikzpicture}
\end{center}
\legende{Répartition approximative de la population active camerounaise par secteur d'activité (ordres de grandeur : agriculture environ 60--65 \%, industrie environ 10 \%, services environ 25--30 \%). Source : estimations INS/Banque mondiale, valeurs arrondies.}
\end{popfigure}

Ce diagramme montre l'essentiel : malgré l'urbanisation rapide et le développement des services, le Cameroun reste un pays où la majorité des bras travaille la terre. Comprendre l'agriculture traditionnelle, c'est donc comprendre la vie quotidienne de la majorité des Camerounais.

\courssub{De petites exploitations familiales}

L'unité de base de l'agriculture traditionnelle est l'\cle{exploitation familiale}. Sa caractéristique la plus frappante est sa petitesse : la grande majorité des exploitations camerounaises couvre \textbf{moins de 2 hectares}, souvent entre 0,5 et 1,5 hectare, parfois morcelée en plusieurs parcelles dispersées autour du village.

Cette petite taille a des conséquences directes :
\begin{itemize}
\item elle rend la mécanisation coûteuse et difficile à rentabiliser (un tracteur pour 1 hectare ?) ;
\item elle limite la production commercialisable, donc les revenus du planteur ;
\item elle explique la forte densité de main-d'œuvre par hectare : toute la famille travaille sur une surface réduite ;
\item elle fragilise les paysans face aux aléas climatiques : une mauvaise saison sur 1 hectare, et c'est toute la sécurité alimentaire du ménage qui vacille.
\end{itemize}

\begin{aretenir}[Le portrait-robot de l'exploitation traditionnelle]
Petite surface (généralement $<$ 2 ha), main-d'œuvre familiale, outillage manuel, faible capital, polyculture vivrière, autoconsommation dominante, vente des surplus. Ce portrait domine la campagne camerounaise, avec des variantes régionales fortes que nous allons maintenant cartographier.
\end{aretenir}

\courssub{Une diversité de systèmes agricoles selon les zones agro-écologiques}

Le Cameroun s'étire sur plus de \SI{1200}{km} du sud au nord, de la forêt équatoriale humide à la steppe soudanienne sèche. Cette diversité climatique et végétale commande une diversité de systèmes agricoles : on ne cultive pas de la même façon à Ebolowa, à Foumban et à Maroua.

\begin{popfigure}
\begin{center}
\includegraphics[width=\linewidth,height=9.5cm,keepaspectratio]{N04/cmr_koppen.png}
\end{center}
\legende{Climats du Cameroun selon Köppen (légende en anglais), fond pour lire les grandes zones agro-écologiques : climats tropicaux humides (Af, Am, bleus) au Sud et sur le littoral, où dominent forêt et cultures de rente et vivrières ; savane (Aw, bleu clair) au centre et au nord ; climats arides (BSh orange, BWh rouge) à l'Extrême-Nord ; taches jaunes (Csb) sur les hautes terres de l'Ouest. L'essentiel est la répartition nord-sud des systèmes agricoles, qui suit la baisse des pluies. \\ Source : Wikimedia Commons, Beck et al., CC BY 4.0.}
\end{popfigure}

Trois grands ensembles se dégagent :

\begin{enumerate}
\item \textbf{La zone forestière humide du Sud} (régions du Sud, du Centre, de l'Est, du Littoral et du Sud-Ouest) : les pluies sont abondantes (plus de \SI{1500}{mm} par an) et réparties sur deux saisons. Le système dominant y est l'\cle{agriculture itinérante sur brûlis} : on défriche une lisière de forêt, on brûle la végétation (les cendres fertilisent le sol), on cultive deux ou trois ans (manioc, plantain, macabo, igname, arachide), puis on abandonne la parcelle à la jachère forestière pendant dix à quinze ans. S'y ajoutent les cultures de rente pérennes : cacao, café, palmiers à huile.
\item \textbf{Les Hautes Terres de l'Ouest} (régions de l'Ouest et du Nord-Ouest) : l'altitude (1000 à \SI{2000}{m}) tempère le climat et les sols volcaniques sont fertiles. La densité de population y est très forte (plus de 100 habitants au km\textsuperscript{2} dans les pays bamiléké), ce qui a poussé vers des \cle{cultures intensives} : rotation des cultures, jachères raccourcies, utilisation du fumier, maraîchage (choux, carottes, pommes de terre à Santa), maïs et haricot en association.
\item \textbf{La zone soudanienne du Nord} (régions de l'Adamaoua, du Nord et de l'Extrême-Nord) : une longue saison sèche et une seule saison des pluies imposent des \cle{cultures sèches} (mil, sorgho, niébé, arachide, coton) et l'\cle{élevage} (bovins zébus, ovins, caprins), souvent associés dans un système agro-pastoral : les troupeaux fertilisent les champs après la récolte. La culture attelée (traction animale) y est plus répandue qu'ailleurs.
\end{enumerate}

\begin{savaistu}[Le Cameroun, « l'Afrique en miniature »]
On surnomme souvent le Cameroun \emph{l'Afrique en miniature} : on y retrouve, concentrés sur un seul territoire, presque tous les paysages, climats et systèmes agricoles du continent — forêt équatoriale, savanes, hautes terres volcaniques, zones sèches sahéliennes. C'est pourquoi l'étude de l'agriculture camerounaise est une porte d'entrée vers la compréhension de toute l'agriculture africaine.
\end{savaistu}

\courssub{Observer et décrire une photographie agricole : la situation-problème d'ouverture}

Revenons à nos deux photographies d'ouverture : le champ de brûlis en zone forestière du Sud, et le champ de maïs intensifié à Foumbot. Au Bac, tu seras souvent confronté à ce type de documents. Voici la démarche rigoureuse pour les exploiter.

\begin{methode}[Décrire une photographie agricole en quatre étapes]
\begin{enumerate}
\item \textbf{Localiser} : identifier le lieu, la région, la zone agro-écologique (indices : végétation, relief, habitat).
\item \textbf{Décrire objectivement} : les outils visibles (houe, machette, tracteur ?), les cultures (une seule ou plusieurs espèces associées ?), l'état du champ (brûlis récent, rangs réguliers ?), les acteurs (famille, salariés ?) et leurs gestes.
\item \textbf{Interpréter} : dégager le système agricole (itinérant sur brûlis, vivrier intensifié, plantation industrielle ?) et le niveau d'amélioration technique (intrants, semences, mécanisation).
\item \textbf{Conclure} : relier le document à la notion du cours et à la problématique (ici : la coexistence de plusieurs degrés d'amélioration de l'agriculture traditionnelle).
\end{enumerate}
\end{methode}

\begin{exoresolu}[Commentaire comparé de deux photographies agricoles]
\textbf{Énoncé.} Photo A : en forêt du Sud, près d'Ebolowa, une parcelle fraîchement brûlée ; on distingue des troncs noircis, des cendres, et une femme qui plante des boutures de manioc à l'aide d'un bâton fouisseur. Photo B : à Foumbot (Ouest), un champ de maïs en rangs réguliers, un planteur épand de l'engrais, une bâche publicitaire d'une coopérative annonce « semences hybrides ». Décris chaque photo et dégage ce qu'elles révèlent de l'agriculture camerounaise.
\tcblower
\textbf{Solution détaillée.}
\emph{Photo A (Ebolowa).} Lieu : zone forestière humide du Sud. Description : parcelle défrichée puis brûlée (troncs calcinés, cendres qui fertilisent le sol), outillage rudimentaire (bâton fouisseur), culture vivrière (manioc), main-d'œuvre familiale. Interprétation : il s'agit de l'agriculture itinérante sur brûlis, forme classique de l'agriculture traditionnelle de subsistance, adaptée aux sols forestiers pauvres mais exigeant de vastes espaces pour la jachère.

\emph{Photo B (Foumbot).} Lieu : Hautes Terres de l'Ouest. Description : monoculture de maïs en lignes régulières, usage d'engrais et de semences hybrides achetées via une coopérative, production orientée vers la vente. Interprétation : il s'agit d'une agriculture familiale intensifiée — une agriculture de « seconde génération » — qui reste de petite taille et de type familial, mais qui a intégré des techniques modernes.

\emph{Conclusion.} Les deux photos montrent que l'agriculture camerounaise n'est pas uniforme : elle reste partout dominée par les petits exploitants familiaux, mais le degré d'amélioration technique varie fortement selon les régions. L'agriculture traditionnelle n'est donc pas figée : elle se transforme, à des rythmes différents, sous l'effet de l'outillage amélioré, des nouvelles techniques culturales et de l'adoption de nouvelles cultures — ce que la suite de la leçon va détailler.
\end{exoresolu}

\courssub{Complément utile au Bac : agriculture traditionnelle et agriculture moderne}

Pour briller dans une dissertation ou un commentaire de document, il faut savoir situer l'agriculture traditionnelle par rapport à l'autre grand type d'agriculture présent au Cameroun : l'\cle{agriculture moderne} (ou d'entreprise). Ce point est un \textbf{complément} : il ne constitue pas le cœur de la leçon, mais il est précieux pour les comparaisons.

\begin{center}
\begin{tabularx}{\linewidth}{|l|X|X|}
\hline
\rowcolor{popblueL} \textbf{Critère} & \textbf{Agriculture traditionnelle} & \textbf{Agriculture moderne} \\
\hline
Taille & Moins de 2 ha en général & Centaines ou milliers d'hectares \\
\hline
Main-d'œuvre & Familiale & Salariée \\
\hline
Outils & Machette, houe, daba & Tracteurs, engins, irrigation \\
\hline
Intrants & Peu ou pas & Engrais, pesticides, semences sélectionnées \\
\hline
Production & Vivrière, autoconsommation + surplus & Cultures de rente destinées à la vente et à l'exportation \\
\hline
Exemples au Cameroun & Champs familiaux de manioc, maïs, mil & CDC (bananes, palmiers, hévéas), PHP à Penja, SOSUCAM (sucre), SODECOTON (coton, dans le Grand Nord) \\
\hline
\end{tabularx}
\end{center}

\begin{attention}[Piège de rédaction]
N'oppose pas les deux systèmes comme si l'un était « bon » et l'autre « mauvais ». Les plantations modernes (CDC, PHP, SOSUCAM) emploient des salariés et exportent, mais elles cohabitent avec des millions de petits planteurs qui nourrissent le pays. Et surtout : la frontière bouge. Un planteur de Foumbot qui achète des engrais et vend son maïs à une coopérative n'est plus un pur « traditionnel » ni encore une « plantation moderne » : c'est précisément toute la richesse de la notion d'\emph{agriculture de seconde génération}, que nous étudierons dans la suite de la leçon. Autre piège : situe correctement les entreprises — la SODECOTON opère dans le Grand Nord (Adamaoua, Nord, Extrême-Nord), pas dans l'Ouest.
\end{attention}

\begin{aretenir}[L'essentiel de la section]
L'agriculture traditionnelle camerounaise est une agriculture de subsistance, à faible mécanisation, dominée par le travail familial sur de petites exploitations de moins de 2 hectares. Elle occupe environ 6 actifs sur 10, nourrit le pays et varie fortement selon les zones agro-écologiques : brûlis en forêt du Sud, cultures intensives sur les Hautes Terres de l'Ouest, cultures sèches et élevage au Nord. Loin d'être archaïque, elle connaît de grandes améliorations — outillage, techniques, nouvelles cultures — que les sections suivantes vont analyser en détail.
\end{aretenir}

\coursec{Un outillage amélioré}

Dans la section précédente, nous avons dressé le portrait de l'agriculture traditionnelle camerounaise : une agriculture familiale, de subsistance, pratiquée sur de petites parcelles par des paysans dont le principal capital est la force de leurs bras. Mais ce portrait serait incomplet si l'on s'arrêtait à l'image figée du paysan et de sa houe. Depuis plusieurs décennies, l'outillage agricole camerounais s'est profondément transformé : la machette côtoie désormais la charrue, le motoculteur et même, dans certaines plaines, le tracteur. Cette section étudie cette révolution discrète mais réelle : celle de l'outil.

\courssub{L'outillage traditionnel de base : l'héritage des générations}

\begin{definition}[L'outillage agricole]
L'\cle{outillage agricole} désigne l'ensemble des instruments, outils et machines utilisés par l'agriculteur pour préparer le sol, semer, entretenir les cultures, récolter et transformer les produits. Le niveau technique de cet outillage est un excellent indicateur du niveau de développement d'une agriculture.
\end{definition}

Pendant des siècles, le paysan camerounais n'a disposé que d'un petit nombre d'outils manuels, fabriqués localement par les forgerons villageois :

\begin{itemize}
\item la \cle{houe} (ou hilaire), outil polyvalent par excellence : elle sert à retourner la terre, à former les billons, à sarcler et à butter les plants ;
\item la \cle{machette}, indispensable pour défricher, abattre les herbes hautes, couper les régimes de bananes ou ouvrir les cosses de cacao ;
\item la \cle{daba}, petite houe à manche court, utilisée surtout par les femmes pour les travaux fins (sarclage, binage) ;
\item la \cle{hache}, pour abattre les arbres lors du défrichement, notamment dans le système de l'itinérance sur brûlis ;
\item le \cle{couteau de récolte}, petit outil à lame courbe pour couper les épis de mil, de sorgho ou les panicules.
\end{itemize}

Cet outillage présente deux qualités majeures : il est \textbf{peu coûteux} et \textbf{parfaitement adapté} aux parcelles morcelées, aux pentes et aux cultures associées. Mais il a une faiblesse décisive : il repose entièrement sur l'énergie musculaire humaine. Or, une journée de travail à la houe épuise l'homme et ne permet de préparer qu'une surface très limitée (de l'ordre de quelques ares par jour selon la nature du sol). Tant que l'outil reste manuel, la surface cultivée par actif agricole reste faible, et la production plafonne.

\begin{savaistu}[Le forgeron, personnage central du village]
Dans la plupart des sociétés camerounaises, le \cle{forgeron} occupe une place sociale particulière, parfois entourée de prestige ou de crainte. C'est lui qui fond, martèle et affûte les lames des houes et des machettes, et c'est encore lui, aujourd'hui, qui répare les charrues et soude les châssis des motoculteurs. La modernisation de l'outillage n'a donc pas fait disparaître l'artisanat local : elle l'a transformé.
\end{savaistu}

\courssub{L'amélioration des outils manuels : la première étape}

La modernisation n'a pas commencé par le tracteur, mais par de simples perfectionnements des outils existants. C'est la \textbf{première génération d'améliorations} :

\begin{itemize}
\item les \textbf{houes à manche renforcé} et à lame en acier importé ou recyclé (ressorts de camion, rails), plus résistantes et mieux affûtées ;
\item les \textbf{pulvérisateurs à dos}, qui permettent de traiter les cultures (coton, cacao, maraîchage) avec des pesticides et des herbicides — l'outil devient alors le complément indispensable de l'intrant chimique ;
\item la \textbf{brouette}, qui soulage le transport des récoltes, du fumier et des plants sur la parcelle ;
\item les \textbf{charrues à bras} (traînées par l'homme), puis les \textbf{charrues attelées}, qui ouvrent la voie à la traction animale.
\end{itemize}

Ces améliorations semblent modestes, mais leurs effets sont considérables : un pulvérisateur permet de traiter en une journée une surface qui aurait demandé plusieurs jours de désherbage manuel. L'outil amélioré libère du temps de travail, que le paysan peut réinvestir dans l'extension de ses parcelles.

\courssub{La traction animale : la révolution des savanes du Nord}

\begin{definition}[Traction animale]
La \cle{traction animale} est l'utilisation d'animaux de trait — principalement des \textbf{bœufs} et des \textbf{ânes}, parfois des chevaux ou des mulets — pour tirer la charrue et les autres instruments aratoires (herse, buttoir, charrette). Elle substitue l'énergie animale à l'énergie musculaire humaine pour les travaux les plus pénibles.
\end{definition}

Pourquoi la traction animale s'est-elle développée dans les savanes de l'\textbf{Adamaoua}, du \textbf{Nord} et de l'\textbf{Extrême-Nord}, et non dans la forêt du Sud ? Le raisonnement géographique s'impose :

\begin{enumerate}
\item \textbf{Le relief} : les savanes septentrionales offrent de vastes plaines et plateaux peu accidentés, praticables par un attelage, contrairement aux collines et aux forêts denses du Centre et du Sud.
\item \textbf{Le cheptel} : ces régions sont des zones d'élevage traditionnel ; les bœufs et les ânes y sont nombreux et disponibles.
\item \textbf{La végétation} : la savane se défriche plus facilement que la forêt, et le sol s'y prête bien au passage de la charrue.
\item \textbf{Les cultures} : le coton, le maïs, le sorgho et le niébé, cultivés en lignes, sont compatibles avec le labour attelé.
\item \textbf{Un frein sanitaire au Sud} : la mouche tsé-tsé, vecteur de la \textbf{trypanosomiase animale (nagana)} et de la maladie du sommeil chez l'homme, décime les bovins en zone forestière humide et y interdit durablement l'élevage de trait.
\end{enumerate}

Introduite dès l'époque coloniale puis relancée dans les années 1970-1980 (notamment dans le sillage de la \textbf{SODECOTON} dans le bassin cotonier), la culture attelée s'est diffusée dans tout le bassin cotonier du Nord et de l'Adamaoua. L'attelage est composé d'une paire de bœufs dressés, d'un joug et d'une charrue à soc simple. Dans l'Extrême-Nord, l'âne, sobre et résistant, est aussi utilisé pour le travail des champs et le transport des récoltes sur charrette.

\begin{exemplebox}[Un gain de productivité mesurable]
Une charrue à traction animale permet de cultiver \textbf{3 à 5 fois plus de surface qu'une houe} en une journée de travail. Concrètement : là où un homme à la houe prépare environ 0,1 à 0,2 hectare par jour, un attelage bovin laboure 0,5 à 1 hectare. C'est ce gain qui a permis aux paysans du bassin cotonier d'étendre leurs surfaces de coton et de maïs, et de dégager un surplus commercial.
\end{exemplebox}

\courssub{La motorisation partielle : motoculteurs, tracteurs et moulins}

La troisième étape de l'amélioration est la \cle{motorisation}, c'est-à-dire l'emploi de moteurs pour remplacer la force humaine et animale. Au Cameroun, elle reste \textbf{partielle et inégalement répartie} :

\begin{itemize}
\item le \textbf{motoculteur}, tracteur à deux roues que l'on guide à pied, adapté aux petites et moyennes parcelles ; on le rencontre surtout dans les périmètres maraîchers et rizicoles ;
\item le \textbf{tracteur}, rare en propriété individuelle : il est le plus souvent utilisé en \textbf{location} (auprès de coopératives, de sociétés de développement ou de gros exploitants), pour le labour des grandes plaines (bassin de la Bénoué, plaine de la Mbo à l'Ouest/Littoral, périmètres de la SEMRY) ;
\item les \textbf{décortiqueuses} (riz) et les \textbf{moulins à grains} (maïs, mil, sorgho), installés au cœur des villages : ils ne travaillent pas le sol, mais ils libèrent les femmes de l'interminable pilage au mortier — c'est une mécanisation de la \emph{transformation}, non de la production.
\end{itemize}

\begin{attention}[Erreur fréquente au Bac]
N'\textbf{écris} jamais que l'agriculture camerounaise est « mécanisée ». La mécanisation reste \textbf{marginale} dans les exploitations familiales : l'immense majorité des paysans camerounais travaille encore avec la houe et la machette. Le tracteur concerne surtout les périmètres aménagés, les sociétés de développement et une minorité d'exploitants aisés. Formulation correcte : « l'outillage s'est amélioré, mais la mécanisation demeure embryonnaire et concentrée ».
\end{attention}

\courssub{Le rôle des artisans locaux et des forgerons}

Un outillage, même amélioré, ne sert à rien s'il ne peut être fabriqué, entretenu et réparé localement. C'est ici qu'interviennent les \textbf{artisans ruraux} :

\begin{itemize}
\item les \textbf{forgerons} fabriquent les lames de houes, de daba et de machettes, souvent à partir d'acier de récupération (ressorts de camions, rails), et réparent les socs de charrue ;
\item les \textbf{charpentiers et menuisiers} construisent les jougs, les mancherons et les charrettes ;
\item les \textbf{petits mécaniciens} des villes et des marchés ruraux (les « quartiers mécaniques » de Garoua, Ngaoundéré ou Maroua) assurent l'entretien des motoculteurs et des moulins.
\end{itemize}

Ce tissu artisanal est un maillon essentiel de l'intégration nationale : il diffuse l'innovation technique jusque dans les villages, crée des emplois non agricoles en milieu rural et évite une dépendance totale aux importations de matériel.

\courssub{Les limites de la mécanisation au Cameroun}

Pourquoi, malgré ses avantages évidents, la mécanisation progresse-t-elle si lentement ? Plusieurs obstacles se combinent :

\begin{enumerate}
\item \textbf{Le coût d'achat et d'entretien} : un tracteur neuf coûte plusieurs dizaines de millions de FCFA ; même un motoculteur dépasse les moyens d'un exploitant familial. S'y ajoutent le carburant, les pièces détachées importées et la main-d'œuvre spécialisée.
\item \textbf{Le morcellement des parcelles} : l'exploitation familiale camerounaise mesure souvent moins de 2 hectares, fragmentée en plusieurs lopins éloignés ; un tracteur y est économiquement absurde.
\item \textbf{Le relief accidenté} : dans les hautes terres de l'Ouest ou les collines du Centre, les pentes fortes interdisent le passage des machines.
\item \textbf{Le manque de crédit agricole} : sans accès au prêt bancaire ni à des coopératives solides, le paysan ne peut investir ; les structures de financement rural restent faibles.
\item \textbf{La disponibilité de la main-d'œuvre familiale} : tant que la famille fournit gratuitement sa force de travail, l'investissement mécanique n'est pas jugé prioritaire.
\end{enumerate}

\begin{methode}[Classer un outillage selon son niveau de technicité]
Pour analyser un document (photographie, texte, tableau) portant sur l'outillage agricole camerounais, procède en trois étapes :
\begin{enumerate}
\item \textbf{Identifier} les outils visibles ou cités (houe, charrue, motoculteur, tracteur...).
\item \textbf{Classer} chaque outil selon la source d'énergie : \emph{manuel} (énergie humaine), \emph{animal} (traction animale), \emph{motorisé} (énergie thermique).
\item \textbf{Localiser} : relie chaque niveau à sa zone de prédilection — outillage manuel dominant dans le Sud forestier et les hautes terres de l'Ouest ; traction animale dans les savanes de l'Adamaoua, du Nord et de l'Extrême-Nord ; motorisation dans les périmètres aménagés (SEMRY, plaine de la Mbo) et les grandes plaines.
\end{enumerate}
Conclus toujours par la nuance : les trois niveaux \emph{coexistent} aujourd'hui au Cameroun.
\end{methode}

\begin{popfigure}
\begin{center}
\begin{tikzpicture}[scale=0.9, font=\small]
% --- Panneau traditionnel ---
\fill[popcream] (-5.2,-2.6) rectangle (-0.3,2.6);
\node[font=\small\bfseries, popdark] at (-2.75,2.25) {Outillage traditionnel};
\node[font=\footnotesize\itshape, popmuted] at (-2.75,1.85) {énergie humaine};
% Houe
\draw[line width=2pt, popdark] (-4.6,-0.4) -- (-3.7,0.9);
\draw[line width=2.5pt, popink] (-3.7,0.9) -- (-3.1,0.55);
\node[font=\footnotesize] at (-4.1,-0.75) {houe};
% Machette
\draw[line width=2pt, popdark] (-1.9,-0.4) -- (-1.9,0.5);
\draw[line width=2.5pt, popink] (-1.9,0.5) .. controls (-1.75,1.1) and (-1.55,1.25) .. (-1.35,1.3);
\node[font=\footnotesize] at (-1.7,-0.75) {machette};
% Paysan stylisé
\draw[poporange, line width=1.2pt] (-3.0,-1.9) circle (0.22);
\draw[poporange, line width=1.2pt] (-3.0,-1.68) -- (-3.0,-1.1);
\draw[poporange, line width=1.2pt] (-3.0,-1.1) -- (-3.25,-0.75);
\draw[poporange, line width=1.2pt] (-3.0,-1.1) -- (-2.75,-0.75);
\draw[poporange, line width=1.2pt] (-3.0,-1.5) -- (-2.6,-1.35);
\node[font=\footnotesize, popmuted] at (-2.75,-2.3) {force des bras};
% --- Panneau amélioré ---
\fill[popblueL] (0.3,-2.6) rectangle (5.2,2.6);
\node[font=\small\bfseries, popdark] at (2.75,2.25) {Outillage amélioré};
\node[font=\footnotesize\itshape, popmuted] at (2.75,1.85) {énergie animale et motorisée};
% Charrue + boeuf
\draw[popgreen, line width=1.4pt] (0.8,0.1) ellipse (0.55 and 0.32);
\draw[popgreen, line width=1.4pt] (1.3,0.35) circle (0.18);
\draw[popgreen, line width=1.2pt] (0.55,-0.2) -- (0.55,-0.6);
\draw[popgreen, line width=1.2pt] (1.0,-0.2) -- (1.0,-0.6);
\draw[popgreen, line width=1.2pt] (1.42,0.5) -- (1.55,0.68);
\node[font=\footnotesize] at (1.0,-0.95) {bœuf de trait};
\draw[line width=1.6pt, popink] (1.6,-0.35) -- (2.5,-0.35);
\draw[line width=1.8pt, popink] (2.5,-0.35) -- (2.85,-0.75);
\draw[line width=1.8pt, popink] (2.5,-0.35) -- (2.85,0.15);
\node[font=\footnotesize] at (2.7,-1.05) {charrue};
% Motoculteur
\draw[line width=1.6pt, poppurple] (3.6,-0.2) -- (4.5,-0.2) -- (4.7,0.5);
\draw[poppurple, line width=1.4pt] (3.75,-0.55) circle (0.28);
\draw[poppurple, line width=1.4pt] (4.35,-0.55) circle (0.28);
\draw[line width=1.6pt, poppurple] (4.5,-0.2) -- (4.95,-0.6);
\node[font=\footnotesize] at (4.25,-1.05) {motoculteur};
\node[font=\footnotesize, popmuted] at (2.75,-2.3) {3 à 5 fois plus de surface / jour};
\end{tikzpicture}
\end{center}
\legende{Schéma comparatif de l'outillage agricole au Cameroun : à gauche, les outils manuels traditionnels (houe, machette) reposant sur la seule énergie humaine ; à droite, l'outillage amélioré (charrue à traction animale, motoculteur) qui multiplie la surface travaillée par journée de travail.}
\end{popfigure}

\courssub{Étude de cas : la motorisation des rizières de la SEMRY à Yagoua}

\begin{exoresolu}[Étude de cas : la SEMRY, vitrine de la mécanisation rizicole]
\textbf{Énoncé.} La SEMRY (Société d'expansion et de modernisation de la riziculture de Yagoua), créée en 1971 dans l'Extrême-Nord, aménage et gère plusieurs dizaines de milliers d'hectares de rizières irriguées dans la plaine inondable du Logone (barrage de Maga). À partir de ces informations et de tes connaissances : (1) montre en quoi la SEMRY illustre l'amélioration de l'outillage agricole ; (2) explique pourquoi ce modèle reste difficile à généraliser à l'ensemble du pays.
\tcblower
\textbf{Solution détaillée.}
(1) La SEMRY illustre le niveau le plus avancé de l'amélioration de l'outillage. Les travaux lourds de préparation du sol (nivellement, labour) y sont réalisés par des \textbf{tracteurs}, possibles grâce à la plaine parfaitement plate et aux parcelles regroupées et régulières. Les paysans utilisent aussi des motoculteurs, des pulvérisateurs et des charrettes ; en aval, des \textbf{unités de décorticage} transforment le riz paddy en riz blanc destiné au marché national. L'eau, maîtrisée par les digues et canaux issus du barrage de Maga, complète ce dispositif technique. La SEMRY montre donc la combinaison gagnante : \emph{motorisation + maîtrise de l'eau + encadrement + crédit en intrants}.

(2) Ce modèle peine à se généraliser pour plusieurs raisons : il exige des investissements publics massifs (barrages, canaux, pistes) que l'État ne peut reproduire partout ; il suppose un relief de plaine, absent des hautes terres et de la zone forestière ; il repose sur un encadrement et un crédit que l'exploitant familial isolé ne trouve pas. La SEMRY reste donc une \textbf{exception réussie} qui confirme la règle : la mécanisation camerounaise est localisée, liée aux périmètres aménagés, et non généralisée aux exploitations familiales.
\end{exoresolu}

\begin{popfigure}
\begin{center}
\begin{tikzpicture}
\begin{axis}[
  width=0.85\linewidth, height=6.2cm,
  ybar, bar width=7pt,
  ymin=0, ymax=5200,
  xlabel={Décennie (ordre de grandeur indicatif)},
  ylabel={Nombre d'unités en service},
  symbolic x coords={1970,1980,1990,2000,2010,2020},
  xtick=data,
  legend style={at={(0.02,0.97)}, anchor=north west, font=\footnotesize},
  ymajorgrids=true, grid style={popline!50},
  every axis plot/.append style={fill opacity=0.85}
]
\addplot+[ybar, fill=popblue, draw=popblue] coordinates {(1970,400) (1980,900) (1990,1400) (2000,1800) (2010,2400) (2020,3000)};
\addplot+[ybar, fill=poporange, draw=poporange] coordinates {(1970,100) (1980,500) (1990,1200) (2000,2200) (2010,3500) (2020,4800)};
\legend{Tracteurs, Motoculteurs}
\end{axis}
\end{tikzpicture}
\end{center}
\legende{Évolution indicative du parc de tracteurs et de motoculteurs au Cameroun depuis les années 1970 (ordres de grandeur reconstitués à partir de données FAO et ministérielles ; les chiffres précis varient selon les sources). Lecture : le parc progresse régulièrement, les motoculteurs plus vite que les tracteurs, mais il demeure très faible au regard des millions d'exploitations familiales — preuve que la mécanisation reste marginale.}
\end{popfigure}

\begin{exercice}[Pour t'entraîner]
Classe les outils suivants selon leur niveau de technicité (manuel, animal, motorisé) et localise la zone du Cameroun où chacun domine : daba, attelage bovin, décortiqueuse de riz, machette, motoculteur, charrue à âne. Rédige ensuite une phrase de synthèse sur l'état de la mécanisation agricole camerounaise.
\end{exercice}

\begin{corrige}[Exercice]
\emph{Manuel} : daba et machette — dominent dans la zone forestière du Sud, du Centre et du Littoral, ainsi que sur les hautes terres de l'Ouest. \emph{Animal} : attelage bovin (Adamaoua, Nord) et charrue à âne (Extrême-Nord). \emph{Motorisé} : décortiqueuse de riz (périmètres de la SEMRY, villages rizicoles) et motoculteur (plaines aménagées, maraîchage périurbain). \emph{Synthèse} : l'outillage agricole camerounais s'est nettement amélioré, mais les trois niveaux de technicité coexistent, et la mécanisation reste marginale, concentrée dans les périmètres aménagés et les savanes septentrionales.
\end{corrige}

\begin{aretenir}[L'essentiel de la section]
\begin{itemize}
\item L'outillage traditionnel (houe, machette, daba, hache, couteau de récolte) repose sur la seule énergie humaine et limite les surfaces cultivées.
\item L'amélioration a suivi trois étapes : perfectionnement des outils manuels (pulvérisateur, brouette), traction animale (savanes du Nord : bœufs et ânes), motorisation partielle (motoculteurs, tracteurs en location, moulins villageois).
\item Les artisans et forgerons locaux fabriquent et réparent l'outillage : ils diffusent l'innovation en milieu rural.
\item La mécanisation bute sur le coût, le morcellement foncier, le relief, l'absence de crédit : elle reste marginale, sauf dans les périmètres aménagés comme la SEMRY à Yagoua.
\end{itemize}
\end{aretenir}

\coursec{L'évolution des techniques culturales}

\noindent Après avoir vu comment l'outillage s'améliore progressivement — de la houe à la charrue attelée, voire au petit tracteur dans les plaines du Nord — il faut comprendre qu'un outil, aussi performant soit-il, ne produit rien sans la \cle{technique culturale} qui l'accompagne. C'est l'ensemble des savoir-faire, des gestes et des décisions agronomiques (choix de la semence, date de semis, fertilisation, protection des cultures, travail du sol) qui transforme un champ en récolte. Au Cameroun, cette évolution est spectaculaire : on passe d'une agriculture extensive, fondée sur la régénération naturelle de la fertilité par la jachère, à une agriculture de plus en plus intensive, pilotée par la recherche et l'encadrement. Observons cette mutation à travers des situations concrètes vécues par nos paysans.

\courssub{De la jachère longue à la jachère raccourcie : la pression démographique}

\begin{definition}[Jachère]
Période de repos d'une terre cultivée pendant laquelle la végétation naturelle (herbes, arbustes, arbres) recolonise la parcelle, permettant la régénération naturelle de la fertilité du sol : reconstitution de la matière organique, remontée des éléments minéraux par les racines profondes, limitation biologique des parasites et des adventices.
\end{definition}

\noindent Traditionnellement, dans les zones de forêt humide (Sud, Est, Centre) comme dans les savanes guinéennes (Adamaoua, Ouest), le paysan pratique l'\cle{agriculture itinérante sur brûlis}. Il défriche une parcelle de forêt ou de jachère longue (10 à 20 ans), brûle la biomasse (les cendres apportent potasse et phosphore), cultive 2 à 3 ans, puis abandonne le champ à la jachère pour aller en ouvrir un autre plus loin. Ce système est durable \emph{tant que la densité de population reste faible} et que l'espace « libre » existe.

\begin{popfigure}
\begin{tikzpicture}[scale=0.9, every node/.style={font=\small}]
% Cycle TRADITIONNEL (Long)
\node[align=center, font=\bfseries] at (2,5.5) {Cycle traditionnel (faible pression)};
\draw[fill=popgreen!30, draw=popgreen, thick, rounded corners] (0.5,4.5) rectangle (3.5,5.2);
\node[align=center] at (2,4.85) {Forêt /\\ Jachère longue (15--20 ans)};
\draw[->, thick, poporange] (3.5,4.85) -- (4.5,4.85) node[midway, above, font=\scriptsize, align=center] {Défrichement\\ \& brûlis};
\draw[fill=popgold!50, draw=poporange, thick, rounded corners] (4.5,4.5) rectangle (7.5,5.2);
\node[align=center] at (6,4.85) {Culture\\ (2--3 ans)};
\draw[->, thick, poporange] (7.5,4.85) -- (8.5,4.85) node[midway, above, font=\scriptsize] {Abandon};
\draw[fill=popgreenL, draw=popgreen, thick, rounded corners] (8.5,4.5) rectangle (11.5,5.2);
\node[align=center] at (10,4.85) {Jachère longue\\ Régénération complète};
\draw[->, thick, poporange, dashed] (10,4.5) .. controls (10,3.5) and (2,3.5) .. (2,4.5) node[midway, below, align=center, font=\scriptsize] {Retour au même emplacement après 15--20 ans};

% Cycle ACTUEL (Raccourci)
\node[align=center, font=\bfseries] at (2,2.5) {Cycle actuel (forte pression démographique)};
\draw[fill=popgreenL, draw=popgreen, thick, rounded corners] (0.5,1.5) rectangle (3.5,2.2);
\node[align=center] at (2,1.85) {Jachère courte (3--5 ans)\\ Régénération incomplète};
\draw[->, thick, poporange] (3.5,1.85) -- (4.5,1.85) node[midway, above, font=\scriptsize, align=center] {Défrichement\\ \& brûlis};
\draw[fill=popgold!50, draw=poporange, thick, rounded corners] (4.5,1.5) rectangle (7.5,2.2);
\node[align=center] at (6,1.85) {Culture\\ (2--3 ans)};
\draw[->, thick, poporange] (7.5,1.85) -- (8.5,1.85) node[midway, above, font=\scriptsize, align=center] {Réoccupation\\ rapide};
\draw[fill=popgreen!15, draw=popgreen!50!poporange, thick, rounded corners] (8.5,1.5) rectangle (11.5,2.2);
\node[align=center] at (10,1.85) {Jachère très courte /\\ Culture continue};
\draw[->, thick, poporange, dashed] (10,1.5) .. controls (10,0.5) and (2,0.5) .. (2,1.5) node[midway, below, align=center, font=\scriptsize] {Retour rapide (cycle total de 5--8 ans)};

% Indicateurs de dégradation
\draw[decorate, decoration={brace, amplitude=5pt}, thick, popdark] (0.5,1.2) -- (3.5,1.2) node[midway, below=3pt, align=center, font=\scriptsize] {Perte de fertilité, érosion,\\ envahissement par \emph{Imperata}};
\draw[decorate, decoration={brace, amplitude=5pt}, thick, popdark] (8.5,1.2) -- (11.5,1.2) node[midway, below=3pt, align=center, font=\scriptsize] {Sol épuisé,\\ rendements en chute};
\end{tikzpicture}
\legende{Schéma comparatif du cycle cultural : raccourcissement drastique de la phase de jachère sous l'effet de la pression démographique (ex. : Hauts-Plateaux de l'Ouest et du Nord-Ouest). La jachère longue (15--20 ans) assure la durabilité ; la jachère courte (3--5 ans) entraîne une spirale de dégradation.}
\end{popfigure}

\noindent Aujourd'hui, dans les Hauts-Plateaux de l'Ouest (pays bamiléké, bamoun), au Nord-Ouest ou dans la vallée du Mbam, la densité dépasse souvent \num{100} à \num{200} hab/km\textsuperscript{2}, avec des pointes locales bien supérieures. La « terre neuve » a disparu. Le paysan est contraint de raccourcir la jachère à \num{3}--\num{5} ans, voire de cultiver en permanence (\cle{culture continue}). Conséquence directe : la matière organique chute, la structure du sol se dégrade, l'érosion emporte l'horizon superficiel, et les mauvaises herbes pérennes comme l'\emph{Imperata cylindrica} (herbe à épée) s'installent, rendant le travail du sol extrêmement pénible. C'est ce « piège de la jachère raccourcie » qui force l'adoption de techniques de substitution de la fertilité naturelle.

\courssub{L'intensification par les intrants : engrais et traitements phytosanitaires}

\noindent Face à l'épuisement des sols, le paysan n'a d'autre choix que d'apporter \emph{de l'extérieur} les éléments nutritifs que la jachère ne fournit plus. On distingue deux voies complémentaires :

\begin{itemize}
    \item \textbf{La fumure organique :} fumier de bovins, de porcs, de volailles ; compost de déchets ménagers et de résidus de culture. Dans l'Ouest, l'élevage porcin et avicole en stabulation se développe en partie pour produire du fumier. Au Nord, le fumier déposé par les troupeaux en transhumance est récupéré dans les parcs de nuit installés sur les champs (contrats entre agriculteurs et éleveurs).
    \item \textbf{Les engrais minéraux (chimiques) :} NPK (Azote--Phosphore--Potasse), urée (46 \% d'azote), engrais phosphatés. Leur usage s'est généralisé dans les filières encadrées (coton, riz, maïs). L'État, à travers des programmes de subvention des intrants et les sociétés de développement, rend le sac de \SI{50}{kg} plus accessible (ordre de grandeur : vendu environ \num{13 000}--\num{15 000} FCFA subventionné au lieu de \num{20 000}--\num{25 000} FCFA au prix du marché ; les montants varient selon les campagnes).
\end{itemize}

\noindent Parallèlement, la pression parasitaire augmente avec la concentration des cultures. Les \cle{traitements phytosanitaires} deviennent indispensables :
\begin{itemize}
    \item \textbf{Insecticides :} lutte contre les foreurs du maïs (\emph{Busseola fusca}), les mirides (capsides) du cacaoyer (\emph{Sahlbergella singularis}), les jassides du cotonnier.
    \item \textbf{Herbicides :} glyphosate (avant semis), atrazine (maïs), pour réduire la pénibilité du sarclage manuel et lutter contre \emph{Imperata} et \emph{Striga} (plante parasite du sorgho et du maïs, surtout au Nord).
    \item \textbf{Fongicides :} traitement des semences (coton, riz) et traitements foliaires (cacaoyer : pourriture brune des cabosses due à \emph{Phytophthora} ; bananier : cercosporiose noire ou \emph{Black Sigatoka}).
\end{itemize}

\begin{exemplebox}[Exemple chiffré : le saut de rendement du maïs]
Dans les Hauts-Plateaux de l'Ouest (par exemple autour de Dschang), un paysan non encadré, utilisant une semence locale, sans engrais, sur sol fatigué, récolte environ \textbf{1,0 à 1,5 t/ha}. Le même paysan, encadré par une coopérative, utilisant une variété améliorée (type \emph{CMS 8704}), environ \SI{200}{kg/ha} de NPK 20-10-10 à la plantation puis \SI{100}{kg/ha} d'urée en couverture, et un désherbage précoce, atteint couramment \textbf{3,5 à 4,5 t/ha}. Le rendement est \textbf{multiplié par 3 à 4}. L'outillage (houe ou charrue) n'explique qu'une faible part de ce gain ; c'est le « paquet technique » (semence + engrais + protection + conduite culturale) qui fait la différence.
\end{exemplebox}

\begin{attention}[Erreur fréquente au Bac]
Ne pas écrire : « Les rendements augmentent grâce aux nouveaux outils (charrue, tracteur). »
\emph{Correction :} l'outillage permet de travailler plus vite et sur de plus grandes surfaces (extensification), mais l'augmentation du rendement \textbf{à l'hectare} (intensification) vient des \textbf{techniques culturales} : semences améliorées, fertilisation organo-minérale, protection phytosanitaire, densité de semis optimisée, dates de semis respectées. Un champ labouré au tracteur mais semé tard, avec de mauvaises graines et sans engrais, rendra moins qu'un champ bien conduit à la houe.
\end{attention}

\courssub{La révolution des semences : recherche (IRAD) et vulgarisation}

\noindent Le cœur de l'intensification, c'est la \cle{semence}. Une variété améliorée porte en elle un potentiel génétique supérieur : précocité (elle échappe à la sécheresse de fin de cycle), résistance aux maladies et aux parasites, tolérance à l'acidité des sols, meilleur rapport grain/paille.

\begin{definition}[Vulgarisation agricole]
Ensemble des activités de diffusion des innovations techniques (nouvelles variétés, itinéraires techniques, utilisation des intrants) auprès des producteurs, réalisées par des agents de vulgarisation (conseillers agricoles, animateurs ruraux) afin d'améliorer la productivité et les revenus des exploitations familiales.
\end{definition}

\noindent Au Cameroun, l'\cle{IRAD} (Institut de Recherche Agricole pour le Développement), dont le siège est à Nkolbisson (Yaoundé), avec des antennes et stations régionales (Bambui au Nord-Ouest, Maroua à l'Extrême-Nord, Ekona au Sud-Ouest, Foumban à l'Ouest, Bertoua à l'Est, Garoua au Nord), est le moteur de la création variétale. Il travaille en partenariat avec des centres internationaux (IITA au Nigeria, CIMMYT au Mexique, AfricaRice) et des instituts étrangers (CIRAD, France).

\begin{savaistu}[L'IRAD, creuset des variétés camerounaises]
L'IRAD a mis au point ou contribué à diffuser des variétés adaptées à nos agro-écologies :
\begin{itemize}
    \item \textbf{Maïs :} variétés de la série \emph{CMS} (ex. \emph{CMS 8704}, \emph{CMS 9015}), précoces (environ 100--110 jours), certaines tolérantes au \emph{Striga} ; variétés QPM (\emph{Quality Protein Maize}) riches en lysine et tryptophane, intéressantes pour la nutrition.
    \item \textbf{Manioc :} variétés améliorées à haut rendement, tolérantes à la mosaïque africaine du manioc et à faible teneur en composés cyanogènes.
    \item \textbf{Riz :} variétés \emph{NERICA} (\emph{New Rice for Africa}), issues du croisement \emph{Oryza sativa} $\times$ \emph{Oryza glaberrima}, précoces (90--100 jours) et tolérantes à la sécheresse, diffusées notamment par la SEMRY et l'UNVDA.
    \item \textbf{Coton :} variétés sélectionnées avec l'appui de la recherche (résistance aux jassides, fibre de bonne qualité, haut rendement), diffusées par la SODECOTON.
\end{itemize}
La diffusion se fait via les \textbf{sociétés de développement} (SODECOTON, SEMRY, UNVDA, CDC, HEVECAM, SOCAPALM), les \textbf{ONG}, les \textbf{chambres d'agriculture} et les \textbf{coopératives}. Le paysan achète la semence certifiée (sachet scellé et étiqueté) auprès des distributeurs agréés.
\end{savaistu}

\courssub{Techniques de conservation des sols et maîtrise de l'eau}

\noindent L'intensification ne se résume pas aux intrants. Sur les pentes fortes de l'Ouest (souvent 20 à 40 \%) et face à l'érosion hydrique (ruissellement violent en saison des pluies), la conservation des eaux et des sols (\cle{CES}) est vitale.

\begin{itemize}
    \item \textbf{Le billonnage (culture sur billons) :} buttes de terre surélevées (environ \SI{30}{--}\SI{40}{cm} de hauteur) séparées par des sillons. Avantages : drainage des excès d'eau (évite l'asphyxie racinaire), réchauffement plus rapide du sol, concentration de la matière organique et des engrais sur la butte. Très utilisé pour l'igname, le taro (macabo), la patate douce et le manioc dans l'Ouest et le Sud-Ouest.
    \item \textbf{Les cultures en courbes de niveau :} labour, semis et billons perpendiculaires à la pente. Cela ralentit le ruissellement et favorise l'infiltration. Le tracé s'effectue avec des outils simples (niveau à eau, cadre en A).
    \item \textbf{L'agroforesterie et les haies vives :} plantation de lignes d'arbres ou d'arbustes légumineux (\emph{Leucaena leucocephala}, \emph{Calliandra calothyrsus}, \emph{Gliricidia sepium}) en bordure ou à l'intérieur des parcelles. Rôles : brise-vent, apport d'azote (les feuilles servent d'engrais vert), bois de feu, piquets, lutte anti-érosive. Promue notamment par le Centre mondial d'agroforesterie (ICRAF, antenne de Yaoundé) et divers projets de développement rural.
    \item \textbf{Le paillage :} couvrir le sol (résidus de culture, herbes fauchées, feuilles) pour limiter l'évaporation, l'érosion par battance et l'échauffement du sol, tout en nourrissant la vie microbienne.
\end{itemize}

\noindent La \cle{maîtrise de l'eau} permet la \cle{culture de contre-saison} (saison sèche) et sécurise la saison des pluies :
\begin{itemize}
    \item \textbf{Petits périmètres irrigués :} forages et puits équipés de motopompes (thermiques ou solaires), distribution par aspersion ou goutte-à-goutte pour le maraîchage (tomate, oignon, piment, chou) à l'Extrême-Nord (Logone-et-Chari, Mayo-Danay) et dans l'Adamaoua.
    \item \textbf{Barrages et retenues d'eau :} le barrage de \textbf{Lagdo} sur la Bénoué (Nord) régule les crues et alimente les périmètres rizicoles de la SEMRY ; le barrage de \textbf{Mbakaou} (Adamaoua) régule le débit de la Sanaga ; de nombreux petits barrages collinaires existent au Nord et dans l'Adamaoua.
    \item \textbf{Aménagement des bas-fonds (bas-fonds hydromorphes) :} nivellement, canaux d'irrigation et de drainage, diguettes. La \textbf{SEMRY} (Société d'Expansion et de Modernisation de la Riziculture de Yagoua) aménage de vastes périmètres rizicoles irrigués autour de Yagoua, Maga et dans la plaine du Logone (plusieurs dizaines de milliers d'hectares au total) ; l'\textbf{UNVDA} (\emph{Upper Nun Valley Development Authority}, basée à Ndop) développe la riziculture dans la haute vallée du Noun (Nord-Ouest), en lien avec la retenue de Bamendjin sur le Noun. Ces aménagements permettent jusqu'à deux cycles de riz par an (saison des pluies + contre-saison).
\end{itemize}

\courssub{L'encadrement paysan : levier de la diffusion technique}

\noindent La technique ne s'adopte pas par décret. Elle nécessite un \cle{encadrement} de proximité : conseil technique, crédit intrants (avance de semences et d'engrais remboursée à la récolte), collecte et commercialisation. Trois piliers structurent cet encadrement au Cameroun :

\begin{enumerate}
    \item \textbf{Les sociétés de développement :} structures parapubliques ou à mission de service public, généralement spécialisées par filière.
    \begin{itemize}
        \item \textbf{SODECOTON} (Société de Développement du Coton du Cameroun) : Grand Nord (Extrême-Nord, Nord, Adamaoua). Elle encadre de l'ordre de \num{250 000} à \num{300 000} producteurs (ordre de grandeur, variable selon les campagnes). Elle fournit les intrants à crédit, assure l'encadrement technique, collecte le coton-graine, l'égrène et commercialise la fibre. Modèle « coton-vivrier » : rotation coton/maïs/sorgho/arachide, le coton laissant au sol une fumure résiduelle profitable aux cultures vivrières suivantes.
        \item \textbf{SEMRY} (Yagoua) : riz irrigué (périmètres de Yagoua, Maga, plaine du Logone). Aménagement hydro-agricole, crédit intrants, mécanisation partielle (labour, battage), décorticage et commercialisation du riz.
        \item \textbf{UNVDA} (Ndop) : riz pluvial et irrigué dans la haute vallée du Noun (Nord-Ouest).
        \item \textbf{CDC} (\emph{Cameroon Development Corporation}) : banane, palmier à huile, hévéa (Sud-Ouest), avec des programmes d'encadrement des planteurs villageois (\emph{outgrowers}).
    \end{itemize}
    \item \textbf{Les coopératives et Groupements d'Initiative Commune (GIC) :} organisations paysannes de base (encadrées par la loi n\textsuperscript{o} 92/006 du 14 août 1992 relative aux sociétés coopératives et aux groupements d'initiative commune). Elles permettent l'achat groupé des intrants, la mutualisation du matériel (décortiqueuse, tricycle), l'accès au crédit (banques, établissements de microfinance) et la vente groupée (meilleur pouvoir de négociation).
    \item \textbf{Les chambres d'agriculture :} la \textbf{Chambre d'Agriculture, des Pêches, de l'Élevage et des Forêts du Cameroun (CAPEF)} et ses antennes régionales représentent les producteurs, participent à l'élaboration des politiques agricoles et appuient la professionnalisation des exploitations familiales.
\end{enumerate}

\begin{popfigure}
\begin{tikzpicture}
\begin{axis}[
    width=\linewidth,
    height=0.55\linewidth,
    xlabel={Année},
    ylabel={Rendement coton-graine (kg/ha)},
    title={Évolution schématique des rendements cotonniers selon l'encadrement (Grand Nord)},
    grid=major,
    grid style={dashed, gray!30},
    legend pos=north west,
    xmin=1980, xmax=2020,
    ymin=400, ymax=1800,
    xtick={1980,1990,2000,2010,2020},
    ytick={600,800,1000,1200,1400,1600},
    tick label style={font=\small},
    label style={font=\small},
    legend style={font=\scriptsize, fill opacity=0.8, draw opacity=1, draw=popdark},
]
\addplot[poporange, thick, mark=square*, mark size=2.5] coordinates {
    (1980, 750) (1990, 720) (2000, 680) (2010, 650) (2020, 600)
};
\addplot[popblue, thick, mark=*, mark size=2.5] coordinates {
    (1980, 800) (1990, 1050) (2000, 1250) (2010, 1400) (2020, 1500)
};
\legend{Sans encadrement (extensif, jachère courte), Avec encadrement SODECOTON (paquet technique complet)}
\end{axis}
\end{tikzpicture}
\legende{Courbe schématique d'évolution des rendements du coton-graine au Grand Nord du Cameroun (ordres de grandeur établis d'après les rapports SODECOTON/MINADER ; les valeurs réelles varient fortement selon les années climatiques). L'écart se creuse avec la généralisation des variétés améliorées et de la fertilisation. Dans les zones bien encadrées, le rendement moyen atteint ou dépasse environ \num{1200}--\num{1500} kg/ha.}
\end{popfigure}

\courssub{TP guidé : construire une courbe d'évolution des rendements}

\noindent Le géographe doit savoir transformer un tableau statistique en graphique lisible pour analyser une dynamique. Voici la démarche rigoureuse.

\begin{methode}[Construire une courbe d'évolution à partir d'un tableau statistique]
\begin{enumerate}
    \item \textbf{Analyser le tableau :} identifier la variable étudiée (rendement en kg/ha ou t/ha), l'unité, la période (années) et les séries (ex. : « paysans encadrés » / « paysans non encadrés »).
    \item \textbf{Choisir le type de graphique :} courbe pour une évolution temporelle continue ; histogramme pour des données discontinues. Ici : \textbf{courbe} à deux séries superposées.
    \item \textbf{Définir les échelles (axes) :}
    \begin{itemize}
        \item Abscisses (axe X) : le temps (années), à intervalles réguliers (ex. : 1 cm = 2 ans).
        \item Ordonnées (axe Y) : le rendement, à échelle régulière partant de \textbf{zéro} (ou d'une valeur minimale clairement indiquée) pour ne pas déformer la pente (ex. : 1 cm = 0,5 t/ha).
    \end{itemize}
    \item \textbf{Reporter les points :} pour chaque année, placer un point à l'intersection (année ; rendement), avec des symboles distincts par série (ronds, carrés).
    \item \textbf{Relier les points :} traits continus fins, sans prolonger au-delà des données.
    \item \textbf{Légender impérativement :} titre complet (quoi ? où ? quand ? source), noms des axes avec unités, légende des séries (couleurs/symboles).
\end{enumerate}
\end{methode}

\begin{exoresolu}[Construction de la courbe des rendements du maïs (Hauts-Plateaux de l'Ouest)]
\textbf{Énoncé :} à partir du tableau ci-dessous, construire la courbe d'évolution du rendement du maïs (t/ha) pour un paysan encadré (paquet technique : semence améliorée + engrais + herbicide) et un paysan non encadré, entre 2010 et 2020. Commenter la tendance.

\begin{center}
\begin{tabularx}{\linewidth}{|c|X|X|}\hline
\rowcolor{popblueL}\textbf{Année} & \textbf{Rendement paysan non encadré (t/ha)} & \textbf{Rendement paysan encadré (t/ha)} \\ \hline
2010 & 1,1 & 1,2 \\
2012 & 1,0 & 2,1 \\
2014 & 0,9 & 2,8 \\
2016 & 0,8 & 3,4 \\
2018 & 0,8 & 3,8 \\
2020 & 0,7 & 4,1 \\ \hline
\end{tabularx}
\end{center}

\tcblower
\textbf{Solution détaillée :}
\begin{enumerate}
    \item \textbf{Choix du graphique :} courbe à deux séries superposées (comparaison directe).
    \item \textbf{Axes :} X = années (2010 à 2020, pas de 2 ans) ; Y = rendement (0 à 4,5 t/ha, pas de 0,5 t/ha).
    \item \textbf{Report :} points (2010 ; 1,1), (2012 ; 1,0), \dots{} pour la série 1 (carrés orange) ; points (2010 ; 1,2), (2012 ; 2,1), \dots{} pour la série 2 (ronds bleus).
    \item \textbf{Traçé :} relier les points de la série 1 (trait orange) et de la série 2 (trait bleu épais).
    \item \textbf{Légende :} « Évolution comparée des rendements du maïs (t/ha) selon le mode d'encadrement, Hauts-Plateaux de l'Ouest, 2010--2020. Source : données d'enquête (hypothétiques, à visée pédagogique). »
\end{enumerate}
\textbf{Commentaire type :} la courbe « non encadré » montre une \textbf{baisse continue} (de 1,1 à 0,7 t/ha, soit une chute d'environ 36 \%), signe de l'épuisement des sols sur jachère raccourcie sans apports. La courbe « encadré » montre une \textbf{hausse forte et soutenue} (de 1,2 à 4,1 t/ha, soit un rendement multiplié par environ 3,4), preuve de l'efficacité du paquet technique. L'écart entre les deux séries passe de 0,1 t/ha en 2010 à 3,4 t/ha en 2020 : l'encadrement technique inverse la spirale de dégradation.
\end{exoresolu}

\begin{aretenir}[Bilan : les piliers de l'amélioration des techniques culturales]
L'évolution des techniques culturales au Cameroun repose sur la substitution de la fertilité naturelle (jachère longue) par une fertilité \textbf{gérée} :
\begin{enumerate}
    \item \textbf{Raccourcissement forcé de la jachère} sous l'effet de la pression démographique $\to$ nécessité d'intrants.
    \item \textbf{Paquet technique « semences améliorées (IRAD) + engrais organo-minéraux + protection phytosanitaire »} $\to$ gains de rendement massifs (facteur 3 à 4 sur le maïs, forts gains sur le coton et le riz).
    \item \textbf{Conservation des eaux et des sols (billonnage, courbes de niveau, agroforesterie, paillage)} $\to$ durabilité de l'intensification sur les pentes.
    \item \textbf{Maîtrise de l'eau (irrigation, bas-fonds aménagés)} $\to$ double culture, sécurisation des récoltes, maraîchage de contre-saison.
    \item \textbf{Encadrement institutionnel (SODECOTON, SEMRY, UNVDA, coopératives, CAPEF)} $\to$ vecteur indispensable de l'accès aux intrants, au conseil et au marché.
\end{enumerate}
L'outillage (section précédente) est le \textbf{bras} ; la technique culturale est le \textbf{cerveau} de l'agriculture modernisée.
\end{aretenir}

\coursec{L'adoption de nouvelles cultures}

Dans la section précédente, nous avons vu comment les paysans camerounais ont fait évoluer leurs \emph{techniques} : meilleure préparation du sol, associations culturales raisonnées, jachères raccourcies mais mieux gérées, petites irrigations. Mais améliorer la manière de cultiver ne suffit pas : il faut aussi savoir \emph{quoi} cultiver. Or, l'une des mutations les plus profondes de l'agriculture camerounaise depuis un siècle tient à un fait simple : une grande partie des plantes que le paysan camerounais cultive aujourd'hui \textbf{n'étaient pas là autrefois}. Le maïs de ton champ de Bafoussam, le manioc de ta grand-mère à Ebolowa, le cacao de la plantation familiale d'Obala, le coton des champs de Garoua : toutes ces plantes sont des « nouvelles venues », adoptées, testées, généralisées. C'est cette dynamique d'\cle{adoption de nouvelles cultures} que nous étudions maintenant.

\courssub{La notion d'adoption : comprendre le mécanisme}

\begin{definition}[Adoption d'une culture]
En géographie rurale, l'\cle{adoption} désigne le processus par lequel une plante \emph{exogène} (venue d'ailleurs) ou une variété \emph{améliorée} (sélectionnée par la recherche) est introduite dans un système de culture traditionnel, y est expérimentée, puis s'y diffuse jusqu'à devenir une culture ordinaire des paysans. L'adoption comprend trois étapes : l'\textbf{introduction} (la plante arrive), l'\textbf{expérimentation} (quelques pionniers l'essaient), la \textbf{diffusion} (elle se généralise de village en village).
\end{definition}

L'intuition est simple : un paysan ne change pas de culture par plaisir du risque. Il adopte une nouvelle plante seulement si elle lui apporte quelque chose de concret : plus de nourriture, plus d'argent, moins de travail, ou une meilleure adaptation à son terroir. L'adoption est donc un \emph{choix rationnel}, fait par des millions de petits producteurs, et non une simple obéissance à l'administration coloniale ou à l'État. C'est ce qui explique que certaines cultures imposées ont échoué, tandis que d'autres, librement choisies, ont connu un succès spectaculaire.

\begin{definition}[Culture de rente]
Une \cle{culture de rente} (ou culture commerciale, culture de rapport) est une culture destinée \textbf{principalement à la vente} — sur le marché national ou à l'exportation — et non à la consommation directe de la famille paysanne. Elle s'oppose à la \cle{culture vivrière}, destinée d'abord à nourrir le producteur. Au Cameroun : cacao, café, coton, hévéa, palmier à huile, banane d'exportation sont des cultures de rente ; manioc, igname, taro, mil sont des cultures vivrières (le maïs et la banane plantain jouent les deux rôles).
\end{definition}

\courssub{Les cultures vivrières adoptées ou généralisées}

\begin{savaistu}[Des plantes venues d'Amérique]
Le maïs, le manioc, le cacao et l'arachide sont tous originaires d'\textbf{Amérique}. Ils ont été introduits en Afrique par les navigateurs portugais après le XV\textsuperscript{e} siècle, dans le cadre des échanges colombiens entre le Nouveau et l'Ancien Monde. Le manioc, arrivé sur les côtes du golfe de Guinée au XVI\textsuperscript{e} siècle, n'a conquis l'intérieur du Cameroun qu'au XIX\textsuperscript{e} et au XX\textsuperscript{e} siècle. Autrement dit : l'« agriculture traditionnelle » que nous décrivons est déjà le produit de plusieurs siècles d'adoptions !
\end{savaistu}

Pourquoi ces plantes américaines ont-elles été adoptées si massivement ? Parce qu'elles répondaient mieux que les cultures anciennes (mil, sorgho, igname, fonio) à certains besoins.

\begin{itemize}
\item \textbf{Le maïs} : cycle court (3 à 4 mois), rendements élevés, grains faciles à conserver et à moudre. Il s'est généralisé dans les hautes terres de l'Ouest et du Nord-Ouest, où il est devenu la base de l'alimentation (maïs grillé, \emph{corn-fufu}), ainsi que dans l'Adamaoua.
\item \textbf{Le manioc} : rustique, il tolère les sols pauvres et les sécheresses, se conserve en terre pendant des mois (véritable « garde-manger » vivant), et ses feuilles se mangent aussi. Il est devenu la culture vivrière reine de la zone forestière du Centre, du Sud, de l'Est et du Littoral (\emph{bâton de manioc}, \emph{water-fufu}).
\item \textbf{La pomme de terre} : introduite plus tardivement, elle s'est parfaitement adaptée à la fraîcheur des hautes terres de l'Ouest et du Nord-Ouest (Bafoussam, Dschang, Santa, Foumbot), où elle est devenue à la fois vivrière et commerciale, écoulée jusque dans les grandes villes du Sud.
\item \textbf{Le riz} : culture de plaine inondable, il s'est développé dans le Nord, notamment dans la plaine du Logone et autour de Yagoua et de Kaélé, avec l'appui de la SEMRY (Société d'expansion et de modernisation de la riziculture de Yagoua), créée en 1971. Il répond à la demande croissante des villes.
\end{itemize}

\courssub{Les cultures de rente : de l'introduction coloniale à l'appropriation paysanne}

Les cultures de rente ont une histoire différente : elles ont été introduites à l'époque coloniale (allemande puis française et britannique) pour approvisionner les industries européennes. Mais leur destin a varié : certaines sont restées des cultures de plantations, d'autres ont été \emph{appropriées} par les paysans eux-mêmes.

\begin{center}
\begin{tabularx}{\linewidth}{|l|X|X|}
\hline
\rowcolor{popblueL} \textbf{Culture} & \textbf{Zone principale d'adoption} & \textbf{Forme dominante de production} \\
\hline
Cacao & Forêt du Centre, du Sud, de l'Est, du Sud-Ouest & Paysanne (petites plantations familiales) \\
\hline
Café (arabica et robusta) & Hautes terres de l'Ouest (arabica), Littoral et Sud-Ouest (robusta) & Paysanne \\
\hline
Coton & Savanes du Nord (Nord, Extrême-Nord, Adamaoua) & Paysanne, encadrée par la SODECOTON \\
\hline
Arachide & Savanes du Nord et de l'Adamaoua & Paysanne (vivrière et de rente) \\
\hline
Palmier à huile & Zone forestière humide (Littoral, Sud-Ouest) & Mixte : agro-industries (SOCAPALM) et villageois \\
\hline
Hévéa & Littoral, Sud, Sud-Ouest & Mixte : plantations (HEVECAM) et villageois \\
\hline
Banane d'exportation & Littoral (plaine de la Mungo, Penja) & Agro-industrielle (PHP) et planteurs villageois \\
\hline
\end{tabularx}
\end{center}

\begin{exemplebox}[Le cacao, une réussite paysanne]
Le Cameroun figure parmi les \textbf{premiers producteurs mondiaux de cacao} (ordre de grandeur : 5\textsuperscript{e} ou 6\textsuperscript{e} place mondiale, avec une production d'environ \num{300000} tonnes par an selon les campagnes — chiffre indicatif, variable d'une année à l'autre). Point essentiel : cette production est \textbf{dominée par les petits planteurs paysans}, qui cultivent des vergers familiaux de 1 à 5 hectares, et non par de grandes plantations. Le cacao a transformé l'économie des villages forestiers du Centre et du Sud : c'est lui qui a payé les maisons en dur, les frais de scolarité et les tôles ondulées de générations de familles.
\end{exemplebox}

\courssub{Les facteurs d'adoption : pourquoi une culture s'impose-t-elle ?}

L'observation des réussites (cacao, coton, maïs) et des échecs permet de dégager quatre grands facteurs d'adoption.

\begin{enumerate}
\item \textbf{La demande du marché} : il faut un débouché. Le cacao et le café trouvaient preneur à l'exportation ; le maïs et la pomme de terre trouvent preneur dans les villes camerounaises en croissance rapide.
\item \textbf{Des prix rémunérateurs} : le paysan compare le revenu espéré à celui de ses cultures anciennes. Quand le prix du cacao chute durablement, les jeunes délaissent les vergers ; quand il remonte, on replante.
\item \textbf{L'encadrement} : la présence d'organismes qui fournissent plants, intrants, formation et crédit a été décisive — la SODECOTON pour le coton, la SEMRY pour le riz, les coopératives et les offices pour le cacao et le café.
\item \textbf{L'adaptation au milieu} : une culture ne s'impose durablement que là où le climat et le sol lui conviennent. C'est le facteur écologique, et il est non négociable.
\end{enumerate}

\begin{attention}[Erreur fréquente : mal placer le coton]
Ne place \textbf{jamais} le coton en zone forestière du Sud ! Le coton est une culture des \cle{savanes du Nord} (régions du Nord, de l'Extrême-Nord et de l'Adamaoua). Il exige une \textbf{saison sèche bien marquée} : l'excès d'humidité favorise les maladies et pourrit les capsules, et la récolte doit se faire au sec. Symétriquement, le cacao exige chaleur et humidité toute l'année : il ne peut pas pousser dans les savanes sèches. Au Baccalauréat, confondre les localisations des cultures de rente coûte des points sur la carte comme dans le commentaire.
\end{attention}

\courssub{Les conséquences de l'adoption des nouvelles cultures}

L'adoption massive des cultures de rente a bouleversé les campagnes camerounaises en profondeur.

\begin{itemize}
\item \textbf{La monétarisation des campagnes} : l'argent est entré dans les villages. Le paysan vend une partie de sa récolte, achète des biens manufacturés, paie l'école, investit. L'économie de subsistance pure recule.
\item \textbf{La spécialisation régionale} : chaque région se profile par ses cultures dominantes — cacao au Centre-Sud, café à l'Ouest, coton au Nord, banane au Littoral. Cela crée des paysages agricoles contrastés et des complémentarités commerciales entre régions (un facteur d'intégration nationale).
\item \textbf{La dépendance aux cours mondiaux} : revers de la médaille, le revenu paysan dépend désormais de prix fixés à Londres ou à New York. L'effondrement des cours du café et du cacao à la fin des années 1980 a plongé des régions entières dans la crise. La spécialisation rend vulnérable.
\item \textbf{Des tensions foncières et alimentaires} : la place donnée aux cultures de rente réduit parfois celle des cultures vivrières, et la valeur des terres plantées en cacao attise les conflits fonciers.
\end{itemize}

\courssub{Étude de cas 1 : la diffusion du cacao en zone forestière du Centre et du Sud}

Introduit sous l'administration allemande au tout début du XX\textsuperscript{e} siècle, le cacao trouve dans la forêt équatoriale du Centre et du Sud des conditions idéales : pluies abondantes et régulières (environ \num{1500} à \num{2000}~mm par an), températures élevées toute l'année, sols fertiles. La diffusion suit un schéma classique : quelques planteurs pionniers, souvent des salariés de retour au village ou des notables, plantent les premiers vergers dans les années 1920-1930 ; leurs revenus visibles convainquent les voisins ; après l'indépendance, la « cacaoculture » devient la base de l'économie rurale de toute la zone forestière, d'Obala à Sangmélima. Le cacao s'intègre au système traditionnel : on le plante souvent sur d'anciens brûlis, sous l'ombre des grands arbres conservés, en association avec le manioc et le plantain les premières années. C'est le modèle de l'\emph{agriculture de plantation villageoise}, typiquement camerounais.

\courssub{Étude de cas 2 : le coton dans les savanes du Nord}

Dans les savanes soudaniennes du Nord, la saison sèche longue (octobre à avril) interdit le cacao mais convient au coton. Introduit à l'époque coloniale, le coton ne décolle vraiment qu'après l'indépendance grâce à un encadrement dense : la SODECOTON (créée en 1974) fournit aux paysans semences, engrais et pesticides à crédit, collecte le coton graine au village, l'égrenne dans ses usines (Garoua, Maroua, Kaélé...) et garantit un prix d'achat fixé avant les semis. Ce contrat de filière sécurise le paysan : c'est la clé de l'adoption. Le coton est devenu la « richesse blanche » du Grand Nord, première source de revenus monétaires de centaines de milliers de familles et l'un des premiers produits d'exportation agricole du pays (ordre de grandeur : environ \num{300000} tonnes de coton graine par campagne, chiffre indicatif).

\begin{exoresolu}[Expliquer une adoption]
\textbf{Énoncé.} À partir de l'étude de cas du coton, dégage en quatre arguments rédigés les conditions qui ont permis l'adoption du coton par les paysans du Nord-Cameroun.
\tcblower
\textbf{Solution détaillée.}
\begin{enumerate}
\item \emph{Adaptation au milieu} : les savanes du Nord offrent au coton la saison sèche prolongée et l'ensoleillement dont il a besoin, là où la zone forestière humide lui serait hostile.
\item \emph{Demande du marché} : le coton trouve un débouché sûr, à la fois pour l'exportation de la fibre et pour l'huilerie locale à partir des graines.
\item \emph{Encadrement} : la SODECOTON fournit intrants à crédit, encadrement technique, collecte au village et égrenage, ce qui lève les obstacles matériels qui auraient découragé les paysans isolés.
\item \emph{Prix rémunérateur et garanti} : le prix d'achat est annoncé avant les semis et le paiement est assuré, ce qui sécurise le revenu et rend le coton plus attractif que les seules cultures vivrières.
\end{enumerate}
Conclusion : l'adoption résulte de la rencontre entre un milieu favorable, un marché et un encadrement — jamais d'un seul facteur.
\end{exoresolu}

\courssub{Cartographier les cultures du Cameroun}

\begin{methode}[Légender une carte de répartition agricole]
\begin{enumerate}
\item \textbf{Choisir des figurés simples et distincts} : une couleur ou un pictogramme par culture ; jamais deux figurés proches pour deux cultures différentes.
\item \textbf{Placer les figurés dans les bonnes zones} : vérifier la cohérence écologique (cacao en forêt humide, coton en savane sèche, pomme de terre sur les hautes terres fraîches).
\item \textbf{Rédiger une légende ordonnée} : regrouper d'abord les cultures vivrières, puis les cultures de rente ; chaque figuré apparaît dans la légende exactement comme sur la carte.
\item \textbf{Donner un titre précis} : thème + espace + date (ex. : « Les principales cultures du Cameroun vers 2020 »), plus l'orientation (flèche du nord) et, si possible, une échelle indicative.
\end{enumerate}
\end{methode}

\begin{popfigure}
\begin{center}
\includegraphics[width=\linewidth,height=10.5cm,keepaspectratio]{N04/cmr_map.jpg}
\end{center}
\legende{Fond de carte pour situer les principales cultures du Cameroun (relief, régions, villes, routes ; légende en anglais). On y lit les trois grands ensembles du cours : la zone forestière humide du Sud (Sud, Centre, Littoral, Est : cacao, manioc, banane, palmier à huile) ; les hautes terres et savanes de transition de l'Ouest, du Nord-Ouest et de l'Adamaoua (café, maïs, pomme de terre) ; les savanes sèches du Nord et de l'Extrême-Nord (coton, arachide, mil). Les cultures sont à placer par l'élève. \\ Source : Wikimedia Commons, CIA, domaine public.}
\end{popfigure}

\begin{popfigure}
\begin{center}
\begin{tikzpicture}
\begin{axis}[
  ybar, bar width=22pt,
  width=0.92\linewidth, height=6.2cm,
  ymin=0, ymax=420,
  ylabel={Production (milliers de tonnes, ordre de grandeur)},
  symbolic x coords={Cacao, Coton, Banane dessert, Café},
  xtick=data,
  nodes near coords,
  every node near coord/.append style={font=\scriptsize},
  ymajorgrids=true, grid style={popline!40},
  axis line style={popdark},
  tick label style={font=\small},
  label style={font=\small},
]
\addplot[fill=popblue!70, draw=popblue] coordinates {(Cacao,290) (Coton,300) (Banane dessert,250) (Café,35)};
\end{axis}
\end{tikzpicture}
\end{center}
\legende{\textbf{Production indicative des principales cultures de rente du Cameroun} (moyennes récentes, ordres de grandeur en milliers de tonnes ; la banane dessert concerne l'exportation via PHP ; le café a fortement reculé depuis les années 1980). Sources : INS, ministère de l'Agriculture — valeurs variables selon les campagnes.}
\end{popfigure}

\begin{experience}[TP guidé : légender une carte de répartition des cultures]
\textbf{Matériel} : fond de carte vierge du Cameroun (limites, principales villes), crayons de couleur, atlas ou manuel.
\begin{enumerate}
\item Trace la limite approximative entre la zone forestière (Sud) et les savanes (Nord) : elle passe grossièrement entre le 6\textsuperscript{e} et le 7\textsuperscript{e} degré de latitude Nord.
\item Choisis tes figurés : par exemple un carré vert pour le cacao, un rond brun pour le café, un triangle blanc pour le coton, un point jaune pour le maïs, une croix pour le manioc, un losange pour la banane, une étoile pour le palmier à huile.
\item Place chaque figuré dans sa zone : cacao et manioc au Centre-Sud-Est ; café et maïs à l'Ouest et au Nord-Ouest ; banane et palmier au Littoral ; coton au Nord et à l'Adamaoua.
\item Rédige la légende en deux blocs ordonnés (vivrières / de rente), ajoute le titre, la flèche du nord et une échelle indicative.
\item Vérification croisée avec un camarade : chaque figuré de la carte figure-t-il dans la légende ? Chaque culture est-elle dans une zone écologiquement cohérente ?
\end{enumerate}
\textbf{Interprétation attendue} : la carte révèle la complémentarité Nord-Sud des agricultures camerounaises, fondement des échanges intérieurs et de l'intégration nationale.
\end{experience}

\begin{aretenir}[L'essentiel de la section]
\begin{itemize}
\item L'\cle{adoption} est le processus d'introduction, d'expérimentation puis de diffusion d'une culture exogène ou améliorée dans les systèmes paysans.
\item Cultures vivrières adoptées : \textbf{maïs} (Ouest), \textbf{manioc} (Sud forestier), \textbf{pomme de terre} (hautes terres), \textbf{riz} (Nord, SEMRY).
\item Cultures de rente : \textbf{cacao} et \textbf{café} (Sud et Ouest, production paysanne), \textbf{coton} (savanes du Nord, SODECOTON), \textbf{arachide}, \textbf{palmier à huile}, \textbf{hévéa}, \textbf{banane dessert}.
\item Quatre facteurs d'adoption : marché, prix rémunérateurs, encadrement, adaptation au milieu.
\item Conséquences : monétarisation des campagnes, spécialisation régionale, mais dépendance aux cours mondiaux.
\item Piège majeur : le coton appartient aux savanes, le cacao à la forêt — ne jamais les inverser.
\end{itemize}
\end{aretenir}

\begin{exercice}[Pour t'entraîner]
\begin{enumerate}
\item Définis « culture de rente » et donne trois exemples camerounais avec leur zone de production.
\item Explique pourquoi le manioc s'est diffusé plus facilement que le riz dans la zone forestière.
\item Sur un fond de carte vierge, place le cacao, le café, le coton et la banane, puis rédige une légende conforme à la méthode.
\item Montre, en cinq lignes, en quoi l'adoption des cultures de rente est à la fois un progrès et une source de fragilité pour les campagnes camerounaises.
\end{enumerate}
\end{exercice}

\begin{corrige}[Exercice — éléments de réponse]
\begin{enumerate}
\item Culture destinée principalement à la vente ; cacao (Centre-Sud), café (Ouest), coton (Nord).
\item Le manioc est rustique, sans besoin d'eau aménagée ni d'intrants, et se conserve en terre ; le riz exige des plaines inondables aménagées, coûteuses, absentes de la forêt du Sud.
\item Cacao au Centre-Sud, café à l'Ouest, coton au Nord, banane au Littoral ; légende en deux blocs (vivrières / rente), titre daté, flèche du nord.
\item Progrès : revenus monétaires, scolarisation, équipement des villages, insertion dans les échanges. Fragilité : dépendance aux cours mondiaux, réduction possible des cultures vivrières, risque de crise en cas d'effondrement des prix (années 1980-1990).
\end{enumerate}
\end{corrige}

\coursec{L'agriculture itinérante sur brûlis et l'agriculture vivrière}

Nous venons de voir que les paysans camerounais ont adopté de nouvelles cultures (maïs, manioc, arachide, cacaoyer...) venues d'ailleurs. Mais ces cultures nouvelles se sont insérées dans des \emph{systèmes} agricoles anciens, profondément adaptés aux milieux. Deux de ces systèmes dominent encore aujourd'hui l'agriculture traditionnelle camerounaise : l'\cle{agriculture itinérante sur brûlis} dans la zone forestière du Sud, et l'\cle{agriculture vivrière} qui constitue la finalité première de la plupart des exploitations. C'est à ces deux systèmes, complémentaires et souvent imbriqués, qu'est consacrée cette section.

\courssub{L'agriculture itinérante sur brûlis : définition et mécanisme}

\textbf{Une observation pour commencer.} Parcourons en pensée la piste qui relie un village boulou à la route Yaoundé--Ebolowa, dans la région du Sud. De part et d'autre, on ne voit pas de grands champs permanents comme dans les plaines de l'Adamaoua, mais une mosaïque : ici une parcelle de manioc et de bananiers plantains au milieu de troncs noircis, là une jeune brousse où la forêt reprend ses droits, plus loin la forêt dense intacte. Comment expliquer ce paysage en patchwork ? C'est toute la logique de l'agriculture itinérante sur brûlis.

\begin{definition}[Agriculture itinérante sur brûlis]
L'\cle{agriculture itinérante sur brûlis} est un système de culture dans lequel la végétation naturelle d'une parcelle forestière est \textbf{coupée puis brûlée} avant les semis, les cendres fertilisant le sol ; la parcelle, cultivée pendant quelques années seulement, est ensuite \textbf{abandonnée} (mise en jachère) lorsque ses rendements s'épuisent, et le paysan défriche une nouvelle parcelle. L'agriculture est dite \emph{itinérante} parce que les champs se déplacent dans l'espace au fil du temps.
\end{definition}

\textbf{Les étapes du cycle.} Le cycle complet se déroule en six étapes bien ordonnées, que l'on peut observer chaque année dans les terroirs du Sud-Cameroun :

\begin{enumerate}
\item \textbf{Le choix de la parcelle} : le chef de famille repère une portion de forêt (ou de vieille jachère) dont la végétation annonce un sol encore fertile ; le droit coutumier règle l'accès à la terre.
\item \textbf{L'abattage} : pendant la saison sèche, les arbres sont coupés à la hache et à la machette (les plus gros troncs sont parfois simplement ébranchés ou laissés sur place).
\item \textbf{Le brûlis} : après séchage, la végétation coupée est incendiée. Le feu détruit les adventices et les graines d'herbes, et les cendres répandues enrichissent le sol.
\item \textbf{Les semis} : dès les premières pluies, on plante, souvent en \emph{association culturale} (polyculture), banane plantain, manioc, taro, macabo, arachide, maïs...
\item \textbf{Les cultures} : la parcelle est exploitée pendant \textbf{2 à 3 ans} en général, avec des entretiens limités (sarclage à la houe).
\item \textbf{L'abandon et la jachère} : quand les rendements chutent et que les herbes envahissent le champ, la parcelle est abandonnée ; la forêt secondaire recolonise le sol pendant une jachère longue, tandis que le paysan ouvre un nouveau champ ailleurs.
\end{enumerate}

\begin{popfigure}
\begin{center}
\begin{tikzpicture}[scale=1]
\foreach \i/\txt in {0/{1. Choix de la parcelle},60/{2. Abattage (saison sèche)},120/{3. Brûlis (feu)},180/{4. Semis (premières pluies)},240/{5. Cultures (2 à 3 ans)},300/{6. Abandon, jachère (10 à 20 ans)}}{
\node[draw=popblue, fill=popblueL, rounded corners=2pt, text width=3.1cm, align=center, font=\small] at (\i:3.6) {\txt};
}
\foreach \a/\b in {30/330,90/30,150/90,210/150,270/210,330/270}{
\draw[-{Stealth[length=3mm]}, poporange, very thick] (\a:2.1) arc (\a:\b:2.1);
}
\node[align=center, font=\small\itshape, text=popdark] at (0,0) {Cycle de l'agriculture\\itinérante sur brûlis};
\end{tikzpicture}
\end{center}
\legende{Schéma : les six étapes du cycle de l'agriculture itinérante sur brûlis en zone forestière du Sud-Cameroun. Le cycle complet dure de 12 à plus de 20 ans, dont seulement 2 à 3 ans de culture effective.}
\end{popfigure}

\begin{savaistu}[Les cendres, un engrais naturel]
Les cendres du brûlis ne sont pas un simple résidu : elles apportent au sol de la \textbf{potasse} (K), ainsi que du calcium et du phosphore, c'est-à-dire un \emph{engrais minéral naturel}. Sous forêt dense humide, les éléments nutritifs sont stockés surtout dans la végétation et non dans le sol (les pluies abondantes lessivent en profondeur les sols ferrallitiques) : brûler la biomasse, c'est donc \emph{transférer} en quelques heures vers le sol des nutriments accumulés pendant des décennies par la forêt. Voilà pourquoi les rendements sont bons la première année... puis s'effondrent quand ce stock s'épuise.
\end{savaistu}

\courssub{Localisation, cultures associées et conditions de viabilité}

\textbf{Localisation.} L'agriculture itinérante sur brûlis est le système caractéristique de la \textbf{zone forestière du Sud} du Cameroun : les régions du \cle{Sud}, de l'\cle{Est}, du \cle{Centre} et la partie forestière du \cle{Littoral}. C'est le domaine de la forêt dense humide (guinéo-congolaise), des sols ferrallitiques pauvres et acides, et d'une population longtemps clairsemée. On le retrouve, sous des formes voisines, dans les zones de savane arborée soudanienne (régions du Nord et de l'Extrême-Nord), où le brûlis concerne alors les herbes et les arbustes de savane.

\textbf{Les cultures associées.} Le champ de brûlis est rarement occupé par une seule plante : c'est une \cle{polyculture} en \emph{association culturale}, qui imite la structure étagée de la forêt. On y trouve typiquement :
\begin{itemize}
\item la \textbf{banane plantain} et le \textbf{manioc}, piliers de l'alimentation en zone forestière ;
\item le \textbf{taro} et le \textbf{macabo} (deux tubercules de la famille des Aracées, genres \emph{Colocasia} et \emph{Xanthosoma}) ;
\item l'\textbf{arachide}, le \textbf{maïs}, le \textbf{gombo}, des légumes-feuilles (folong, ndolé...) et des plantes condimentaires.
\end{itemize}
Cette association présente de réels avantages : elle protège le sol du ruissellement, limite la propagation des maladies et des ravageurs, étale les récoltes dans le temps et diversifie l'alimentation familiale.

\textbf{Conditions de viabilité.} Ce système n'est durable que sous deux conditions strictes :
\begin{enumerate}
\item une \textbf{faible densité de population} (de l'ordre de quelques habitants au km\textsuperscript{2} ; historiquement, l'Est et le Sud comptaient souvent moins de 10 hab/km\textsuperscript{2}) ;
\item une \textbf{large disponibilité de terres forestières}, permettant des jachères longues de 10 à 20 ans, seules capables de reconstituer la fertilité du sol.
\end{enumerate}

\begin{exemplebox}[Un exemple chiffré : le temps de culture et le temps de jachère]
En zone forestière du Sud-Cameroun, une parcelle de brûlis n'est cultivée que \textbf{2 à 3 ans} avant d'être laissée en jachère pendant \textbf{10 à 20 ans}. Autrement dit, pour nourrir durablement une famille, il faut disposer d'un \emph{parcours} de 6 à 10 parcelles en rotation. Si la jachère tombe à 3 ou 4 ans (pression foncière), la forêt n'a pas le temps de se reconstituer : la fertilité n'est plus restaurée et le système bascule dans la dégradation.
\end{exemplebox}

\courssub{Limites et conséquences de l'agriculture itinérante sur brûlis}

Tant que les jachères restent longues, le système est équilibré : la forêt régénère le sol et le bilan carbone est à peu près neutre. Mais il devient destructeur lorsque l'équilibre se rompt.

\begin{itemize}
\item \textbf{Déforestation} : chaque nouveau brûlis rogne la forêt. Autour des villes (Yaoundé, Douala), le long des axes routiers et dans les zones de front pionnier (Est, Sud), le raccourcissement des jachères transforme la forêt dense en forêt secondaire, puis en savane arbustive dégradée. Le Cameroun perd chaque année une fraction de son couvert forestier (ordre de grandeur : de l'ordre de 0,5 \% à 1 \% par an selon les sources et les périodes, chiffre à manier avec prudence).
\item \textbf{Baisse des rendements} : quand la jachère se raccourcit sous l'effet de la croissance démographique, le sol n'a plus le temps de reconstituer son stock de nutriments ; les rendements du manioc, du plantain ou du maïs chutent, poussant à défricher toujours plus : c'est un \emph{cercle vicieux}.
\item \textbf{Érosion des sols} : sur les parcelles récemment brûlées, le sol nu est exposé aux pluies violentes de la zone équatoriale ; le ruissellement emporte la fine couche d'humus, surtout sur les versants (Ouest forestier, hautes terres bamiléké où le brûlis de savane sévit aussi).
\item \textbf{Autres effets} : perte de biodiversité, destruction des habitats de la faune, émissions de carbone, conflits fonciers entre agriculteurs, exploitants forestiers et agro-industriels.
\end{itemize}

\begin{attention}[Erreur fréquente : le brûlis, une « destruction gratuite » ?]
Il est faux de présenter le brûlis comme une pratique « primitive » ou irrationnelle. C'est une \textbf{technique rationnelle et adaptée} aux sols pauvres de la zone forestière : sans engrais minéraux disponibles, le feu est le seul moyen de transférer rapidement les nutriments de la végétation vers le sol, et il assainit la parcelle (adventices, parasites). Le brûlis ne devient \emph{nuisible} que lorsque la \textbf{pression démographique} raccourcit les jachères en dessous du seuil de régénération. Au Baccalauréat, évite donc tout jugement de valeur : analyse les \emph{conditions} de viabilité du système.
\end{attention}

\courssub{L'agriculture vivrière : nourrir d'abord la famille}

\textbf{Une scène de marché.} Samedi matin au marché de Mbalmayo : des femmes déballent des sacs de manioc, des régimes de plantain, des cosses d'arachide, des tubercules de macabo. Une partie de cette production vient des champs familiaux situés à quelques kilomètres ; l'essentiel de la récolte, lui, n'a jamais quitté le village : il a été consommé sur place. C'est le cœur de l'agriculture vivrière.

\begin{definition}[Agriculture vivrière]
L'\cle{agriculture vivrière} est une agriculture dont la production est destinée \textbf{en priorité à la consommation de la famille} qui la produit (autoconsommation) ; seuls les \textbf{surplus} éventuels sont vendus sur les marchés locaux pour se procurer un revenu monétaire (sel, huile, savon, scolarité des enfants...). Elle s'oppose, conceptuellement, à l'agriculture commerciale (ou de rente), tournée vers la vente et l'exportation.
\end{definition}

\textbf{Les cultures principales.} L'agriculture vivrière camerounaise repose sur des cultures vivrières de base, réparties selon les grandes zones agro-écologiques :
\begin{itemize}
\item en \textbf{zone forestière} (Sud, Centre, Est, Littoral) : manioc, banane plantain, taro, macabo, igname, maïs, arachide ;
\item en \textbf{zone de savane soudanienne} (Nord, Extrême-Nord, Adamaoua) : mil, sorgho, maïs, niébé, arachide, fonio ;
\item dans les \textbf{hautes terres de l'Ouest} : maïs, haricot, pomme de terre, taro, banane.
\end{itemize}
Le \cle{manioc} est aujourd'hui la première culture vivrière du pays (plusieurs millions de tonnes par an ; ordre de grandeur : environ 5 à 6 millions de tonnes ces dernières années, à manier avec prudence), devant le plantain, le maïs et le mil-sorgho.

\textbf{Autoconsommation et surplus.} Dans l'exploitation vivrière typique, la plus grande part de la récolte est consommée par le ménage ou stockée (manioc transformé en bâtons, en farine ou en gari ; mil conservé en greniers). Les surplus sont écoulés sur les \emph{marchés locaux} périodiques (marchés hebdomadaires villageois) ou vendus aux revendeuses qui alimentent les villes. Ce revenu, modeste mais vital, finance les dépenses monétaires de la famille. On parle parfois d'\emph{agriculture vivrière marchande} lorsque la part vendue devient significative : la frontière avec l'agriculture commerciale est donc graduelle, pas absolue.

\begin{aretenir}[Le rôle central des femmes]
L'agriculture vivrière camerounaise repose majoritairement sur le \textbf{travail des femmes} : elles assurent l'essentiel des semis, des sarclages, des récoltes, du transport (portage sur la tête), de la transformation (gari, huile, feuilles de manioc) et de la commercialisation des surplus sur les marchés. On estime couramment qu'elles fournissent de l'ordre de \textbf{60 à 80 \%} de la main-d'œuvre agricole vivrière (ordre de grandeur prudent). Pourtant, elles disposent rarement du titre de propriété foncière coutumier, qui revient aux hommes, et leur accès au crédit et aux intrants reste limité : c'est un enjeu majeur d'équité et de modernisation agricole.
\end{aretenir}

\courssub{TP guidé : construire un croquis de terroir villageois en zone forestière}

\begin{methode}[Construire un croquis géographique]
\begin{enumerate}
\item \textbf{Dégager l'essentiel} : à partir du document (carte, photo, récit), identifier les 4 ou 5 éléments qui expriment l'organisation de l'espace (ici : village, champs, jachères, forêt, piste, cours d'eau). Éliminer le détail anecdotique.
\item \textbf{Choisir des figurés conventionnels} : surface colorée ou hachurée pour les unités paysagères, points ou petits carrés pour les habitats, traits pour les axes, traits bleus ondulés pour l'hydrographie.
\item \textbf{Hiérarchiser l'information} : l'élément principal doit être le plus visible ; ne pas tout mettre au même niveau.
\item \textbf{Soigner le cadre} : titre explicite en haut, flèche du nord, échelle indicative si possible.
\item \textbf{Rédiger la légende} : numérotée, placée sous le croquis, avec des intitulés précis (jamais « zone 1 » sans explication).
\end{enumerate}
\end{methode}

\begin{experience}[TP : croquis d'un terroir villageois en zone forestière du Sud]
\textbf{Matériel} : feuille blanche, crayons de couleur, règle ; photographie aérienne ou description d'un terroir (fournie par l'enseignant).
\textbf{Protocole} : (1) repérer sur le document le village et la piste ; (2) délimiter les champs en culture (brûlis récents), les jachères jeunes (brousse) et anciennes (forêt secondaire), la forêt dense ; (3) reporter ces unités sur le croquis avec des figurés simples ; (4) rédiger titre, orientation et légende numérotée.
\textbf{Observations attendues} : les champs se concentrent près du village et le long de la piste ; les jachères forment une couronne intermédiaire ; la forêt dense occupe les marges du terroir ; le cours d'eau sert de repère et de limite.
\textbf{Interprétation} : le terroir est organisé en \emph{anneaux concentriques} autour du village, reflet du cycle de l'agriculture itinérante : plus une parcelle est loin, moins elle est fréquemment cultivée.
\end{experience}

\begin{popfigure}
\begin{center}
\begin{tikzpicture}[scale=0.95]
% forêt dense (fond)
\fill[popgreen!55] (-5.2,-3.2) rectangle (5.2,3.4);
% jachère ancienne
\fill[popgreen!30] (-3.6,-2.2) rectangle (3.4,2.4);
% jachère jeune
\fill[popgold!45] (-2.4,-1.4) rectangle (2.2,1.6);
% champ en culture
\fill[poporange!55] (-1.1,-0.7) rectangle (1.0,0.6);
% village
\foreach \x/\y in {-0.5/1.15,0.1/1.3,0.6/1.1,-0.1/0.95}{
\draw[fill=popdark] (\x,\y) rectangle ++(0.22,0.18);
\draw[popdark] (\x-0.03,\y+0.18) -- (\x+0.11,\y+0.34) -- (\x+0.25,\y+0.18);
}
% piste
\draw[popdark, thick, dashed] (-5.2,1.6) -- (-1,1.1) -- (0.1,1.15) -- (2,0.9) -- (5.2,1.3);
% rivière
\draw[popblue, very thick] (-5.2,-2.6) .. controls (-2,-2.0) and (-1,-2.9) .. (1.5,-2.3) .. controls (3,-1.9) and (4,-2.7) .. (5.2,-2.2);
% nord
\draw[-{Stealth[length=3mm]}, very thick] (4.5,2.5) -- (4.5,3.2);
\node[above] at (4.5,3.2) {\small N};
% labels
\node[font=\small\itshape] at (0,0) {champ en culture};
\node[font=\small\itshape] at (0.05,1.55) {village};
\node[font=\small\itshape] at (2.9,1.9) {jachère jeune};
\node[font=\small\itshape] at (-3.0,2.7) {jachère ancienne};
\node[font=\small\itshape, text=white] at (-3.4,-2.8) {forêt dense};
\node[font=\small\itshape, text=popblue] at (3.4,-2.55) {rivière};
\node[font=\small\itshape, rotate=2] at (-3.3,1.45) {piste};
\end{tikzpicture}
\end{center}
\legende{Croquis modèle : un terroir villageois en zone forestière du Sud-Cameroun. Légende : 1. village (cases) ; 2. piste (trait discontinu) ; 3. champ en culture sur brûlis récent (orange) ; 4. jachère jeune, brousse de 1 à 5 ans (jaune) ; 5. jachère ancienne, forêt secondaire (vert clair) ; 6. forêt dense (vert foncé) ; 7. rivière (bleu). Échelle indicative : le terroir s'étend sur quelques kilomètres.}
\end{popfigure}

\courssub{Exercice résolu et synthèse}

\begin{exoresolu}[Commentaire de croquis : le terroir de Nkometou (Centre)]
Énoncé. Un élève a réalisé le croquis d'un terroir au nord de Yaoundé : village central, deux champs de manioc-plantain proches, trois jachères d'âges croissants, forêt dense en périphérie, piste menant à la route nationale. La jachère moyenne y est passée de 15 ans (années 1980) à 5 ans aujourd'hui. \textbf{Question} : montrer que ce terroir illustre à la fois la rationalité et les limites actuelles de l'agriculture itinérante sur brûlis.
\tcblower
\textbf{Solution détaillée.}
\emph{1. Un système rationnel.} Le croquis montre une rotation organisée : les champs en culture côtoient des jachères d'âges différents, ce qui prouve que le paysan ne défriche pas au hasard mais gère un \emph{parcours} de parcelles. Le brûlis, en zone de sols ferrallitiques pauvres, est le seul moyen paysan de restituer rapidement au sol les nutriments stockés par la forêt (potasse des cendres) ; la polyculture (manioc, plantain, macabo) protège le sol et sécurise l'alimentation.
\emph{2. Des limites apparues avec la pression démographique.} Le raccourcissement de la jachère de 15 à 5 ans signifie que la forêt secondaire n'a plus le temps de reconstituer la fertilité : les rendements baissent, les adventices s'installent, et le paysan est contraint de défricher plus loin, aux dépens de la forêt dense visible en périphérie du croquis.
\emph{3. Conclusion.} Le terroir de Nkometou illustre un système longtemps équilibré, devenu fragile sous l'effet de la croissance démographique et de la proximité urbaine (marché de Yaoundé) : il appelle des améliorations (jachères améliorées, engrais, cultures pérennes) qui seront étudiées avec l'agriculture de seconde génération.
\end{exoresolu}

\begin{exercice}[Pour t'entraîner]
1. Définis en deux phrases l'agriculture vivrière et explique pourquoi elle ne s'oppose pas totalement à l'agriculture commerciale.
2. À l'aide du croquis de terroir ci-dessus, rédige une légende numérotée complète, puis explique en cinq lignes pourquoi les champs se situent près du village.
3. Calcule : si une famille cultive 2 ans chaque parcelle et laisse une jachère de 12 ans, combien de parcelles doit comporter au minimum son parcours de rotation ?
\end{exercice}

\begin{corrige}[Exercice n°3]
Pendant qu'une parcelle est cultivée 2 ans puis reposée 12 ans, le cycle complet d'une parcelle dure $2 + 12 = 14$ ans. Pour disposer chaque année d'une nouvelle parcelle à défricher, il faut :
\[ \frac{14 \text{ ans}}{2 \text{ ans de culture}} = 7 \text{ parcelles}. \]
Le parcours de rotation doit donc comporter au minimum \textbf{7 parcelles} (une en culture, six à des stades divers de jachère). C'est ce calcul simple qui montre pourquoi le système exige de vastes espaces forestiers et s'effondre quand la terre vient à manquer.
\end{corrige}

\begin{aretenir}[L'essentiel de la section]
\begin{itemize}
\item L'\cle{agriculture itinérante sur brûlis} (zone forestière du Sud, de l'Est, du Centre et du Littoral forestier) suit un cycle : choix de la parcelle, abattage, brûlis, semis en polyculture (plantain, manioc, taro, macabo, arachide), 2 à 3 ans de culture, puis abandon pour une jachère de 10 à 20 ans.
\item Elle n'est viable qu'à \textbf{faible densité de population} et avec de \textbf{vastes réserves forestières} ; le raccourcissement des jachères entraîne déforestation, baisse des rendements et érosion.
\item L'\cle{agriculture vivrière} vise d'abord l'autoconsommation familiale (manioc, maïs, mil, sorgho, igname, plantain), les surplus étant vendus sur les marchés locaux ; elle repose majoritairement sur le travail des \textbf{femmes}.
\end{itemize}
\end{aretenir}

\coursec{L'agriculture de seconde génération et bilan des améliorations}

Dans la section précédente, nous avons étudié deux systèmes agricoles qui dominent encore les campagnes camerounaises : l'\cle{agriculture itinérante sur brûlis}, fondée sur la jachère longue, et l'\cle{agriculture vivrière}, tournée d'abord vers l'autoconsommation. Toutes deux reposent sur un outillage sommaire et des rendements modestes. Pourtant, si tu traverses aujourd'hui les savanes de la région du Nord en saison cotonnière, ou les collines des Hautes Terres de l'Ouest couvertes de champs de maïs, tu constates que le paysage agricole a changé : parcelles bien alignées, sacs d'engrais, pulvérisateurs, camions de collecte. Une troisième forme d'agriculture est née, que les géographes appellent l'\cle{agriculture de seconde génération}. C'est elle qui couronne les « grandes améliorations » de l'agriculture traditionnelle camerounaise. Cette dernière section la définit, en présente les exemples majeurs, puis dresse le bilan complet des transformations de l'agriculture camerounaise.

\courssub{Définition et caractéristiques de l'agriculture de seconde génération}

\textbf{Intuition.} Imagine un planteur de la région de l'Ouest. Son père cultivait le maïs uniquement pour nourrir la famille, avec une houe et des graines prélevées sur la récolte précédente. Lui, en revanche, achète des semences améliorées à la coopérative, épand de l'engrais NPK, traite son champ contre les insectes, et vend les trois quarts de sa récolte au collecteur qui approvisionne le marché de Bafoussam. Avec l'argent gagné, il paie la scolarité de ses enfants et achète les intrants de la campagne suivante. Il reste un paysan — sa famille fournit l'essentiel du travail — mais il n'est plus un simple paysan vivrier : il est entré dans l'économie de marché.

\begin{definition}[Agriculture de seconde génération]
L'\cle{agriculture de seconde génération} est une agriculture paysanne \textbf{améliorée et modernisée}, qui demeure à dimension \textbf{familiale} (main-d'œuvre familiale, exploitations de petite ou moyenne taille), mais qui utilise des \textbf{intrants modernes} (semences améliorées, engrais, pesticides) et qui est \textbf{intégrée au marché} : elle combine cultures vivrières et cultures de rente, et vend une part croissante de sa production. Elle constitue une étape intermédiaire entre l'agriculture vivrière traditionnelle et l'agriculture de plantation (dite « de troisième génération »).
\end{definition}

Ses caractéristiques essentielles sont les suivantes :

\begin{itemize}
\item \textbf{L'utilisation d'intrants} : semences sélectionnées (maïs hybride, riz amélioré), engrais minéraux (NPK, urée) et organiques, produits phytosanitaires (herbicides, insecticides, fongicides).
\item \textbf{L'encadrement technique} : sociétés de développement (SODECOTON, SEMRY), services du MINADER, ONG et projets agricoles fournissent formation, conseil et suivi.
\item \textbf{L'accès au crédit} : crédit de campagne (intrants avancés puis remboursés à la récolte), microfinance, tontines, comptes épargne-crédit des coopératives.
\item \textbf{La commercialisation organisée} : collecte à domicile ou en points de vente, coopératives et GIC (Groupements d'Initiative Commune), contrats avec les sociétés d'encadrement, circuits périurbains.
\item \textbf{La combinaison des cultures} : une part vivrière (autoconsommation) et une part de rente (vente), ce qui sécurise la famille tout en dégageant un revenu monétaire.
\end{itemize}

\begin{attention}[Ne pas opposer totalement traditionnel et moderne]
Erreur fréquente au Baccalauréat : présenter l'agriculture de seconde génération comme une agriculture « moderne » qui aurait rompu avec la tradition. C'est faux. Elle reste une agriculture \textbf{familiale}, peu mécanisée, pratiquée sur de petites surfaces, souvent en culture associée. Sa modernité est \emph{relative} : intrants, encadrement, marché. C'est une \textbf{étape intermédiaire}, non une rupture. Rédige toujours cette nuance dans tes devoirs.
\end{attention}

\courssub{Les exemples camerounais de l'agriculture de seconde génération}

\textbf{Le coton des savanes du Nord, sous encadrement de la SODECOTON.} La Société de Développement du Coton du Cameroun (SODECOTON), créée en 1974 et basée à Garoua, fournit aux planteurs des régions du Nord, de l'Adamaoua et de l'Extrême-Nord semences, engrais et pesticides à crédit, encadre les champs, puis collecte, égrenne et achète toute la production à un prix fixé avant la campagne. Le coton est ainsi devenu la principale source de revenus monétaires des paysans des savanes septentrionales, avec des productions de coton-graine de l'ordre de \SI{200000}{} à plus de \SI{300000}{\tonne} selon les campagnes (ordre de grandeur, variable selon les années).

\textbf{Le riz de la SEMRY dans l'Extrême-Nord.} La Société d'Expansion et de Modernisation de la Riziculture de Yagoua (SEMRY) aménage des périmètres irrigués dans la plaine du Logone (Yagoua, Maga, Pouss) grâce aux retenues d'eau comme le barrage de Maga. Les riziculteurs y pratiquent deux cycles culturaux par an sur certains périmètres, avec semences améliorées et mécanisation partielle (motoculteurs). Le riz de la SEMRY alimente les marchés de tout le Grand Nord.

\textbf{Le maïs intensif des Hautes Terres de l'Ouest.} Dans les départements des Hauts-Plateaux, du Noun et de la Menoua, le maïs est cultivé en rotation ou en association avec le haricot et la pomme de terre, avec engrais et semences améliorées. Les rendements y dépassent nettement la moyenne nationale. Le surplus est collecté par les GIC et les commerçants pour les marchés de Bafoussam, Douala et Yaoundé.

\textbf{Le maraîchage périurbain de Yaoundé et Douala.} Autour des grandes villes, les bas-fonds sont mis en valeur pour le maraîchage (tomates, poivrons, laitues, gombos, aubergines), souvent de contre-saison grâce à l'eau disponible en saison sèche. La proximité des marchés urbains garantit des débouchés rapides et rémunérateurs.

\begin{exemplebox}[Le coton, moteur monétaire des savanes du Nord]
Grâce à l'encadrement de la SODECOTON, un planteur de la région du Nord peut cultiver \SI{1}{} à \SI{3}{\hectare} de coton sans avancer d'argent : les intrants lui sont fournis à crédit et leur coût est déduit du prix de sa récolte à la livraison. Ce système a fait du coton la principale source de revenus monétaires des paysans des savanes du Grand Nord, finançant la scolarisation, la santé, la construction de cases en dur et l'achat d'équipements (motos, charrettes, radios solaires).
\end{exemplebox}

\courssub{Les acteurs de cette mutation agricole}

L'agriculture de seconde génération est portée par des acteurs nouveaux :

\begin{itemize}
\item \textbf{Les jeunes agriculteurs formés} : issus des lycées agricoles, des CEP (Centres d'Enseignement Pratique) ou des programmes de formation du MINADER, ils appliquent les itinéraires techniques modernes.
\item \textbf{Les coopératives et les GIC} : ils mutualisent l'achat des intrants, la location de matériel et la vente groupée, ce qui renforce le pouvoir de négociation des paysans.
\item \textbf{Les agropasteurs} : notamment dans l'Adamaoua et le Nord, ils combinent élevage et cultures, utilisant le fumier du bétail comme engrais organique — un bel exemple d'intégration agriculture-élevage.
\item \textbf{Les « cadres » investisseurs} : fonctionnaires, commerçants ou retraités urbains qui investissent leurs revenus dans des fermes (maïs, volailles, maraîchage), faisant affluer capitaux et exigences de rentabilité vers les campagnes.
\end{itemize}

\begin{savaistu}[Le maraîchage de contre-saison, une chance pour les jeunes urbains]
Dans les bas-fonds de Yaoundé (Mvolyé, Nkolondom, Etoudi) et de Douala (Bonabéri, PK zones humides), de jeunes citadins pratiquent le maraîchage de contre-saison : pendant la saison sèche, l'eau des bas-fonds permet d'irriguer des cultures qui manquent alors sur les marchés, et qui se vendent donc cher. Sur quelques centaines de mètres carrés seulement, un jeune maraîcher peut dégager un revenu régulier supérieur à celui de nombreux emplois informels. L'agriculture urbaine et périurbaine est devenue un véritable secteur d'insertion professionnelle.
\end{savaistu}

\courssub{La chaîne de valeur de l'agriculture de seconde génération}

La grande nouveauté de ce système est que la production s'inscrit dans une \cle{chaîne de valeur} : chaque étape ajoute de la valeur et génère un revenu qui, réinvesti, relance le cycle suivant.

\begin{popfigure}
\begin{center}
\begin{tikzpicture}[
  box/.style={draw=popblue, thick, fill=popblueL, rounded corners=2pt, minimum width=2.6cm, minimum height=1.1cm, align=center, font=\small},
  arr/.style={-{Stealth[length=3mm]}, very thick, poporange},
  note/.style={font=\scriptsize\itshape, text=popdark, align=center}
]
\node[box, align=center] (intr) at (0,0){\textbf{Intrants}\\semences, engrais,\\pesticides, crédit};
\node[box, align=center] (prod) at (4.2,0){\textbf{Production}\\champ familial\\encadré};
\node[box, align=center] (coll) at (8.4,0){\textbf{Collecte}\\GIC, coopérative,\\SODECOTON, SEMRY};
\node[box, align=center] (march) at (0,-2.6){\textbf{Marché}\\local, urbain,\\exportation};
\node[box, align=center] (rev) at (4.2,-2.6){\textbf{Revenu}\\monétaire\\du producteur};
\node[box, fill=popgreenL, draw=popgreen, align=center] (reinv) at (8.4,-2.6){\textbf{Réinvestissement}\\intrants, équipement,\\scolarité, épargne};
\draw[arr] (intr) -- (prod);
\draw[arr] (prod) -- (coll);
\draw[arr] (coll.south) |- (march.east);
\draw[arr] (march) -- (rev);
\draw[arr] (rev) -- (reinv);
\draw[arr, dashed] (reinv.north) -- (coll.south);
\draw[arr, dashed] (reinv.west) ++(0,0.35) -| (intr.south);
\node[note] at (4.2,1.15) {Le revenu réinvesti boucle le cycle : c'est le moteur de la modernisation paysanne.};
\end{tikzpicture}
\end{center}
\legende{Schéma de la chaîne de valeur de l'agriculture de seconde génération : 1. intrants fournis souvent à crédit ; 2. production familiale encadrée ; 3. collecte organisée ; 4. vente sur les marchés ; 5. revenu monétaire ; 6. réinvestissement qui relance la campagne suivante (flèches pointillées).}
\end{popfigure}

\courssub{Bilan des grandes améliorations : progrès et limites}

\textbf{Les progrès sont réels et mesurables.}

\begin{itemize}
\item \textbf{Hausse des rendements et des productions} : le maïs passe d'environ \SI{1}{\tonne\per\hectare} en système vivrier traditionnel à \SI{3}{\tonne\per\hectare} et plus en système amélioré ; le coton, le riz irrigué et le maraîchage connaissent des gains comparables.
\item \textbf{Autosuffisance alimentaire relative} : le Cameroun couvre l'essentiel de sa consommation en produits de base (maïs, manioc, plantain, mil), même si des importations persistent (riz, blé, poisson).
\item \textbf{Revenus monétaires paysans} : les cultures de rente et la vente des surplus vivriers ont monétisé les campagnes, finançant scolarité, santé et habitat.
\item \textbf{Intégration nationale} : les flux de produits agricoles relient les régions entre elles (maïs de l'Ouest vers Douala, oignons et riz du Nord vers le Sud), tissant la solidarité économique du pays.
\end{itemize}

\textbf{Mais des limites persistent.}

\begin{itemize}
\item \textbf{Faible mécanisation} : la houe et la machette dominent encore ; le tracteur reste rare et concentré dans quelques périmètres.
\item \textbf{Enclavement} : nombre de pistes rurales sont impraticables en saison des pluies, ce qui renchérit le transport et fait perdre une partie des récoltes.
\item \textbf{Fluctuation des prix} : l'effondrement des cours du cacao et du café dans les années 1980-1990 a ruiné des milliers de planteurs ; les prix du maïs ou des oignons varient fortement d'une saison à l'autre.
\item \textbf{Vieillissement des planteurs} : l'exode rural prive les campagnes de bras jeunes ; l'âge moyen des exploitants augmente.
\item \textbf{Changement climatique} : retard ou irrégularité des pluies, sécheresses dans le Grand Nord, nouvelles maladies des plantes fragilisent les systèmes améliorés, très dépendants du calendrier des pluies.
\end{itemize}

\begin{aretenir}[Bilan de la leçon]
L'agriculture traditionnelle camerounaise s'est profondément transformée : outillage amélioré, techniques culturales perfectionnées, adoption de nouvelles cultures. Trois systèmes coexistent aujourd'hui : l'agriculture itinérante sur brûlis (jachère longue, faibles rendements), l'agriculture vivrière (autoconsommation d'abord) et l'agriculture de seconde génération (familiale, mais dotée d'intrants, encadrée et tournée vers le marché). Cette dernière a fait progresser rendements, revenus et autosuffisance alimentaire, sans éliminer les contraintes : mécanisation faible, enclavement, prix instables, vieillissement, climat.
\end{aretenir}

\courssub{Activité de synthèse : comparer les trois systèmes agricoles}

\begin{methode}[Construire un tableau comparatif]
\begin{enumerate}
\item \textbf{Définir des critères communs} : choisis 3 à 5 critères pertinents et identiques pour tous les objets comparés (ici : outillage, techniques, cultures, destination de la production).
\item \textbf{Remplir chaque case avec des faits précis} : un fait court et vérifiable par case, jamais de phrase vague (« c'est mieux » est interdit).
\item \textbf{Vérifier l'équilibre} : chaque système doit être renseigné pour chaque critère ; une case vide est une faute de méthode.
\item \textbf{Conclure par une phrase de synthèse} : dégage la logique d'ensemble (ici : une gradation de l'autoconsommation vers le marché).
\end{enumerate}
\end{methode}

\begin{center}
\begin{tabularx}{\linewidth}{|l|X|X|X|}
\hline
\rowcolor{popblueL} \textbf{Critères} & \textbf{Itinérante sur brûlis} & \textbf{Vivrière} & \textbf{Seconde génération} \\
\hline
Outillage & Houe, machette, hache, feu & Houe, machette, parfois charrue & Houe, machette, pulvérisateur, parfois motoculteur ou charrue \\
\hline
Techniques & Défriche-brûlis, jachère longue, rotation des parcelles & Jachère courte, cultures associées, fumure organique & Semences améliorées, engrais, pesticides, rotation raisonnée, irrigation \\
\hline
Cultures & Mil, sorgho, manioc, igname & Manioc, plantain, maïs, arachide, taro & Coton, riz, maïs intensif, maraîchage + vivrier \\
\hline
Destination & Autoconsommation quasi totale & Autoconsommation + petits surplus vendus & Vente d'une grande part + autoconsommation \\
\hline
Encadrement & Aucun & Faible & Fort (SODECOTON, SEMRY, MINADER, GIC) \\
\hline
\end{tabularx}
\end{center}

\textbf{Phrase de synthèse :} des savanes de l'Extrême-Nord aux Hautes Terres de l'Ouest, on passe ainsi d'un système tourné vers la survie de la famille à un système tourné vers le marché, l'agriculture de seconde génération occupant la position intermédiaire et dynamique de cette évolution.

\courssub{TP guidé : diagramme comparatif des rendements du maïs}

\begin{experience}[Construire un diagramme en barres groupées]
\textbf{Objectif :} représenter et comparer les rendements du maïs (en \SI{}{\tonne\per\hectare}) selon le système agricole.

\textbf{Matériel :} papier millimétré ou logiciel, règle, crayons de couleur, données du tableau ci-dessous (ordres de grandeur indicatifs).

\textbf{Données :}
\begin{center}
\begin{tabular}{|l|c|}
\hline
\rowcolor{popblueL} \textbf{Système agricole} & \textbf{Rendement du maïs (\SI{}{\tonne\per\hectare})} \\
\hline
Itinérant sur brûlis & 0,8 \\
\hline
Vivrière améliorée & 1,5 \\
\hline
Seconde génération & 3,0 \\
\hline
\end{tabular}
\end{center}

\textbf{Protocole :}
\begin{enumerate}
\item Trace deux axes : en abscisse, les trois systèmes ; en ordonnée, le rendement de 0 à 3,5 \SI{}{\tonne\per\hectare} (graduation tous les 0,5).
\item Pour chaque système, élève une barre de largeur égale et de hauteur correspondant au rendement ; laisse un espace identique entre les barres.
\item Colorie chaque barre d'une couleur différente et rédige la légende.
\item Inscris le titre, les noms des axes avec leurs unités, et la source des données.
\end{enumerate}

\textbf{Observations :} la barre de la seconde génération est près de quatre fois plus haute que celle du brûlis.

\textbf{Interprétation :} l'apport d'intrants et l'encadrement multiplient les rendements sans changer la nature familiale de l'exploitation ; c'est la preuve chiffrée des « grandes améliorations ».
\end{experience}

\begin{popfigure}
\begin{center}
\begin{tikzpicture}
\begin{axis}[
  ybar, bar width=22pt,
  width=0.85\linewidth, height=7cm,
  ymin=0, ymax=3.6,
  ylabel={Rendement du maïs (\SI{}{\tonne\per\hectare})},
  symbolic x coords={Brûlis, Vivrière, 2e génération},
  xtick=data,
  ytick={0,0.5,1,1.5,2,2.5,3,3.5},
  nodes near coords,
  every node near coord/.append style={font=\small\bfseries, color=popdark},
  axis lines=left,
  ymajorgrids=true, grid style={popline, dashed},
  x tick label style={font=\small},
  ylabel style={font=\small}
]
\addplot[fill=poporangeL, draw=poporange] coordinates {(Brûlis,0.8)};
\addplot[fill=popgreenL, draw=popgreen] coordinates {(Vivrière,1.5)};
\addplot[fill=popblueL, draw=popblue] coordinates {(2e génération,3.0)};
\end{axis}
\end{tikzpicture}
\end{center}
\legende{Rendements indicatifs du maïs selon le système agricole au Cameroun (ordres de grandeur, en tonnes par hectare) : 1. agriculture itinérante sur brûlis (environ 0,8 t/ha) ; 2. agriculture vivrière améliorée (environ 1,5 t/ha) ; 3. agriculture de seconde génération avec intrants (environ 3 t/ha). Source : ordres de grandeur FAO/MINADER.}
\end{popfigure}

\begin{exoresolu}[Commenter le diagramme des rendements]
\textbf{Énoncé.} À partir du diagramme ci-dessus : 1) calcule le taux d'accroissement du rendement du maïs entre le système sur brûlis et la seconde génération ; 2) propose deux explications à cet écart ; 3) indique une limite de la seconde génération.
\tcblower
\textbf{Solution détaillée.}

\begin{align*}
\text{1) Taux de variation } &= \frac{3{,}0 - 0{,}8}{0{,}8} \times 100 \\
&= \frac{2{,}2}{0{,}8} \times 100 \\
&= 275\,\%.
\end{align*}
Le rendement est multiplié par 3,75, soit une hausse d'environ 275 \%.

2) Explications : d'une part, l'utilisation de semences améliorées et d'engrais qui augmentent la fertilité et le potentiel des plantes ; d'autre part, l'encadrement technique (itinéraires culturaux, traitements contre les ravageurs) qui réduit les pertes.

3) Limite : ce système dépend de l'achat d'intrants coûteux et des prix de vente ; une chute des prix ou une mauvaise pluviométrie peut endetter le producteur (on peut aussi citer l'enclavement ou le vieillissement des planteurs).
\end{exoresolu}

\begin{exercice}[À toi de jouer]
1) Rédige une définition personnelle de l'agriculture de seconde génération en deux phrases maximum, en utilisant obligatoirement les mots : familiale, intrants, marché.

2) Recopie et complète le tableau comparatif des trois systèmes en ajoutant une ligne « Principales contraintes ».

3) Un jeune diplômé de Bafoussam hésite entre un emploi urbain précaire et l'installation comme maraîcher périurbain. Rédige un paragraphe argumenté (8 à 10 lignes) présentant deux avantages et deux risques du choix agricole.
\end{exercice}

\begin{corrige}[Exercice]
1) Exemple de définition attendue : « L'agriculture de seconde génération est une agriculture paysanne familiale qui utilise des intrants modernes (semences améliorées, engrais, pesticides) grâce à l'encadrement et au crédit. Orientée vers le marché, elle vend une part importante de sa production tout en conservant des cultures vivrières. »

2) Ligne « Principales contraintes » : itinérant sur brûlis — déforestation, épuisement des sols, faibles rendements ; vivrière — outillage sommaire, faible revenu monétaire, enclavement ; seconde génération — coût des intrants, endettement, fluctuation des prix, dépendance à l'encadrement.

3) Éléments de réponse : avantages — débouchés urbains proches et rémunérateurs (maraîchage de contre-saison dans les bas-fonds), revenu régulier possible sur petite surface, insertion sans diplôme supplémentaire coûteux ; risques — fluctuation des prix, coût des intrants, maladies des cultures, concurrence, accès au foncier. La copie doit mobiliser le vocabulaire de la leçon (chaîne de valeur, intrants, commercialisation) et conclure par un jugement nuancé.
\end{corrige}

\newpage

\coursec{Activité d'intégration}

\coursec{Leçon 04 : L'agriculture traditionnelle : de grandes améliorations}

\courssub{Objectifs de l'activité}

À l'issue de cette activité d'intégration, tu seras capable de :
\begin{itemize}
\item \cle{observer} et \cle{décrire} des documents géographiques (photographies de champs, tableau statistique) et de les \cle{classer} selon le type d'agriculture représenté ;
\item \cle{localiser} deux villages camerounais sur un fond de carte et \cle{légender} une carte avec les cultures dominantes ;
\item \cle{construire} un diagramme en barres comparant des rendements et des parts de production vendue ;
\item \cle{réaliser} le croquis d'un terroir villageois avec titre, légende et orientation ;
\item \cle{rédiger} un paragraphe argumenté mobilisant les notions de la leçon (outillage amélioré, techniques culturales, nouvelles cultures, agriculture de seconde génération).
\end{itemize}

\courssub{Situation}

\textbf{Deux villages, deux agricultures : enquête géographique au Cameroun.}

Dans le cadre d'une enquête de terrain sur les transformations de l'agriculture camerounaise, une équipe du lycée a visité deux villages :

\begin{itemize}
\item \textbf{Mbankomo}, village de la zone forestière du Centre (département de la Méfou-et-Akono), à une vingtaine de kilomètres de Yaoundé. On y pratique une agriculture itinérante sur brûlis, essentiellement vivrière : manioc, macabo, plantain, arachide, maïs, avec quelques cacaoyers en héritage des anciennes plantations. L'outillage reste simple (machette, houe, daba), la mécanisation est quasi absente, et l'essentiel de la récolte est consommé par la famille.
\item \textbf{Pitoa}, village de la savane du Nord (département du Bénoué), au nord de Garoua. On y pratique une culture de rente encadrée : le coton, sous le contrat de la \cle{SODECOTON} (Société de développement du coton du Cameroun), qui fournit semences sélectionnées, engrais et pesticides à crédit, encadre les producteurs par des techniciens agricoles et achète toute la production à un prix garanti. L'attelage bovin (culture attelée) y est répandu.
\end{itemize}

\textbf{Document 1. Photographies.}
Photo A : une parcelle forestière récemment défrichée, troncs calcinés encore visibles, cendres au sol, jeunes plants de manioc et de maïs semés en mélange ; une femme travaille à la houe. Photo B : un champ de coton aligné en rangs réguliers, un paysan guidant une charrue tirée par deux bœufs, un pulvérisateur à dos posé au bord du champ, et au loin un camion de collecte.

\textbf{Document 2. Tableau statistique comparatif (données indicatives de terrain, ordres de grandeur).}

\begin{center}
\begin{tabularx}{\linewidth}{|l|X|X|}
\hline
\rowcolor{popblueL} \textbf{Indicateur} & \textbf{Mbankomo (Centre)} & \textbf{Pitoa (Nord)} \\
\hline
Superficie moyenne par exploitation & 1 à 2 ha (parcelles dispersées) & 2 à 5 ha (champs groupés) \\
\hline
Culture dominante & Manioc, plantain (vivrier) & Coton (rente) + maïs, sorgho \\
\hline
Rendement de la culture principale & manioc : environ 10 t/ha & coton-graine : environ 1,2 à 1,5 t/ha \\
\hline
Outillage & machette, houe, daba & charrue à traction animale, pulvérisateur \\
\hline
Intrants utilisés & quasi aucun (cendres du brûlis) & semences sélectionnées, engrais NPK et urée, pesticides (à crédit SODECOTON) \\
\hline
Part de la production vendue & environ 20 \% (excédents au marché de Yaoundé) & environ 80 \% (coton vendu à la SODECOTON) \\
\hline
Main-d'œuvre & familiale & familiale + travail salarié saisonnier \\
\hline
\end{tabularx}
\end{center}

\textbf{Document 3.} Un fond de carte du Cameroun (contour simplifié, grandes villes : Yaoundé, Douala, Garoua, Bafoussam, Maroua) à compléter.

\courssub{Consignes}

Par groupes de quatre, vous traiterez les tâches suivantes, dans l'ordre :

\begin{enumerate}
\item \textbf{Observer et classer.} Décrivez en quelques lignes chacune des deux photographies (paysage, outils, organisation du champ), puis indiquez à quel type d'agriculture étudié dans la leçon correspond chacune d'elles. Justifiez votre classement par deux indices visibles sur la photo.
\item \textbf{Localiser et légender.} Sur le fond de carte du Cameroun, placez Mbankomo et Pitoa, tracez une flèche du nord, puis légendez la carte en indiquant la culture dominante de chaque zone (cacao-vivrier au Centre, coton au Nord) à l'aide de deux couleurs.
\item \textbf{Construire un diagramme.} À partir du document 2, construisez un diagramme en barres comparant, pour les deux villages, la part de la production vendue (en \%). Donnez un titre et une légende à votre diagramme.
\item \textbf{Construire un croquis.} Réalisez le croquis du terroir de Mbankomo en y faisant figurer : le village (habitat groupé), la parcelle cultivée de l'année, les jachères de différents âges, la forêt non défrichée, la piste vers Yaoundé. N'oubliez ni le titre, ni la flèche du nord, ni la légende numérotée.
\item \textbf{Rédiger.} En un paragraphe argumenté de 10 à 15 lignes, expliquez pourquoi l'agriculture de Pitoa est dite « de seconde génération » et dégagez au moins trois améliorations qui la distinguent de l'agriculture de Mbankomo (outillage, techniques culturales, rapport au marché).
\end{enumerate}

Restitution orale de 5 minutes par groupe, puis correction collective au tableau.

\courssub{Ressources à mobiliser}

\begin{aretenir}[Ce qu'il faut avoir en tête avant de commencer]
\begin{itemize}
\item \cle{Agriculture itinérante sur brûlis} : défrichement, brûlis (les cendres fertilisent le sol), culture 2 à 3 ans, puis abandon en jachère quand le sol s'épuise ; outillage simple, faibles rendements, production vivrière.
\item \cle{Agriculture vivrière} : l'essentiel de la production est destiné à l'autoconsommation ; seuls les excédents sont vendus.
\item \cle{Agriculture de seconde génération} : agriculture traditionnelle améliorée et encadrée (encadrement technique, intrants, crédit, débouché garanti), intégrée au marché, sans être encore une agriculture de plantation.
\item \cle{Améliorations de l'agriculture traditionnelle} : outillage amélioré (traction animale, petites machines), évolution des techniques culturales (semences sélectionnées, engrais, traitements, jachère améliorée), adoption de nouvelles cultures (coton, maïs amélioré, soja, etc.).
\item Savoir-faire cartographiques : titre, orientation, légende numérotée ou fléchée, échelle indicative ; un croquis n'est pas une carte : il sélectionne et organise l'information.
\end{itemize}
\end{aretenir}

\courssub{Démarche et corrigé commenté}

\begin{exoresolu}[Deux villages, deux agricultures : corrigé complet]
Énoncé : comparer les systèmes agricoles de Mbankomo et de Pitoa à partir des trois documents, construire les productions graphiques demandées et rédiger le paragraphe d'explication.
\tcblower
\textbf{Étape 1 — Description et classement des photographies.}

Photo A : paysage forestier défriché, troncs noircis et cendres (indices du brûlis), cultures en association (manioc + maïs), outil manuel (houe), aucune trace de mécanisation. Il s'agit de l'\cle{agriculture itinérante sur brûlis}, de type vivrier.

Photo B : rangs réguliers de coton, charrue à traction bovine (outillage amélioré), pulvérisateur (traitements phytosanitaires), camion de collecte (intégration au marché). Il s'agit d'une \cle{agriculture de seconde génération}, encadrée par la SODECOTON.

\textbf{Étape 2 — Localisation et légendage de la carte.}

Mbankomo se place au sud du Cameroun, près de Yaoundé (zone forestière du Centre) ; Pitoa se place au nord de Garoua (zone soudanienne du Nord). On colore la zone Centre-Sud en vert (cacao, cultures vivrières) et la zone Nord en jaune (coton), avec flèche du nord et titre.

\begin{popfigure}
\begin{center}
\includegraphics[width=\linewidth,height=10.5cm,keepaspectratio]{N04/cmr_map.jpg}
\end{center}
\legende{Carte du Cameroun (relief ombré, régions, villes, routes et voies ferrées ; légende en anglais). Elle permet de localiser les deux terroirs étudiés : au Centre, la zone forestière du Centre-Sud autour de Yaoundé (cacao, cultures vivrières, village de Mbankomo, près de la capitale) et, au Nord, la région de Garoua (coton de la SODECOTON, Pitoa, près de Garoua). \\ Source : Wikimedia Commons, CIA, domaine public.}
\end{popfigure}

\textbf{Étape 3 — Diagramme en barres : part de la production vendue.}

\begin{popfigure}
\begin{tikzpicture}
\begin{axis}[
ybar, bar width=1.1cm,
width=0.75\linewidth, height=6cm,
ymin=0, ymax=100,
ylabel={Part de la production vendue (en \%)},
symbolic x coords={Mbankomo, Pitoa},
xtick=data,
nodes near coords,
nodes near coords align={vertical},
every node near coord/.append style={font=\small},
ymajorgrids=true, grid style={popline!40},
]
\addplot[fill=popgreen, draw=popdark] coordinates {(Mbankomo,20)};
\addplot[fill=popgold, draw=popdark] coordinates {(Pitoa,80)};
\end{axis}
\end{tikzpicture}
\legende{Part de la production vendue à Mbankomo (environ 20 \%) et à Pitoa (environ 80 \%) — données indicatives du document 2.}
\end{popfigure}

Le contraste est net : à Mbankomo, l'essentiel de la récolte est autoconsommé (agriculture vivrière) ; à Pitoa, l'essentiel est vendu (agriculture de rente intégrée au marché).

\textbf{Étape 4 — Croquis du terroir de Mbankomo.}

\begin{popfigure}
\begin{tikzpicture}[scale=1]
% Forêt
\fill[popgreenL] (0,0) rectangle (10,6.4);
\foreach \x/\y in {0.7/5.6, 1.6/5.9, 2.6/5.5, 8.6/5.7, 9.4/5.2, 0.9/0.7, 8.8/0.8, 9.5/1.6}
\node[font=\scriptsize, popgreen!60!black] at (\x,\y) {$\clubsuit$};
% Jachères
\fill[poppinkL] (0.6,2.2) rectangle (3.0,4.6);
\fill[poporangeL] (3.4,2.4) rectangle (5.4,4.4);
% Parcelle cultivée
\fill[popgoldL] (5.9,2.6) rectangle (7.9,4.6);
\foreach \x in {6.2,6.6,7.0,7.4,7.8}
\draw[popdark!60] (\x,2.7) -- (\x,4.5);
% Village
\foreach \x/\y in {4.0/1.2, 4.5/1.4, 5.0/1.1, 4.3/0.8, 4.8/0.7}
\draw[popdark, fill=popcream] (\x,\y) rectangle ++(0.35,0.28);
% Piste
\draw[thick, dashed, popink] (4.5,0.6) -- (4.5,-0.2) -- (4.5,-0.6) node[below, font=\scriptsize]{vers Yaoundé};
\draw[thick, dashed, popink] (4.5,1.6) -- (4.6,2.4);
% Nord
\draw[-{Stealth}, thick, popdark] (9.4,3.4) -- (9.4,4.4) node[above, font=\scriptsize]{N};
% Légende fléchée
\node[font=\scriptsize, anchor=west] at (0.6,4.9) {3. Jachères (ancienne et récente)};
\node[font=\scriptsize, anchor=west] at (5.9,4.9) {2. Parcelle cultivée de l'année};
\node[font=\scriptsize, anchor=west] at (3.4,1.9) {1. Village (habitat groupé)};
\node[font=\scriptsize, anchor=west] at (6.6,0.9) {4. Forêt non défrichée};
\end{tikzpicture}
\legende{Croquis du terroir de Mbankomo : rotation des parcelles sur brûlis. 1. Village ; 2. Parcelle cultivée ; 3. Jachères ; 4. Forêt. Échelle indicative : le terroir s'étend sur quelques kilomètres.}
\end{popfigure}

\textbf{Étape 5 — Paragraphe argumenté (corrigé type).}

L'agriculture pratiquée à Pitoa est dite « de seconde génération » parce qu'elle reste une agriculture paysanne familiale, mais profondément améliorée et encadrée, à mi-chemin entre l'agriculture traditionnelle et l'agriculture de plantation. Trois améliorations la distinguent de celle de Mbankomo. D'abord, l'\cle{outillage est amélioré} : la charrue à traction bovine remplace la houe, ce qui augmente les superficies cultivées (2 à 5 ha contre 1 à 2 ha). Ensuite, les \cle{techniques culturales ont évolué} : semences sélectionnées, engrais (NPK, urée) et pesticides fournis à crédit par la SODECOTON, semis en lignes, alors qu'à Mbankomo la fertilité repose sur les seules cendres du brûlis et la jachère. Enfin, le \cle{rapport au marché} a changé : environ 80 \% de la production de Pitoa est vendue grâce au contrat d'achat garanti par la SODECOTON, contre 20 \% à Mbankomo où l'autoconsommation domine. L'adoption d'une nouvelle culture de rente, le coton, a ainsi transformé une agriculture de subsistance en agriculture commerciale encadrée.
\end{exoresolu}

\begin{attention}[Erreurs fréquentes à éviter]
\begin{itemize}
\item Confondre agriculture de seconde génération et agriculture de plantation : à Pitoa, l'exploitation reste familiale et paysanne ; ce n'est pas une grande concession capitaliste comme celles de la CDC ou de la PHP.
\item Oublier le titre, la flèche du nord ou la légende sur le croquis : chaque oubli coûte des points au Baccalauréat.
\item Décrire les photos sans les classer, ou classer sans justifier par des indices visibles.
\item Affirmer que le brûlis « détruit » simplement le sol : il le fertilise à court terme (cendres), c'est le raccourcissement des jachères qui pose problème.
\item Comparer des rendements de cultures différentes (tonnes de manioc contre tonnes de coton) sans précaution : préfère comparer les parts vendues et les pratiques.
\end{itemize}
\end{attention}

\courssub{Bilan de l'activité}

Cette enquête comparée a permis de mettre en évidence trois acquis majeurs de la leçon. D'une part, l'agriculture traditionnelle camerounaise n'est pas figée : entre Mbankomo et Pitoa, on mesure concrètement le chemin parcouru grâce à l'\cle{outillage amélioré}, à l'\cle{évolution des techniques culturales} et à l'\cle{adoption de nouvelles cultures} de rente comme le coton. D'autre part, l'\cle{agriculture de seconde génération} apparaît comme une voie camerounaise de modernisation paysanne : encadrement (SODECOTON), intrants à crédit, débouché garanti, tout en conservant la base familiale de l'exploitation. Enfin, tu as mobilisé l'ensemble des savoir-faire du programme — observer, localiser, légender, construire un diagramme et un croquis, rédiger une argumentation — qui sont précisément ceux évalués dans les épreuves de géographie du Baccalauréat. Garde cette démarche en tête : un bon géographe part toujours du document pour construire sa démonstration.

\newpage

\coursec{Méthodes et exercices résolus}

\coursec{L'agriculture traditionnelle : de grandes améliorations}

\courssub{Observer, localiser et décrire un espace agricole traditionnel}

\begin{methode}[Décrire un paysage agricole ou une photographie de champ]
Face à une photographie, une carte ou un croquis d'espace agricole camerounais (mosaïque de champs de l'Ouest, brûlis en zone forestière, plantation villageoise), la description suit une démarche rigoureuse en cinq étapes :
\begin{enumerate}
\item \textbf{Localiser} : situer l'espace (région, département, zone agro-écologique : hautes terres de l'Ouest, savane du Nord, forêt du Sud) et préciser l'échelle d'observation.
\item \textbf{Identifier le type d'agriculture} : agriculture itinérante sur brûlis, agriculture vivrière sédentaire, arboriculture villageoise (cacaoyers, caféiers), jachère... Justifier le choix par des indices visibles (cendres, champs dispersés, association de cultures, présence d'outils).
\item \textbf{Décrire les faits visibles} : forme et taille des parcelles, cultures présentes, outillage (machette, houe, motoculteur, tracteur), aménagements (irrigation, billons, terrasses, haies), habitat, voies d'accès. Ne décrire que ce qui se voit.
\item \textbf{Interpréter} : relier chaque fait observé à une caractéristique de l'agriculture traditionnelle (faible mécanisation, polyculture, main-d'œuvre familiale) ou à ses améliorations (outillage motorisé, nouvelles cultures, encadrement).
\item \textbf{Conclure} : dégager en une ou deux phrases le problème géographique posé (niveau de modernisation, contraintes naturelles, intégration au marché).
\end{enumerate}
\textbf{Réflexes} : toujours aller du général au particulier ; employer le vocabulaire du cours (\cle{brûlis}, \cle{jachère}, \cle{polyculture vivrière}, \cle{agriculture de seconde génération}) ; interdire toute interprétation non fondée sur un fait observé.
\end{methode}

\begin{exoresolu}[Description d'une photographie de champ en zone forestière du Sud]
\textbf{Énoncé (type Bac)} : Une photographie prise dans la région du Centre, à une trentaine de kilomètres de Mbalmayo, montre une clairière récente en forêt : troncs d'arbres abattus et partiellement carbonisés, souches encore visibles, sol nu entre lesquels une femme plante des bâtons de manioc à l'aide d'une machette ; en bordure, quelques bananiers et un cacaoyer ; au fond, une case en matériaux locaux. Décrivez cette photographie et identifiez le système agricole représenté, en justifiant votre réponse.
\tcblower
\textbf{Solution détaillée.}

\emph{Étape 1 — Localisation.} La photographie est prise dans la région du Centre (environs de Mbalmayo), c'est-à-dire dans la zone agro-écologique forestière du Sud-Cameroun, caractérisée par une pluviométrie abondante (environ \SI{1500}{\milli\meter} à \SI{2000}{\milli\meter} par an) répartie sur quatre saisons.

\emph{Étape 2 — Description des faits visibles.} On observe : (a) une clairière récente ouverte dans la forêt dense ; (b) des troncs abattus et carbonisés, signe d'un feu récent ; (c) un sol nu, sans labour ni billons ; (d) une paysanne plantant des boutures de manioc avec une machette, seul outil visible ; (e) une association de cultures (manioc, bananiers, cacaoyer) ; (f) un habitat villageois proche.

\emph{Étape 3 — Interprétation.} Les troncs carbonisés et la clairière récente traduisent la technique du \cle{défriche-brûlis} : la végétation est coupée puis brûlée, et la cendre fertilise temporairement le sol. L'absence de labour et d'outil autre que la machette indique un outillage rudimentaire et une absence de mécanisation. L'association manioc-bananiers-cacaoyer est typique de la \cle{polyculture} vivrière et commerciale des champs forestiers camerounais.

\emph{Étape 4 — Identification et conclusion.} Il s'agit de l'\cle{agriculture itinérante sur brûlis}, forme traditionnelle de l'agriculture vivrière pratiquée en zone forestière : le champ est cultivé deux à trois ans puis abandonné en jachère forestière (de 10 à 15 ans) le temps que la fertilité se reconstitue. Ce système, adapté à une faible densité de population, devient fragile lorsque la pression démographique raccourcit les jachères, ce qui menace la fertilité des sols et la forêt.
\end{exoresolu}

\courssub{Légender une carte et construire un croquis agricole}

\begin{methode}[Construire un croquis des systèmes agricoles traditionnels du Cameroun]
Le croquis est un schéma simplifié, sélectif et hiérarchisé. Démarche en six étapes :
\begin{enumerate}
\item \textbf{Tracer le fond de carte} : contour simplifié du Cameroun (polygone), flèche du nord, échelle indicative, titre en haut.
\item \textbf{Choisir 2 à 4 informations maximum} (ex. : agriculture itinérante sur brûlis au Sud ; cultures vivrières sédentaires des hautes terres de l'Ouest ; élevage et cultures de saison sèche au Nord ; zones d'agriculture de seconde génération).
\item \textbf{Associer à chaque information un code graphique unique} : aplat de couleur pour les grandes zones, point pour une ville, trait pour un axe, flèche pour un flux. Un même code = une même réalité.
\item \textbf{Placer les faits} avec sobriété : pas de surcharge, noms écrits horizontalement.
\item \textbf{Rédiger la légende} en bas ou à droite, numérotée ou fléchée, dans le même ordre que le titre.
\item \textbf{Vérifier la cohérence} : titre, orientation, légende complète, aucun élément non légendé.
\end{enumerate}
\textbf{Réflexes} : le croquis n'est pas une carte reproduite ; il illustre un raisonnement. Toute couleur non expliquée en légende fait perdre des points.
\end{methode}

\begin{exoresolu}[Croquis des grands systèmes d'agriculture traditionnelle au Cameroun]
\textbf{Énoncé (type Bac)} : À l'aide de vos connaissances, construisez un croquis intitulé « Les grands systèmes d'agriculture traditionnelle au Cameroun ». Vous y ferez figurer : les zones d'agriculture itinérante sur brûlis, les zones de cultures vivrières sédentaires intensives, les zones d'agropastoralisme, ainsi que deux pôles d'agriculture de seconde génération.
\tcblower
\textbf{Solution détaillée.}

\emph{Raisonnement préalable.} La répartition des systèmes agricoles suit les zones agro-écologiques : (1) le Sud forestier humide (régions du Centre, du Sud, de l'Est, du Littoral et du Sud-Ouest) est le domaine de l'agriculture itinérante sur brûlis ; (2) les hautes terres de l'Ouest et du Nord-Ouest, très densément peuplées, pratiquent une agriculture vivrière sédentaire intensive (maïs, haricot, pomme de terre, maraîchage) ; (3) le Grand Nord soudano-sahélien (Extrême-Nord, Nord, Adamaoua) combine cultures de saison des pluies (mil, sorgho, niébé, coton) et élevage : c'est l'agropastoralisme ; (4) l'agriculture de seconde génération (périmètres modernisés, encadrement, motorisation) se repère autour de pôles comme la plaine des Mbo (SOSUCAM à Nkoteng et Bandjock pour la canne à sucre) ou les périmètres rizicoles de la SEMRY à Yagoua.

\emph{Construction du croquis.} On trace le contour simplifié du pays, on place la flèche du nord, puis on colore trois grandes zones et on marque deux pôles par des points :

\begin{popfigure}
\begin{center}
\includegraphics[width=\linewidth,height=10.5cm,keepaspectratio]{N04/cmr_map.jpg}
\end{center}
\legende{Fond de carte pour le croquis des grands systèmes d'agriculture traditionnelle au Cameroun (relief, régions, villes, routes ; légende en anglais). Le croquis attendu y situe : 1 : l'agriculture itinérante sur brûlis (zone forestière du Sud et de l'Est) ; 2 : les cultures vivrières sédentaires intensives (hautes terres de l'Ouest et du Nord-Ouest) ; 3 : l'agropastoralisme (savanes du Nord et de l'Adamaoua) ; 4 : les pôles d'agriculture de seconde génération (canne à sucre de la SOSUCAM, riz de la SEMRY) par des symboles ponctuels. \\ Source : Wikimedia Commons, CIA, domaine public.}
\end{popfigure}

\emph{Vérification finale.} Le croquis comporte un titre, une orientation, quatre codes graphiques tous expliqués dans la légende, et des villes-repères pour la localisation. \textbf{Conclusion} : la carte montre que les systèmes agricoles traditionnels camerounais s'organisent selon un gradient climatique sud-nord, et que l'agriculture de seconde génération apparaît en pôles ponctuels encadrés par de grandes sociétés.
\end{exoresolu}

\courssub{Inventorier, classer et interpréter des documents statistiques}

\begin{methode}[Exploiter un tableau statistique agricole]
\begin{enumerate}
\item \textbf{Lire le titre complet} : quoi (production, rendement, superficie), où (région, pays), quand (années), unité (tonnes, kg/ha, \%).
\item \textbf{Repérer les grandeurs extrêmes} : valeur maximale, valeur minimale, évolutions sensibles.
\item \textbf{Calculer les indicateurs utiles} :
\begin{itemize}
\item Taux de variation entre deux dates : $T = \dfrac{V_{2}-V_{1}}{V_{1}}\times 100$ (en \%).
\item Part d'un élément dans un total : $p = \dfrac{\text{partie}}{\text{total}}\times 100$ (en \%).
\item Rendement : $R = \dfrac{\text{production}}{\text{superficie}}$ (en tonnes/ha ou kg/ha).
\end{itemize}
\item \textbf{Classer} : hiérarchiser les régions ou les cultures (du plus fort au plus faible).
\item \textbf{Interpréter géographiquement} : relier les chiffres aux faits du cours (outillage amélioré, adoption de nouvelles cultures, encadrement, contraintes climatiques).
\item \textbf{Conclure} par une phrase de synthèse qui répond à la consigne.
\end{enumerate}
\textbf{Attention} : toujours écrire l'unité dans le résultat et arrondir de façon cohérente (une décimale suffit en général).
\end{methode}

\begin{exoresolu}[Calcul et interprétation : production de maïs dans l'Ouest]
\textbf{Énoncé (type Bac)} : Le tableau suivant donne la production de maïs (en milliers de tonnes) dans deux régions du Cameroun (ordres de grandeur indicatifs) :
\begin{center}
\begin{tabularx}{\linewidth}{|l|X|X|}\hline
\rowcolor{popblueL} \textbf{Région} & \textbf{Année 1} & \textbf{Année 10} \\ \hline
Ouest & 400 & 620 \\ \hline
Extrême-Nord & 150 & 180 \\ \hline
\end{tabularx}
\end{center}
1) Calculez le taux de variation de la production dans chaque région. 2) Calculez la part de l'Ouest dans la production totale des deux régions à l'année 10. 3) Tirez-en deux enseignements sur l'évolution de l'agriculture vivrière camerounaise.
\tcblower
\textbf{Solution détaillée.}

\emph{1) Taux de variation.} Formule : $T=\dfrac{V_{2}-V_{1}}{V_{1}}\times 100$.
\begin{align*}
T_{\text{Ouest}} &= \frac{620-400}{400}\times 100 = \frac{220}{400}\times 100 = 55\%\\
T_{\text{Extrême-Nord}} &= \frac{180-150}{150}\times 100 = \frac{30}{150}\times 100 = 20\%
\end{align*}
La production de maïs a augmenté de \SI{55}{\%} dans l'Ouest et de \SI{20}{\%} dans l'Extrême-Nord entre l'année 1 et l'année 10.

\emph{2) Part de l'Ouest à l'année 10.} Total année 10 : $620+180=800$ milliers de tonnes.
\[p = \frac{620}{800}\times 100 = 77,5\%\]
L'Ouest fournit \SI{77,5}{\%} de la production des deux régions à l'année 10.

\emph{3) Interprétation.} Premier enseignement : la production vivrière progresse partout, ce qui traduit les améliorations de l'agriculture traditionnelle (outillage amélioré, semences sélectionnées, extension des superficies). Deuxième enseignement : la progression est beaucoup plus forte dans l'Ouest, région de hautes terres densément peuplée où l'agriculture vivrière sédentaire intensive (maïs associé au haricot et au maraîchage) bénéficie de sols volcaniques fertiles, d'une pluviométrie favorable et d'un meilleur encadrement ; l'Extrême-Nord, soumis à l'irrégularité des pluies, progresse plus lentement.

\textbf{Conclusion} : les calculs confirment un dynamisme inégal de l'agriculture vivrière camerounaise, très fort dans les hautes terres de l'Ouest, plus modéré dans les savanes du Nord.
\end{exoresolu}

\courssub{Construire et lire un diagramme statistique agricole}

\begin{methode}[Réaliser un diagramme en bâtons ou circulaire]
\begin{enumerate}
\item \textbf{Choisir le bon diagramme} : bâtons pour comparer des quantités (productions par culture) ; circulaire (camembert) pour montrer des parts d'un total (répartition des superficies).
\item \textbf{Pour un diagramme circulaire}, convertir chaque part en angle : $\alpha = \dfrac{\text{effectif}}{\text{total}}\times 360$ (en degrés). Vérifier que la somme des angles vaut $360^{\circ}$.
\item \textbf{Tracer avec soin} : axes gradués et légendés (diagramme en bâtons) ou secteurs proportionnels (diagramme circulaire), couleurs distinctes.
\item \textbf{Ajouter titre, unités, source} et légende des couleurs.
\item \textbf{Lire le diagramme} : repérer le secteur ou le bâton dominant, les écarts, puis interpréter.
\end{enumerate}
\end{methode}

\begin{exoresolu}[Diagramme circulaire : répartition des superficies vivrières]
\textbf{Énoncé (type Bac)} : Dans un village de la région du Centre, les superficies cultivées se répartissent ainsi (en hectares) : manioc 40, plantain 30, maïs 20, arachide 10. 1) Calculez l'angle de chaque secteur d'un diagramme circulaire représentant cette répartition. 2) Tracez le diagramme. 3) Quel enseignement géographique en tirez-vous ?
\tcblower
\textbf{Solution détaillée.}

\emph{1) Calcul des angles.} Total : $40+30+20+10=100$ ha. Formule : $\alpha=\dfrac{\text{effectif}}{\text{total}}\times 360$.
\begin{align*}
\alpha_{\text{manioc}} &= \frac{40}{100}\times 360 = 144^{\circ} &
\alpha_{\text{plantain}} &= \frac{30}{100}\times 360 = 108^{\circ}\\
\alpha_{\text{maïs}} &= \frac{20}{100}\times 360 = 72^{\circ} &
\alpha_{\text{arachide}} &= \frac{10}{100}\times 360 = 36^{\circ}
\end{align*}
Vérification : $144+108+72+36 = 360^{\circ}$. Le calcul est juste.

\emph{2) Tracé du diagramme.}

\begin{popfigure}
\begin{center}
\begin{tikzpicture}[scale=1.0]
% Secteurs : manioc 144°, plantain 108°, maïs 72°, arachide 36°
\fill[popgreen!60] (0,0) -- (90:2.2) arc (90:234:2.2) -- cycle;
\fill[poporange!60] (0,0) -- (234:2.2) arc (234:342:2.2) -- cycle;
\fill[popgold!70] (0,0) -- (342:2.2) arc (342:414:2.2) -- cycle;
\fill[poppink!60] (0,0) -- (54:2.2) arc (54:90:2.2) -- cycle;
\draw[popdark] (0,0) circle (2.2);
% Étiquettes
\node at (162:1.5) {\small Manioc 40\,\%};
\node at (288:1.5) {\small Plantain 30\,\%};
\node at (18:1.5) {\small Maïs 20\,\%};
\node at (72:2.75) {\small Arachide 10\,\%};
\end{tikzpicture}
\end{center}
\legende{Diagramme circulaire : répartition des superficies cultivées dans un village de la région du Centre (total : 100 ha ; données fictives d'application).}
\end{popfigure}

\emph{3) Interprétation.} Le manioc et le plantain occupent à eux deux \SI{70}{\%} des superficies ($40+30=70$ ha sur 100). \textbf{Conclusion} : le diagramme illustre la dominance des cultures vivrières de base de la zone forestière (manioc, plantain), piliers de la sécurité alimentaire, tandis que le maïs et l'arachide complètent l'assolement : c'est le profil typique de la polyculture vivrière camerounaise.
\end{exoresolu}

\courssub{Commenter un document sur les améliorations de l'agriculture traditionnelle}

\begin{methode}[Rédiger un commentaire de document (texte, photo, tableau)]
\begin{enumerate}
\item \textbf{Introduction} : présenter le document (nature, source, date, thème) et annoncer le problème posé (ex. : comment l'agriculture traditionnelle camerounaise se modernise-t-elle ?).
\item \textbf{Analyse} en deux ou trois parties ordonnées : décrire les faits, les expliquer par les notions du cours (\cle{outillage amélioré}, \cle{techniques culturales}, \cle{nouvelles cultures}, \cle{agriculture de seconde génération}), citer des exemples camerounais précis (SEMRY, SOSUCAM, SODECAO, IRAD).
\item \textbf{Discussion} : montrer les limites (coût des intrants, endettement, faible mécanisation réelle, disparités régionales).
\item \textbf{Conclusion} : répondre au problème et ouvrir (vers l'agriculture de plantation ou la sécurité alimentaire).
\end{enumerate}
\textbf{Réflexes} : un commentaire n'est ni une paraphrase ni un récital du cours ; chaque affirmation doit s'appuyer sur le document ou sur un exemple daté et localisé.
\end{methode}

\begin{exoresolu}[Commentaire : l'agriculture de seconde génération dans la plaine des Mbo]
\textbf{Énoncé (type Bac)} : Document : « Dans la plaine des Mbo (région du Centre), la SOSUCAM cultive environ \SI{20000}{ha} de canne à sucre sur les sites de Nkoteng et Bandjock : périmètres drainés et irrigués, mécanisation complète des travaux, usines de transformation sur place, emploi de milliers de salariés et de nombreux planteurs villageois encadrés. » À partir de ce document et de vos connaissances, montrez en quoi la plaine des Mbo illustre l'agriculture de seconde génération au Cameroun et discutez ses limites.
\tcblower
\textbf{Solution détaillée (plan et développement rédigé).}

\emph{Introduction.} Le document est un extrait de texte descriptif portant sur l'exploitation sucrière de la SOSUCAM dans la plaine des Mbo (Centre). Il pose le problème suivant : en quoi ce type d'exploitation traduit-il la modernisation de l'agriculture camerounaise, et quelles en sont les limites ?

\emph{I. Une agriculture de seconde génération accomplie.} Le document montre d'abord un \cle{outillage amélioré} poussé à son maximum : la mécanisation complète (tracteurs, moissonneuses) remplace la houe et la machette de l'agriculture traditionnelle. Ensuite, les \cle{techniques culturales} sont profondément renouvelées : drainage, irrigation, sélection variétale, usage d'engrais et de produits phytosanitaires, ce qui élève fortement les rendements par rapport aux champs vivriers. Enfin, l'intégration production-transformation (usines sur place) et l'encadrement des planteurs villageois (formation, intrants, débouchés garantis) caractérisent précisément l'\cle{agriculture de seconde génération} : une agriculture traditionnelle améliorée, insérée dans les circuits commerciaux et appuyée par une structure moderne.

\emph{II. Des limites réelles.} La mécanisation intégrale exige des capitaux considérables, inaccessibles au paysan isolé ; le modèle reste donc concentré entre les mains de grandes sociétés. La monoculture de la canne fragilise les sols et dépend des cours mondiaux du sucre ; les conflits fonciers avec les riverains et la précarité de certains emplois saisonniers sont également signalés. Par ailleurs, ce modèle en pôles contraste avec la lente modernisation des campagnes environnantes.

\emph{Conclusion.} La plaine des Mbo illustre la réussite de l'agriculture de seconde génération au Cameroun — outillage motorisé, techniques modernes, intégration au marché — mais révèle aussi ses limites : coûts élevés, monoculture, inégalités d'accès à la modernisation. Elle invite à réfléchir à une diffusion plus large des améliorations vers l'ensemble du monde vivrier.
\end{exoresolu}

\courssub{Distinguer et comparer les systèmes agricoles (étude de cas)}

\begin{methode}[Construire un tableau comparatif et rédiger une comparaison]
\begin{enumerate}
\item \textbf{Définir des critères de comparaison} identiques pour tous les systèmes : milieu naturel, outillage, techniques, cultures, main-d'œuvre, destination de la production, rendements, limites.
\item \textbf{Remplir le tableau} ligne par ligne (un critère = une ligne), jamais l'inverse.
\item \textbf{Dégager les points communs} puis \textbf{les différences} : une comparaison ne juxtapose pas deux descriptions, elle confronte.
\item \textbf{Expliquer les contrastes} par les facteurs géographiques (climat, densité de population, encadrement, marchés).
\item \textbf{Conclure} en hiérarchisant : quel système est le plus modernisé, le plus durable, le plus menacé ?
\end{enumerate}
\end{methode}

\begin{exoresolu}[Comparer brûlis forestier et agriculture de seconde génération]
\textbf{Énoncé (type Bac)} : À l'aide d'un tableau puis d'un paragraphe rédigé, comparez l'agriculture itinérante sur brûlis de la région de l'Est et l'agriculture de seconde génération des périmètres rizicoles de la SEMRY (Yagoua, Extrême-Nord).
\tcblower
\textbf{Solution détaillée.}

\emph{Tableau comparatif.}

\begin{center}
\begin{tabularx}{\linewidth}{|l|X|X|}\hline
\rowcolor{popblueL} \textbf{Critère} & \textbf{Brûlis (Est)} & \textbf{Seconde génération (SEMRY)} \\ \hline
Milieu & Forêt équatoriale humide & Savane soudanienne, plaine inondable du Logone \\ \hline
Outillage & Machette, houe (rudimentaire) & Motoculteurs, tracteurs, matériel de récolte \\ \hline
Techniques & Défriche-brûlis, jachère longue, aucun intrant & Irrigation, engrais, semences sélectionnées, billonnage \\ \hline
Cultures & Manioc, plantain, macabo (polyculture) & Riz (monoculture améliorée) \\ \hline
Main-d'œuvre & Familiale & Familiale encadrée + techniciens \\ \hline
Destination & Autoconsommation, surplus vendus & Vente organisée (marché national, transformation) \\ \hline
Rendements & Faibles et décroissants & Élevés et stabilisés \\ \hline
Limites & Déforestation, fertilité éphémère & Coût des intrants, endettement, dépendance à l'encadrement \\ \hline
\end{tabularx}
\end{center}

\emph{Paragraphe rédigé.} Les deux systèmes relèvent de l'agriculture pratiquée par des paysans camerounais sur de petites exploitations familiales, mais ils s'opposent sur presque tous les critères. Dans l'Est, l'agriculture itinérante sur brûlis repose sur la fertilité éphémère des cendres, un outillage rudimentaire et une polyculture vivrière tournée vers l'autoconsommation ; elle suppose de longues jachères et devient destructrice de forêt lorsque la population croît. À Yagoua, la SEMRY a transformé l'agriculture traditionnelle en agriculture de seconde génération : l'eau est maîtrisée par l'irrigation, les semences sont sélectionnées, la motorisation allège le travail et la production de riz est intégralement orientée vers le marché. Le contraste s'explique par l'encadrement (présent dans le Nord, absent dans l'Est) et par les milieux : la forêt humide impose la jachère, tandis que la plaine du Logone se prête aux périmètres irrigués.

\textbf{Conclusion} : la comparaison montre que l'amélioration de l'agriculture traditionnelle camerounaise passe moins par le changement de cultures que par l'outillage, la maîtrise de l'eau et l'encadrement — conditions qui restent inégalement réparties sur le territoire national.
\end{exoresolu}

\begin{aretenir}[Les trois formes de l'agriculture traditionnelle camerounaise]
\begin{itemize}
\item \textbf{Agriculture itinérante sur brûlis (ou défriche-brûlis) :} pratiquée en zone forestière humide (Sud, Est, Centre). Défrichement au feu, fertilité cendrée éphémère, polyculture vivrière (manioc, plantain, macabo), outillage manuel (machette, houe), jachère longue (10--15 ans), main-d'œuvre familiale, autoconsommation dominante. Menacée par le raccourcissement des jachères (pression démographique).
\item \textbf{Agriculture vivrière sédentaire intensive :} pratiquée dans les hautes terres de l'Ouest et du Nord-Ouest (Bamileke, Bamoun). Forte densité de population, sols volcaniques fertiles, pas de jachère (fumure organique, terrasses, billons), polyculture intensive (maïs, haricot, pomme de terre, maraîchage), outillage amélioré (houe, motoculteur), main-d'œuvre familiale dense, production excédentaire vendue sur les marchés urbains.
\item \textbf{Agriculture de seconde génération :} agriculture paysanne modernisée par l'encadrement (SEMRY, SOSUCAM, SODECAO, UNVDA, MIDENO). Maintien de la structure familiale mais adoption d'outillage motorisé, d'intrants chimiques (engrais, pesticides), de semences sélectionnées (IRAD), de techniques culturales modernes (irrigation, drainage, billonnage) et d'une orientation marchande affirmée. Exemples : riz de la SEMRY (Yagoua), canne à sucre SOSUCAM (Nkoteng, Bandjock), cacao SODECAO (Sud, Centre, Sud-Ouest).
\end{itemize}
\end{aretenir}

\begin{propriete}[Caractères de l'agriculture de seconde génération au Cameroun]
L'agriculture de seconde génération se définit par la combinaison de cinq piliers :
\begin{enumerate}
\item \textbf{Motorisation partielle ou totale :} remplacement de la traction humaine par la traction animale (Nord) ou motorisée (motoculteur, tracteur).
\item \textbf{Maîtrise de l'eau :} irrigation (SEMRY, UNVDA, MIDENO) ou drainage (SOSUCAM), billonnage, calendriers culturaux maîtrisés.
\item \textbf{Intrants et semences améliorées :} engrais minéraux, produits phytosanitaires, variétés à haut rendement (ex. riz NERICA, maïs CMS, canne R570) issues de la recherche (IRAD).
\item \textbf{Encadrement technique et commercial :} vulgarisation, crédit agricole, groupements de producteurs, débouchés garantis (usines, offices de commercialisation).
\item \textbf{Intégration aux marchés :} production monétarisée, transformation locale (usines de décorticage, sucreries), contribution à la sécurité alimentaire nationale et aux exportations.
\end{enumerate}
\textit{N.B. : La démonstration de ces caractères n'est pas exigible en tant que telle au Bac, mais ils constituent la grille de lecture obligatoire pour tout commentaire de document ou étude de cas.}
\end{propriete}

\begin{savaistu}[La recherche agronomique et les projets d'appui : l'IRAD, le PDDA, le PNDP]
L'\textbf{IRAD} (Institut de Recherche Agricole pour le Développement) est le pilier scientifique : il produit les variétés améliorées (maïs CMS 8704, 9022 ; riz NERICA ; manioc TMS ; plantain PITA) et les itinéraires techniques adaptés aux cinq zones agro-écologiques. Depuis les années 2000, l'État pilote la modernisation via le \textbf{PDDA} (Programme de Développement de l'Agriculture) et le \textbf{PNDP} (Programme National de Développement Participatif) : financement de pistes rurales, de magasins de stockage, d'unités de transformation villageoises, d'appui aux organisations paysannes (OP). Ces programmes visent à diffuser les améliorations de la « seconde génération » hors des grands pôles historiques (SEMRY, SOSUCAM) vers l'ensemble du tissu paysan, pour réduire la dépendance aux importations (riz, blé, poisson, huile) et renforcer l'intégration nationale par les flux commerciaux internes.
\end{savaistu}

\begin{attention}[Les 5 erreurs qui coûtent des points au Bac]
\begin{enumerate}
\item \textbf{Confondre agriculture traditionnelle et agriculture de plantation.} La plantation (bananeraies de la CDC/Pamol, hévéa, palmier à huile industriel) est une agriculture commerciale capitaliste, salariale, à grande échelle ; l'agriculture traditionnelle reste paysanne et familiale, même améliorée. Ne mélangez jamais les deux dans une même définition.
\item \textbf{Décrire sans localiser ni conclure.} Une description de photo ou de carte sans région, sans zone agro-écologique et sans phrase de synthèse finale est systématiquement pénalisée : la moitié des points porte sur la localisation et l'interprétation.
\item \textbf{Oublier la légende ou l'orientation du croquis.} Tout code couleur non expliqué, toute absence de flèche du nord ou de titre fait perdre des points de méthode, même si le fond est exact. Relisez la légende avant de rendre la copie.
\item \textbf{Se tromper dans les formules statistiques.} Le taux de variation se calcule sur la valeur de départ : $T=\dfrac{V_{2}-V_{1}}{V_{1}}\times 100$, et non sur $V_{2}$. Pour le diagramme circulaire, vérifiez toujours que la somme des angles vaut $360^{\circ}$. Écrivez l'unité (tonnes, \%, degrés) dans chaque résultat.
\item \textbf{Affirmer sans exemple camerounais précis.} Écrire « l'agriculture s'est modernisée » sans citer la SEMRY (riz, Yagoua), la SOSUCAM (canne, Nkoteng-Bandjock), la SODECAO (cacao) ou l'IRAD (recherche et semences) donne une copie vague. Un exemple localisé et chiffré (même en ordre de grandeur) vaut mieux que trois généralités.
\end{enumerate}
\end{attention}

\newpage

\coursec{Exercices}

\coursec{L'agriculture traditionnelle : de grandes améliorations}

\courssub{L'agriculture traditionnelle camerounaise : cadre et caractéristiques}

\begin{definition}[Agriculture traditionnelle]
Ensemble de pratiques culturales transmises de génération en génération, fondées sur le travail manuel, l'usage d'outils simples (houe, machette, daba), la polyculture vivrière, la jachère longue pour restaurer la fertilité, et une main-d'œuvre familiale. Elle domine encore dans les exploitations familiales (plus de 80\% des exploitations au Cameroun).
\end{definition}

\begin{propriete}[Les deux grands systèmes traditionnels]
\begin{itemize}
\item \textbf{Agriculture itinérante sur brûlis (ou culture itinérante sur brûlis) :} prédomine en zone forestière (Sud, Centre, Est, Sud-Ouest). Le paysan défriche une parcelle, brûle la végétation (apport de cendres = engrais potassique), cultive 1 à 3 ans, puis abandonne la parcelle à la jachère longue (10 à 20 ans autrefois) pour en ouvrir une autre plus loin. Exemple : culture du manioc, macabo, plantain, igname chez les Béti, Boulou, Bassa.
\item \textbf{Agriculture vivrière sédentaire (ou agriculture de case) :} prédomine en zone soudanienne et sahélienne (Nord, Extrême-Nord, Adamaoua). La pression foncière et la sécheresse contraignent à cultiver chaque année les mêmes champs autour du village. Fumure organique (déjections animales), rotation culturale (céréales / légumineuses), cultures associées (mil/sorgho + niébé + arachide). Exemple : mil, sorgho, maïs, riz pluvial, coton chez les Peuls, Mafa, Tupuri.
\end{itemize}
\end{propriete}

\begin{savaistu}[Le terroir villageois en zone forestière]
L'organisation en ceintures concentriques autour du village (jardins de case $\rightarrow$ champs annuels $\rightarrow$ jachères jeunes $\rightarrow$ jachères vieilles/forêt) est le modèle classique décrit par René Dumont et les géographes de l'ORSTOM (actuel IRD) dans les années 1960. La distance moyenne des champs au village est un indicateur direct de la pression foncière.
\end{savaistu}

\courssub{Un outillage amélioré}

\begin{definition}[Outillage amélioré]
Ensemble d'outils et d'équipements intermédiaires entre l'outil purement manuel (houe, daba, machette) et la mécanisation lourde (tracteur). Il comprend la \cle{traction animale} (charrue, hersse, semoir, charrette attelés à des bœufs/zébus/ânes) et la \cle{motorisation légère} (motoculteur, pulvérisateur motorisé, décortiqueuse, broyeur).
\end{definition}

\begin{experience}[Diffusion de la culture attelée au Nord-Cameroun]
\textbf{Contexte :} Dans les années 1970-1980, la SODECOTON (Société de Développement du Coton) et les missions agricoles (CFA, puis SAA) ont vulgarisé la charrue « Manga » et le semoir « Maya » auprès des cotonculteurs de l'Adamaoua, du Nord et de l'Extrême-Nord.
\textbf{Résultat :} La surface labourée par exploitation est passée de 0,5 ha (houe) à 2-3 ha (attelage). Le rendement du coton a gagné 30 à 50\% grâce au labour profond (meilleure infiltration, enracinement) et au semis en lignes (densité maîtrisée, désherbage facilité).
\textbf{Limite actuelle :} Le coût d'un attelage complet (2 bœufs + charrue + hersse + semoir) dépasse souvent 1 500 000 FCFA, hors de portée sans crédit. La mortalité bovine (péripneumonie, sècheresse) fragilise le système.
\end{experience}

\begin{aretenir}[Répartition spatiale de l'outillage (ordre de grandeur)]
\begin{itemize}
\item \textbf{Zone soudano-sahélienne (Nord, Extrême-Nord, Adamaoua) :} Traction animale dominante (coton, maïs, mil/sorgho). Environ 30 à 40\% des exploitations cotonnières équipées en attelage.
\item \textbf{Zone forestière (Centre, Sud, Est, Littoral, Sud-Ouest, Ouest, Nord-Ouest) :} Outillage manuel amélioré (machette, houe, daba, pulvérisateur à dos, motoculteur rare). Relief accidenté, parcelles petites et enclavées, essences forestières (souches) interdisant la charrue.
\item \textbf{Plaines alluviales (Yagoua, Mbo, Noun, Mbam) :} Motorisation légère (motopompes pour riz irrigué, motoculteurs, batteuses) via SEMRY, UNVDA, SODECOTON.
\end{itemize}
\end{aretenir}

\courssub{L'évolution des techniques culturales}

\begin{methode}[Principales améliorations techniques adoptées]
\begin{enumerate}
\item \textbf{Semences sélectionnées / variétés améliorées :} Maïs (CMS 8704, CMS 9015, variétés courtes cycles 90-100 jours), Riz (NERICA, IR 46, TOX 3145), Manioc (variétés TMS résistantes à la mosaïque, rendement 25-30 t/ha vs 8-10 t/ha local), Coton (variétés IRMA, Allen 333). Diffusées par IRAD, SODECOTON, SEMRY, UNVDA.
\item \textbf{Fertilisation :} Engrais minéraux (NPK 20-10-10, Urée 46\% N) sur coton, maïs, riz. Fumure organique (compost, fumier, engrais verts \emph{Mucuna}, \emph{Crotalaria}) promue par l'agroécologie (projets PADFA, PRODER). Dose type coton : 200 kg/ha NPK + 100 kg/ha Urée.
\item \textbf{Protection phytosanitaire :} Traitements insecticides (coton : 4 à 5 passages/an contre charançons, jassides ; maïs : traitement semences + 1 passage contre chenille légionnaire). Pulvérisateur à dos (15-20 L) généralisé.
\item \textbf{Techniques culturales :} Labour profond (20-25 cm) vs grattage (5-10 cm) ; semis en lignes (écartement 0,75-0,90 m) vs semis au hasard ; désherbage chimique (herbicides pré/post-levée) complétant le sarclage manuel ; rotation céréales/légumineuses (maïs/soja, coton/arachide).
\item \textbf{Gestion de l'eau :} Riz irrigué (plaine de Yagoua : 10 000 ha aménagés par SEMRY ; plaine de la Noun : barrage de Mapé, 7 000 ha potentiels) ; cultures de contre-saison (tomate, oignon, maïs) en bas-fonds.
\end{enumerate}
\end{methode}

\begin{attention}[Pièges fréquents au Bac]
\begin{itemize}
\item Ne pas confondre \emph{agriculture de seconde génération} (exploitations familiales modernisées, 5-50 ha, intrants, marché) et \emph{agro-industrie / plantations industrielles} (CDC, SOCAPALM, HEVECAM, PHP : centaines/milliers d'hectares, salariat, capital étranger).
\item La traction animale n'est \textbf{pas} de la mécanisation lourde (tracteur). Elle est un outillage amélioré.
\item « Jachère courte » ne signifie pas « absence de jachère ». En zone forestière, la jachère est passée de 15-20 ans à 3-5 ans, ce qui est insuffisant pour restaurer la fertilité $\rightarrow$ dégradation des sols (acidification, érosion).
\end{itemize}
\end{attention}

\courssub{L'adoption de nouvelles cultures}

\begin{propriete}[Chronologie et géographie des adoptions majeures]
\begin{itemize}
\item \textbf{Époque précoloniale (avant 1884) :} Manioc (Amérique du Sud via Portugais, XVIe s.), Maïs (Amérique, XVIe s.), Plantain (Asie du Sud-Est, ancien), Arachide (Amérique, XVIIIe s.), Piment, Tomate. Ces cultures sont devenues les bases de l'alimentation (manioc = 1ère culture vivrière en tonnage).
\item \textbf{Époque coloniale (1884-1960) :} \textbf{Cultures de rente imposées/encouragées :} Cacao (Brésil $\rightarrow$ Sud, Centre, Littoral, SW), Café arabica (Java $\rightarrow$ Ouest, NO), Café robusta (Congo $\rightarrow$ Centre, Sud), Palmier à huile (Sélection Deli $\rightarrow$ SW, Littoral), Hévéa (Asie $\rightarrow$ SW), Coton (USA/Égypte $\rightarrow$ Nord via CFDT/SODECOTON), Riz (Asie $\rightarrow$ Yagoua, Mbo).
\item \textbf{Depuis l'indépendance (1960-) :} Soja (Nord, Adamaoua - rotation coton), Niébé (partout), Patate douce, Taro, Nouvelles variétés de maïs/riz/manioc (IRAD/IITA), Ananas, Banane dessert (PHP, Boh Plantations), Pomme de terre (Ouest, Adamaoua - zones d'altitude).
\end{itemize}
\end{propriete}

\begin{exemplebox}[Le maïs : de culture marginale à culture stratégique]
Dans les années 1960, le maïs était secondaire (mil/sorgho au Nord, manioc/plantain au Sud). Grâce aux variétés composées (CMS) et hybrides (IRAD), à l'engrais et au labour attelé, la production est passée de $\sim$300 000 t (1970) à plus de \textbf{2,5 millions de tonnes (2022)} (source : MINADER/INS). Il est devenu la 1ère céréale produite et consommée (farine, pâte, boisson, aliment bétail). Zones clés : Ouest, Nord-Ouest, Adamaoua, Nord.
\end{exemplebox}

\courssub{L'agriculture itinérante sur brûlis et l'agriculture vivrière}

\begin{definition}[Jachère]
Période de repos d'une parcelle cultivée, pendant laquelle la végétation spontanée (herbes, arbustes, arbres) recolonise l'espace. Elle permet la reconstitution de la matière organique, de la structure du sol et de la biodiversité (auxiliaires de cultures). \textbf{Durée critique :} $<$ 5-7 ans en zone forestière $\rightarrow$ jachère « jeune » (herbacée/arbustive) inefficace $\rightarrow$ pourrissement des sols, invasion par \emph{Chromolaena odorata} (herbe de la mort), \emph{Imperata cylindrica}.
\end{definition}

\begin{experience}[Étude de cas : Village ewondo (Centre) - Données réelles]
\begin{center}
\begin{tabularx}{\linewidth}{|l|X|X|X|}\hline
\rowcolor{popblueL} Indicateur & Vers 1960 & Vers 1990 & Vers 2020 \\ \hline
Population village & 300 & 750 & 1 800 \\ \hline
Durée jachère (ans) & 15 & 8 & 4 \\ \hline
Rendement manioc (t/ha) & 10 & 8 & 6 \\ \hline
Distance champs/village (km) & 1 & 2,5 & 5 \\ \hline
\end{tabularx}
\end{center}
\textbf{Analyse :} Triplement de la population en 60 ans $\rightarrow$ raccourcissement drastique de la jachère (facteur 3,75) $\rightarrow$ baisse rendement (40\%) $\rightarrow$ allongement distance champs (facteur 5) $\rightarrow$ perte de temps de travail, fatigue, déforestation périphérique. Le système n'est plus en équilibre.
\end{experience}

\begin{aretenir}[Agriculture vivrière aujourd'hui]
Elle reste la base de la \cle{sécurité alimentaire} (autosuffisance relative en racines/tubercules/plantain, déficit en riz/blé/maïs/huile/poisson/lait). Elle s'intensifie : cultures associées (maïs/manioc/arachide), utilisation d'engrais sur manioc/plantain (Ouest), motoculture (basse vallée du Mbam), commercialisation du surplus (plantain, macabo, igname vers Douala/Yaoundé).
\end{aretenir}

\courssub{L'agriculture de seconde génération et bilan des améliorations}

\begin{definition}[Agriculture de seconde génération (2G)]
Exploitation familiale modernisée (5 à 50 ha), gérée comme une entreprise : utilisation d'intrants (semences sélectionnées, engrais, pesticides), outillage amélioré (attelage, motoculteur), main-d'œuvre salariée saisonnière, encadrement technique (SODECOTON, SEMRY, UNVDA, UCCAO, CAPEF), production orientée vers le marché (cultures de rente + vivrières commercialisées), revenu monétaire significatif. Appelée aussi « agriculture paysanne modernisée » ou « agriculture commerciale familiale ».
\end{definition}

\begin{exemplebox}[Exemples localisés d'agriculture 2G]
\begin{itemize}
\item \textbf{Bassin cotonnier (Nord, Extrême-Nord, Adamaoua) :} Paysans « cotonculteurs » encadrés par SODECOTON (crédit intrants, conseil, collecte). Surface moyenne 5-10 ha. Rotation coton/maïs/soja/arachide. Revenus : 500 000 - 2 000 000 FCFA/an.
\item \textbf{Plaine de Yagoua (Extrême-Nord) :} Riziculteurs SEMRY (aménagement hydro-agricole, 2 campagnes/an). Rendement 5-6 t/ha. Motorisation (motopompes, batteuses).
\item \textbf{Hauts-Plateaux de l'Ouest (Bafoussam, Dschang) :} Producteurs de maïs, haricot, pomme de terre, chou, carotte. Encadrement UCCAO, CAPEF. Utilisation engrais, fientes de poulet, motoculteurs. Spéculation sur marchés urbains.
\item \textbf{Zone cacaoyère (Centre, Sud, SW) :} « Cacaoyers de 2ème génération » : vergers rénovés (matériel sélectionné IRAD/IRCC), ombrage rationnel, fumure organique + engrais, traitement phytosanitaire. Rendement 800-1 200 kg/ha vs 300-400 kg/ha vergers vieillissants.
\end{itemize}
\end{exemplebox}

\begin{propriete}[Bilan des améliorations : Progrès et Limites]
\begin{center}
\begin{tabularx}{\linewidth}{|l|X|X|}\hline
\rowcolor{popblueL} \textbf{Dimension} & \textbf{Progrès constatés} & \textbf{Limites persistantes} \\ \hline
\textbf{Rendements} & Maïs : 1,5 $\rightarrow$ 3-4 t/ha (2G) ; Coton : 800 $\rightarrow$ 1 300 kg/ha ; Riz irrigué : 5-6 t/ha. & Écart huge avec potentiel (maïs 8-10 t/ha). Variabilité climatique (sécheresses). \\ \hline
\textbf{Production totale} & Autonomie en manioc, plantain, maïs, sorgho/mil. Export cacao, coton, banane, bois. & Déficit structurel : Riz (import 600 000 t/an), Blé (100\%), Huile, Poisson, Lait, Viande. \\ \hline
\textbf{Revenus paysans} & Hausse pour les 2G (cotton, riz, maraîchage). Monétarisation économies rurales. & Majorité exploitations < 2 ha, subsistance, pauvreté monétaire. Endettement (crédit intrants). \\ \hline
\textbf{Environnement} & Agroforesterie cacaoyère, rotation, engrais verts (projets). & Déforestation (extension cultures), érosion (labour pente), pollution (nitrates, pesticides), perte biodiversité. \\ \hline
\textbf{Infrastructures} & Quelques pistes aménagées (cotton belts). & Enclavement massif (pistes impraticables saison pluies), pertes post-récolte 20-30\%, stockage insuffisant. \\ \hline
\end{tabularx}
\end{center}
\end{propriete}

\courssub{Je vérifie mes connaissances}

\begin{exercice}[QCM : les notions clés]
Pour chaque question, une seule réponse est exacte. Recopie la lettre de la bonne réponse et justifie ton choix en une phrase.

\textbf{1.} La \cle{jachère} est :
\begin{itemize}
\item[a)] une période de repos des terres permettant la reconstitution de la fertilité du sol ;
\item[b)] une technique d'irrigation des rizières de la plaine de Yagoua ;
\item[c)] une plantation industrielle de palmiers à huile ;
\item[d)] un outil agricole à traction animale.
\end{itemize}

\textbf{2.} Le \cle{brûlis} consiste à :
\begin{itemize}
\item[a)] brûler les récoltes invendues sur le marché ;
\item[b)] mettre le feu à la végétation d'une parcelle défrichée pour fertiliser le sol par les cendres ;
\item[c)] chauffer le sol pour tuer les parasites avant les semis ;
\item[d)] brûler les herbes des pâturages pour faire paître le bétail.
\end{itemize}

\textbf{3.} L'\cle{agriculture vivrière} désigne :
\begin{itemize}
\item[a)] une agriculture destinée exclusivement à l'exportation ;
\item[b)] une agriculture mécanisée utilisant des tracteurs ;
\item[c)] une agriculture dont la production sert d'abord à nourrir la famille du producteur ;
\item[d)] une agriculture pratiquée uniquement en zone soudanienne.
\end{itemize}

\textbf{4.} L'\cle{agriculture de seconde génération} se caractérise par :
\begin{itemize}
\item[a)] l'abandon total des cultures vivrières au profit du cacao ;
\item[b)] l'association de techniques améliorées (intrants, matériel, encadrement) à une production orientée vers le marché ;
\item[c)] le retour aux pratiques d'avant la colonisation ;
\item[d)] la culture exclusive du riz dans les plaines inondables.
\end{itemize}
\end{exercice}

\begin{exercice}[Vrai ou faux ?]
Indique si chaque affirmation est vraie ou fausse, puis corrige les affirmations fausses.

\textbf{1.} La houe et la machette restent les seuls outils utilisés par les agriculteurs camerounais aujourd'hui.

\textbf{2.} La traction animale (culture attelée) s'est surtout diffusée dans les régions du Nord-Cameroun (Nord, Extrême-Nord, Adamaoua).

\textbf{3.} L'agriculture itinérante sur brûlis exige que la population reste fixe sur la même parcelle pendant vingt ans.

\textbf{4.} Le maïs, le manioc et le cacao sont des cultures introduites au Cameroun à l'époque coloniale ou précoloniale et adoptées par les paysans.

\textbf{5.} L'agriculture de seconde génération utilise des semences sélectionnées, des engrais et des pesticides.
\end{exercice}

\begin{exercice}[Définitions express]
Définis en deux ou trois lignes chacun des termes suivants, en donnant pour chacun un exemple localisé au Cameroun :
\begin{itemize}
\item agriculture itinérante sur brûlis ;
\item agriculture vivrière ;
\item culture de rente ;
\item agriculture de seconde génération.
\end{itemize}
\end{exercice}

\begin{exercice}[Calculs directs : rendements et production]
Un village de la plaine des Mbo (Ouest) a récolté les données suivantes pour le maïs :

\begin{center}
\begin{tabularx}{\linewidth}{|l|X|X|X|}\hline
\rowcolor{popblueL} Campagne & Surface cultivée (ha) & Production (tonnes) & Rendement (t/ha) \\ \hline
2018 & 120 & 180 & \ldots \\ \hline
2023 & 120 & 300 & \ldots \\ \hline
\end{tabularx}
\end{center}

\begin{enumerate}
\item Calcule le rendement du maïs (en tonnes par hectare) pour chacune des deux campagnes.
\item Calcule le taux d'accroissement de la production entre 2018 et 2023 (en \%).
\item À surface égale, à quoi peut-on attribuer cette hausse du rendement ? Cite deux améliorations techniques possibles.
\end{enumerate}
\end{exercice}

\courssub{Je m'entraîne}

\begin{exercice}[L'outillage agricole : du manuel à l'attelé]
Voici une liste d'outils et d'équipements observés dans les campagnes camerounaises : machette, houe, charrue à traction animale, tracteur, daba, pulvérisateur à dos, semoir, coupe-coupe.

\begin{enumerate}
\item Classe ces outils en trois catégories : outillage traditionnel manuel, outillage amélioré (traction animale ou motorisation légère), mécanisation lourde.
\item Pour chaque catégorie, indique une région du Cameroun où elle domine et justifie ce choix (relief, type de culture, revenu des paysans).
\item Explique en cinq lignes pourquoi le tracteur reste peu répandu dans les exploitations familiales camerounaises.
\end{enumerate}
\end{exercice}

\begin{exercice}[Construire un histogramme : la production de maïs]
La production de maïs d'un groupement de paysans de la région de l'Ouest évolue ainsi :

\begin{center}
\begin{tabularx}{\linewidth}{|l|X|X|X|X|X|}\hline
\rowcolor{popblueL} Campagne & 2019 & 2020 & 2021 & 2022 & 2023 \\ \hline
Production (tonnes) & 80 & 95 & 110 & 150 & 190 \\ \hline
\end{tabularx}
\end{center}

\begin{enumerate}
\item Construis un histogramme représentant cette évolution (axe horizontal : campagnes ; axe vertical : production en tonnes ; échelle conseillée : 1 cm pour 20 tonnes).
\item Calcule le taux d'accroissement global de la production entre 2019 et 2023.
\item Calcule la part de l'augmentation enregistrée entre 2021 et 2023 dans l'augmentation totale de la période.
\item Rédige un commentaire de cinq lignes : que révèle cette évolution des améliorations de l'agriculture traditionnelle (semences sélectionnées, engrais, encadrement) ?
\end{enumerate}
\end{exercice}

\begin{exercice}[Croquis : le terroir villageois en zone forestière]
À partir de la description suivante, extraite d'une photographie aérienne d'un village boulou du Sud-Cameroun, réalise un croquis de terroir villageois.

\emph{Description :} au centre, le village s'étire le long d'une route latéritique orientée est-ouest. Autour des cases, des jardins de case (plantains, arbres fruitiers). Plus loin, une ceinture de champs de manioc et de macabo en cours de culture. Au-delà, des parcelles de jachère arbustive reconnaissables à leurs jeunes arbres, puis la forêt dense qui occupe le reste de l'espace. Une plantation de cacao de 3 hectares se situe au nord du village, le long d'un ruisseau.

\begin{enumerate}
\item Réalise le croquis en respectant l'organisation en ceintures concentriques autour du village.
\item Place une flèche du nord, donne un titre au croquis et rédige une légende ordonnée (au moins cinq éléments).
\item Souligne par un trait de couleur la limite entre l'espace cultivé et la forêt.
\item En trois lignes, explique ce que cette organisation révèle du fonctionnement de l'agriculture itinérante sur brûlis.
\end{enumerate}
\end{exercice}

\begin{exercice}[Légender une carte : les grandes cultures du Cameroun]
Sur un fond de carte du Cameroun reproduit sur ta feuille (contour simplifié ci-dessous à recopier), localise et légende les zones de production suivantes :

\begin{itemize}
\item le coton (Nord, Extrême-Nord, Adamaoua) ;
\item le cacao (Centre, Sud, Littoral, Sud-Ouest) ;
\item le café arabica (Ouest, Nord-Ouest) et le café robusta (Centre, Sud) ;
\item le maïs (Ouest, Nord-Ouest, Adamaoua) ;
\item le riz (plaine de Yagoua dans l'Extrême-Nord, plaines des Mbo et de la Noun dans l'Ouest).
\end{itemize}

\begin{popfigure}
\begin{center}
\includegraphics[width=\linewidth,height=10cm,keepaspectratio]{N04/cmr_reg_muette.jpg}
\end{center}
\legende{Fond de carte du Cameroun à recopier et à légender (limites des régions en trait épais, départements en trait fin, fleuves en bleu). \\ Source du fond de carte : Wikimedia Commons, Flappiefh, CC BY-SA 4.0.}
\end{popfigure}

\begin{enumerate}
\item Reproduis le fond de carte et hachure ou colorie chaque zone de culture avec un symbole distinct.
\item Rédige une légende ordonnée (des cultures de rente vers les cultures vivrières) et donne un titre à ta carte.
\item Formule deux remarques sur la répartition spatiale des cultures (opposition Nord/Sud, rôle du climat).
\end{enumerate}
\end{exercice}

\courssub{Je me prépare au Bac}

\begin{exercice}[Question à réponse construite : les améliorations de l'agriculture traditionnelle]
\textbf{Document 1.} Photographie (décrite) : dans la plaine de l'Adamaoua, près de Ngaoundéré, un paysan guide une charrue tirée par deux zébus ; la parcelle, d'environ 2 hectares, est labourée en lignes régulières avant les semis de maïs.

\textbf{Document 2.} Évolution des rendements du coton au Cameroun (ordres de grandeur, sources : SODECOTON et INS) :

\begin{center}
\begin{tabularx}{\linewidth}{|l|X|X|X|}\hline
\rowcolor{popblueL} Période & Années 1970 & Années 1990 & Années 2010 \\ \hline
Rendement (kg/ha, ordre de grandeur) & 800 & 1 100 & 1 300 \\ \hline
\end{tabularx}
\end{center}

\textbf{Document 3.} Carte (décrite) : les cultures de rente du Cameroun — coton dans les trois régions septentrionales, cacao dans le Centre-Sud et le Littoral, café dans l'Ouest, palmier à huile et hévéa dans le Sud-Ouest et le Littoral, banane dans la plaine de la Njombé-Penja.

\textbf{Travail à faire :}
\begin{enumerate}
\item Décris le document 1 et dégage ce qu'il révèle de l'évolution de l'outillage agricole. (4 points)
\item Calcule le taux d'accroissement du rendement du coton entre les années 1970 et les années 2010, puis propose deux explications de cette hausse. (4 points)
\item À partir de l'étude des documents et de tes connaissances, montre, dans un développement rédigé et structuré (introduction, deux paragraphes, conclusion), que l'agriculture traditionnelle camerounaise a connu de grandes améliorations. (12 points)
\end{enumerate}
\end{exercice}

\begin{exercice}[Étude de cas : l'agriculture itinérante sur brûlis face à la croissance démographique]
Dans un village ewondo de la région du Centre, on a relevé les données suivantes sur trois générations :

\begin{center}
\begin{tabularx}{\linewidth}{|l|X|X|X|}\hline
\rowcolor{popblueL} Indicateur & Vers 1960 & Vers 1990 & Vers 2020 \\ \hline
Population du village (habitants) & 300 & 750 & 1 800 \\ \hline
Durée moyenne de la jachère (années) & 15 & 8 & 4 \\ \hline
Rendement du manioc (t/ha, ordre de grandeur) & 10 & 8 & 6 \\ \hline
Distance moyenne des champs au village (km) & 1 & 2,5 & 5 \\ \hline
\end{tabularx}
\end{center}

\begin{enumerate}
\item Définis l'agriculture itinérante sur brûlis et décris ses étapes successives. (4 points)
\item Calcule le taux d'accroissement de la population du village entre 1960 et 2020, puis le taux de baisse du rendement du manioc sur la même période. (4 points)
\item Établis la relation entre l'évolution de la population, la durée de la jachère et la baisse des rendements. (4 points)
\item Propose, en un paragraphe argumenté, trois améliorations techniques permettant de fixer l'agriculture sur des parcelles permanentes tout en maintenant les rendements (fumure organique, engrais, rotation, agroforesterie cacaoyère). (4 points)
\item Pourquoi dit-on que ce système traditionnel était autrefois « en équilibre avec le milieu » ? Cet équilibre est-il encore possible aujourd'hui ? Justifie ta réponse en quatre lignes. (4 points)
\end{enumerate}
\end{exercice}

\begin{exercice}[Dissertation guidée : l'agriculture de seconde génération]
\textbf{Sujet :} « L'agriculture de seconde génération est-elle la solution à l'alimentation des populations camerounaises ? »

\textbf{Consignes et plan guidé :}
\begin{enumerate}
\item \textbf{Introduction (4 points).} Amène le sujet à partir d'une situation concrète (par exemple : les files de camions de plantain et de maïs approvisionnant les marchés de Douala et de Yaoundé), définis les termes clés (\cle{agriculture de seconde génération}, \cle{sécurité alimentaire}) et pose une problématique claire. Annonce ton plan.
\item \textbf{Première partie (8 points).} Montre que l'agriculture de seconde génération constitue un progrès décisif : semences sélectionnées, engrais et traitements, traction animale et motorisation, encadrement (SODECOTON, SEMRY à Yagoua, UCCAO à Bafoussam), hausse des rendements et des revenus paysans, approvisionnement des villes. Appuie chaque argument sur au moins un exemple localisé.
\item \textbf{Deuxième partie (6 points).} Montre ses limites : coût élevé des intrants et endettement des paysans, dépendance aux importations (engrais, matériel), faible mécanisation, pistes rurales dégradées qui bloquent l'évacuation, concurrence des produits alimentaires importés (riz, huile), risques environnementaux. Conclus sur les conditions à réunir pour qu'elle nourrisse durablement le pays.
\item \textbf{Conclusion (2 points).} Réponds clairement à la problématique et ouvre sur une perspective (agro-industrie, souveraineté alimentaire, « import-substitution »).
\end{enumerate}

\emph{Barème indicatif : 20 points. Soin de la langue et organisation de la copie pris en compte.}
\end{exercice}

\newpage

\coursec{Corrigés des exercices}

\coursec{Leçon 4 : L'agriculture traditionnelle : de grandes améliorations}

\courssub{Exercices d'application et corrigés détaillés}

\begin{corrige}[Exercice 1 — QCM : les notions clés]
\textbf{1. Réponse a).} La jachère est une période de repos d'une parcelle pendant laquelle la végétation spontanée recolonise le sol et reconstitue sa fertilité (matière organique, structure). Les propositions b), c) et d) ne correspondent à aucune définition de la jachère.

\textbf{2. Réponse b).} Le brûlis consiste à incendier la végétation d'une parcelle défrichée : les cendres riches en potasse servent d'engrais naturel avant les semis. C'est l'étape centrale de l'agriculture itinérante sur brûlis.

\textbf{3. Réponse c).} L'agriculture vivrière produit d'abord pour la consommation de la famille du producteur (manioc, plantain, mil, maïs), le surplus éventuel étant vendu. Elle n'est ni réservée à l'exportation (a), ni mécanisée (b), ni limitée au Nord (d).

\textbf{4. Réponse b).} L'agriculture de seconde génération associe techniques améliorées (semences sélectionnées, engrais, attelage, encadrement) et orientation vers le marché, tout en restant une exploitation familiale. Elle n'abandonne pas les cultures vivrières (a), ne revient pas au passé (c) et ne se limite pas au riz (d).

\textbf{Points de vigilance :} au Bac, la justification en une phrase est exigée ; une lettre seule ne rapporte pas tous les points.
\end{corrige}

\begin{corrige}[Exercice 2 — Vrai ou faux ?]
\textbf{1. FAUX.} La houe et la machette restent très répandues, mais elles ne sont plus les seuls outils : la traction animale (charrue, semoir) domine au Nord, et la motorisation légère (motoculteurs, motopompes, pulvérisateurs) progresse dans les plaines aménagées (Yagoua, Mbo, Noun).

\textbf{2. VRAI.} La culture attelée s'est diffusée dans les régions septentrionales (Nord, Extrême-Nord, Adamaoua), portée par la SODECOTON : relief de plaines, cultures en lignes (coton, maïs) et présence de zébus la rendent possible.

\textbf{3. FAUX.} C'est l'inverse : l'agriculture itinérante repose sur le déplacement des champs. Après 1 à 3 ans de culture, la parcelle est abandonnée à la jachère et le paysan en ouvre une autre. C'est la jachère, pas la culture, qui durait 15 à 20 ans.

\textbf{4. VRAI.} Le maïs et le manioc sont arrivés d'Amérique à l'époque précoloniale (via les Portugais, XVIe siècle), le cacao a été introduit à l'époque coloniale ; tous ont été adoptés et appropriés par les paysans camerounais.

\textbf{5. VRAI.} L'agriculture de seconde génération se définit précisément par l'usage d'intrants (semences sélectionnées, engrais NPK et urée, pesticides), d'un outillage amélioré et d'un encadrement technique.

\textbf{Points de vigilance :} une affirmation fausse non corrigée ne vaut que la moitié des points ; toujours reformuler la phrase corrigée.
\end{corrige}

\begin{corrige}[Exercice 3 — Définitions express]
\textbf{Agriculture itinérante sur brûlis :} système de culture de la zone forestière dans lequel le paysan défriche une parcelle, brûle la végétation (les cendres fertilisent le sol), cultive 1 à 3 ans puis abandonne la parcelle à une longue jachère pour en ouvrir une autre. \emph{Exemple :} culture du manioc et du macabo chez les Boulou du Sud.

\textbf{Agriculture vivrière :} agriculture dont la production sert en priorité à nourrir la famille du producteur, le surplus étant vendu sur les marchés locaux. \emph{Exemple :} le mil et le sorgho cultivés autour des villages de l'Extrême-Nord (pays mafa).

\textbf{Culture de rente :} culture destinée à la vente, sur le marché national ou à l'exportation, qui procure un revenu monétaire au producteur. \emph{Exemple :} le coton dans le bassin de Garoua (SODECOTON) ou le cacao dans le Centre.

\textbf{Agriculture de seconde génération :} exploitation familiale modernisée (5 à 50 ha) utilisant intrants, outillage amélioré et encadrement technique, et produisant pour le marché. \emph{Exemple :} les riziculteurs de la plaine de Yagoua encadrés par la SEMRY.

\textbf{Points de vigilance :} chaque définition doit comporter l'exemple localisé ; une définition sans exemple camerounais est incomplète.
\end{corrige}

\begin{corrige}[Exercice 4 — Calculs directs : rendements et production]
\textbf{1. Rendements.} Le rendement est le rapport de la production à la surface :
\[ R_{2018} = \frac{180}{120} = 1{,}5 \text{ t/ha} \qquad R_{2023} = \frac{300}{120} = 2{,}5 \text{ t/ha}. \]

\textbf{2. Taux d'accroissement de la production :}
\[ T = \frac{300 - 180}{180} \times 100 = \frac{120}{180} \times 100 \approx 66{,}7\,\%. \]
La production a augmenté d'environ \textbf{66,7\,\%} entre 2018 et 2023.

\textbf{3. Explications possibles (à surface égale) :} l'adoption de semences sélectionnées de maïs (variétés CMS à cycle court diffusées par l'IRAD) et l'apport d'engrais minéraux (NPK puis urée). On peut aussi citer le semis en lignes, le désherbage amélioré ou l'encadrement par une coopérative.

\textbf{Points de vigilance :} ne pas confondre accroissement de la production (66,7\,\%) et accroissement du rendement, qui vaut ici $\frac{2{,}5-1{,}5}{1{,}5}\times 100 \approx 66{,}7\,\%$ — identique uniquement parce que la surface est constante ; le préciser montre la rigueur du raisonnement.
\end{corrige}

\begin{corrige}[Exercice 5 — L'outillage agricole : du manuel à l'attelé]
\textbf{1. Classement :}
\begin{center}
\begin{tabularx}{\linewidth}{|l|X|}\hline
\rowcolor{popblueL} Catégorie & Outils \\ \hline
Outillage traditionnel manuel & machette, houe, daba, coupe-coupe \\ \hline
Outillage amélioré & charrue à traction animale, semoir, pulvérisateur à dos \\ \hline
Mécanisation lourde & tracteur \\ \hline
\end{tabularx}
\end{center}

\textbf{2. Régions dominantes :}
\begin{itemize}
\item \emph{Outillage manuel :} zone forestière (Centre, Sud, Est) — relief accidenté, parcelles petites et souches d'arbres qui interdisent la charrue ; cultures (manioc, plantain) peu exigeantes en labour.
\item \emph{Outillage amélioré :} Nord, Extrême-Nord, Adamaoua — plaines soudaniennes dégagées, présence de zébus, cultures en lignes (coton, maïs) vulgarisées par la SODECOTON.
\item \emph{Mécanisation lourde :} rare ; quelques plaines aménagées et périmètres de sociétés (SEMRY à Yagoua, UNVDA dans la Noun) et plantations agro-industrielles (PHP, SOCAPALM).
\end{itemize}

\textbf{3. Faible diffusion du tracteur (5 lignes) :} le tracteur coûte très cher à l'achat (plusieurs millions de FCFA) et à l'entretien (carburant, pièces détachées importées) ; les exploitations familiales sont trop petites (souvent moins de 2 ha) pour le rentabiliser ; le relief forestier accidenté, les souches et l'enclavement des pistes limitent son usage ; enfin, l'absence de crédit agricole accessible et de services de location empêche sa diffusion.

\textbf{Points de vigilance :} ne pas classer la charrue attelée dans la mécanisation lourde : la traction animale relève de l'outillage amélioré.
\end{corrige}

\begin{corrige}[Exercice 6 — Construire un histogramme : la production de maïs]
\textbf{1. Histogramme attendu} (modèle à l'échelle 1 cm pour 20 t ; hauteurs : 4 cm ; 4,75 cm ; 5,5 cm ; 7,5 cm ; 9,5 cm) :

\begin{popfigure}
\begin{center}
\begin{tikzpicture}[scale=0.9]
\draw[->,thick,popdark] (0,0) -- (6.6,0) node[below left,font=\small]{Campagnes};
\draw[->,thick,popdark] (0,0) -- (0,5.4) node[above,font=\small]{Production (t)};
\foreach \y/\lab in {1/40,2/80,3/120,4/160,5/200}{\draw[popline] (0,\y) -- (6.4,\y); \node[left,font=\small] at (0,\y) {\lab};}
\fill[popblue!70] (0.3,0) rectangle (1.1,2);
\fill[popblue!70] (1.5,0) rectangle (2.3,2.375);
\fill[popblue!70] (2.7,0) rectangle (3.5,2.75);
\fill[popblue!70] (3.9,0) rectangle (4.7,3.75);
\fill[poporange!80] (5.1,0) rectangle (5.9,4.75);
\node[below,font=\small] at (0.7,0) {2019};
\node[below,font=\small] at (1.9,0) {2020};
\node[below,font=\small] at (3.1,0) {2021};
\node[below,font=\small] at (4.3,0) {2022};
\node[below,font=\small] at (5.5,0) {2023};
\node[above,font=\small] at (0.7,2) {80};
\node[above,font=\small] at (1.9,2.375) {95};
\node[above,font=\small] at (3.1,2.75) {110};
\node[above,font=\small] at (4.3,3.75) {150};
\node[above,font=\small] at (5.5,4.75) {190};
\end{tikzpicture}
\end{center}
\legende{Histogramme de la production de maïs du groupement (2019-2023). Échelle : 1 cm pour 40 t sur le graphique modèle ; sur ta copie, 1 cm pour 20 t.}
\end{popfigure}

\textbf{2. Taux d'accroissement global 2019-2023 :}
\[ T = \frac{190 - 80}{80} \times 100 = \frac{110}{80} \times 100 = 137{,}5\,\%. \]

\textbf{3. Part de l'augmentation 2021-2023 dans l'augmentation totale :}
\[ \text{Augmentation 2021-2023} = 190 - 110 = 80 \text{ t} \;;\quad \frac{80}{110} \times 100 \approx 72{,}7\,\%. \]
Près des trois quarts de la hausse totale se concentrent sur les deux dernières campagnes.

\textbf{4. Commentaire (5 lignes) :} la production du groupement a plus que doublé en cinq campagnes (+137,5\,\%), avec une nette accélération après 2021. Cette évolution traduit les améliorations de l'agriculture traditionnelle : adoption de semences sélectionnées à cycle court, usage croissant d'engrais (NPK, urée, fientes de poulet sur les Hauts-Plateaux de l'Ouest), semis en lignes et encadrement par les coopératives (type UCCAO). L'accélération récente suggère que ces techniques se diffusent et se combinent, transformant progressivement une agriculture de subsistance en agriculture commerciale familiale.

\textbf{Points de vigilance :} axes légendés avec unités, barres de largeur égale et accolées, titre obligatoire ; un diagramme sans titre ni échelle perd des points.
\end{corrige}

\begin{corrige}[Exercice 7 — Croquis : le terroir villageois en zone forestière]
\textbf{1-3. Modèle de croquis attendu :}

\begin{popfigure}
\begin{center}
\begin{tikzpicture}[scale=0.95]
% Forêt dense
\fill[popgreen!25] (0,0) rectangle (10,7);
% Jachère arbustive
\fill[popgreen!45] (2.2,1.4) rectangle (7.8,5.8);
% Champs en culture
\fill[popgold!45] (3.2,2.2) rectangle (6.8,5.0);
% Jardins de case
\fill[poporange!45] (4.0,2.9) rectangle (6.0,4.3);
% Route
\draw[very thick,popdark,dashed] (0,3.6) -- (10,3.6);
\node[font=\small,popdark] at (9.2,3.85) {route latéritique};
% Village
\foreach \x in {4.2,4.7,5.2,5.7}{\draw[fill=popink] (\x,3.45) rectangle (\x+0.25,3.7);}
\node[font=\small] at (5.0,3.15) {village};
% Cacaoyère au nord le long du ruisseau
\draw[very thick,popblue] (2.4,7) .. controls (3.2,6.0) and (3.4,5.4) .. (3.6,4.6);
\node[font=\small,popblue] at (2.2,5.6) {ruisseau};
\fill[poppurple!40] (3.8,4.8) rectangle (5.4,5.7);
\node[font=\small] at (4.6,5.25) {cacaoyère (3 ha)};
% Limite espace cultivé / forêt (trait de couleur)
\draw[very thick,poppink] (2.2,1.4) rectangle (7.8,5.8);
% Nord
\draw[->,thick,popink] (9.4,5.9) -- (9.4,6.7); \node[above,font=\small] at (9.4,6.7) {N};
% Labels ceintures
\node[font=\small] at (5.0,4.55) {jardins de case};
\node[font=\small] at (5.0,2.5) {champs (manioc, macabo)};
\node[font=\small] at (5.0,1.75) {jachère arbustive};
\node[font=\small] at (8.6,6.4) {forêt dense};
\end{tikzpicture}
\end{center}
\legende{Croquis : le terroir villageois en zone forestière (Sud-Cameroun). Légende ordonnée : 1. village et cases ; 2. route latéritique ; 3. jardins de case ; 4. champs en culture (manioc, macabo) ; 5. jachère arbustive ; 6. forêt dense ; 7. cacaoyère et ruisseau ; 8. limite espace cultivé/forêt (trait rose).}
\end{popfigure}

\textbf{Barème indicatif :} organisation en ceintures concentriques respectée (4 pts) ; flèche du nord et titre (2 pts) ; légende ordonnée d'au moins cinq éléments (3 pts) ; limite cultivé/forêt soulignée (1 pt).

\textbf{4. Explication (3 lignes) :} l'organisation en ceintures concentriques traduit la logique de l'agriculture itinérante sur brûlis : les cultures les plus soignées sont proches des cases, les champs annuels occupent les parcelles fraîchement défrichées, et les jachères marquent les parcelles abandonnées en cours de régénération avant retour à la forêt. L'espace est ainsi parcouru par rotation des champs plutôt que par rotation des cultures.

\textbf{Points de vigilance :} un croquis n'est pas une carte : pas de coordonnées, mais un titre, une flèche du nord et une légende \emph{numérotée et ordonnée} (du village vers la forêt) sont exigés.
\end{corrige}

\begin{corrige}[Exercice 8 — Légender une carte : les grandes cultures du Cameroun]
\textbf{1-2. Modèle de carte légendée :}

\begin{popfigure}
\begin{center}
\includegraphics[width=\linewidth,height=9.5cm,keepaspectratio]{N04/cmr_koppen.png}
\end{center}
\legende{Carte : « Les grandes cultures du Cameroun : une répartition commandée par le climat ». Fond de carte : climats du Cameroun selon Köppen (légende en anglais) : climats tropicaux humides (bleus) au Sud et sur le littoral, climat de savane (bleu clair) au centre et au nord, climats arides (orange et rouge) à l'Extrême-Nord, climat plus frais (jaune) sur les hautes terres de l'Ouest. Le croquis attendu y superpose, du sud vers le nord, les cultures de rente (1 à 4) puis les cultures vivrières (5 et 6) en fonction de ces zones climatiques. \\ Source : Wikimedia Commons, Beck et al., CC BY 4.0.}
\end{popfigure}

\textbf{Barème indicatif :} localisation correcte des cinq ensembles (5 pts) ; symboles distincts et légende ordonnée (3 pts) ; titre et orientation (2 pts).

\textbf{3. Deux remarques :}
\begin{itemize}
\item \emph{Opposition Nord/Sud :} les cultures de rente septentrionales sont dominées par le coton en zone soudanienne sèche, tandis que le Sud forestier et humide concentre cacao, café, palmier à huile et hévéa.
\item \emph{Rôle du climat :} la répartition suit les gradients climatiques — coton et mil sous climat soudanien à longue saison sèche, cacao et café robusta sous climat équatorial humide, café arabica et maïs sur les hautes terres plus fraîches de l'Ouest ; le riz occupe les plaines inondables aménagées (Yagoua, Mbo, Noun).
\end{itemize}

\textbf{Points de vigilance :} la légende doit être \emph{ordonnée} (ici : cultures de rente puis vivrières) et chaque symbole doit être distinct ; ne jamais écrire les noms des cultures en travers de la carte sans symbole.
\end{corrige}

\begin{corrige}[Exercice 9 — Question à réponse construite : les améliorations de l'agriculture traditionnelle]
\textbf{1. Description du document 1 (4 pts).} La photographie montre un paysan de la plaine de l'Adamaoua conduisant une charrue tirée par deux zébus sur une parcelle d'environ 2 ha labourée en lignes régulières avant les semis de maïs. Elle révèle le passage de l'outillage purement manuel (houe, daba) à la \cle{traction animale}, un outillage amélioré qui permet de labourer plus profondément, de semer en lignes et de cultiver des surfaces trois à cinq fois plus grandes. \emph{Barème :} description factuelle (2 pts) ; mise en évidence de l'amélioration de l'outillage (2 pts).

\textbf{2. Calcul et explications (4 pts).}
\[ T = \frac{1\,300 - 800}{800} \times 100 = \frac{500}{800} \times 100 = 62{,}5\,\%. \]
Le rendement du coton a progressé d'environ \textbf{62,5\,\%} entre les années 1970 et les années 2010. Explications (deux attendues) : diffusion de semences sélectionnées (variétés IRMA) ; usage d'engrais (NPK + urée) et de traitements insecticides vulgarisés et fournis à crédit par la SODECOTON ; généralisation du labour attelé et des semis en lignes. \emph{Barème :} calcul exact avec formule (2 pts) ; deux explications localisées (2 pts).

\textbf{3. Développement structuré (12 pts).} \emph{Éléments attendus :}

\textbf{Introduction :} amorce (les campagnes camerounaises ne sont plus celles de la seule houe), définition de l'agriculture traditionnelle, problématique (en quoi a-t-elle été transformée ?), annonce du plan.

\textbf{Paragraphe 1 — Un outillage et des techniques modernisés :} traction animale au Nord (charrue « Manga », semoir « Maya » diffusés par la SODECOTON), motorisation légère dans les plaines aménagées (motopompes et batteuses de la SEMRY à Yagoua), semences sélectionnées (maïs CMS, manioc TMS, coton IRMA), engrais et traitements phytosanitaires, semis en lignes et rotations ; conséquence : hausse des rendements (coton +62,5\,\% ; riz irrigué 5-6 t/ha).

\textbf{Paragraphe 2 — De nouvelles cultures et une agriculture tournée vers le marché :} adoption historique du maïs, du manioc, puis des cultures de rente (cacao, café, coton, riz — document 3) ; essor de l'\cle{agriculture de seconde génération} (cotonculteurs du Nord, riziculteurs de Yagoua, maraîchers des Hauts-Plateaux de l'Ouest encadrés par l'UCCAO) qui approvisionne les villes et procure des revenus monétaires.

\textbf{Conclusion :} bilan — la tradition n'a pas disparu mais s'est profondément améliorée ; ouverture sur les limites (enclavement, dépendance aux intrants importés).

\emph{Barème indicatif :} introduction (2 pts) ; paragraphe 1 (4 pts) ; paragraphe 2 (4 pts) ; conclusion (1 pt) ; langue et organisation (1 pt).

\textbf{Points de vigilance :} citer les documents explicitement (« le document 2 montre que... ») ; chaque argument doit s'appuyer sur un exemple localisé ; interdiction de recopier les documents sans analyse.
\end{corrige}

\begin{corrige}[Exercice 10 — Étude de cas : l'agriculture itinérante sur brûlis face à la croissance démographique]
\textbf{1. Définition et étapes (4 pts).} L'agriculture itinérante sur brûlis est un système de culture de la zone forestière dans lequel les champs se déplacent au lieu que les cultures se succèdent sur le même champ. Étapes : (i) choix et défrichement d'une parcelle forestière ou de jachère vieille ; (ii) brûlis de la végétation, dont les cendres fertilisent le sol ; (iii) culture pendant 1 à 3 ans (manioc, macabo, plantain) ; (iv) abandon de la parcelle à la jachère longue (15-20 ans autrefois) et ouverture d'un nouveau champ plus loin.

\textbf{2. Calculs (4 pts).}
\[ T_{\text{population}} = \frac{1\,800 - 300}{300} \times 100 = \frac{1\,500}{300} \times 100 = 500\,\%. \]
\[ T_{\text{rendement}} = \frac{6 - 10}{10} \times 100 = -40\,\% \quad (\text{baisse de } 40\,\%). \]
La population a été multipliée par six (+500\,\%) pendant que le rendement du manioc chutait de 40\,\%.

\textbf{3. Relation entre les évolutions (4 pts).} La multiplication de la population par six exige davantage de terres cultivées chaque année : les parcelles sont remises en culture avant la fin de leur régénération, et la jachère passe de 15 à 4 ans. Or une jachère de moins de 5-7 ans reste herbacée ou arbustive et ne reconstitue ni la matière organique ni la structure du sol ; la fertilité s'épuise, les adventices envahissantes (\emph{Chromolaena odorata}) s'installent, et les rendements baissent (-40\,\%). Pour compenser, les paysans défrichent plus loin (champs à 5 km), ce qui accroît la fatigue et la déforestation : le système est sorti de son équilibre.

\textbf{4. Trois améliorations pour fixer l'agriculture (4 pts).} Premièrement, la fumure organique (compost, fumier, engrais verts comme \emph{Mucuna}) complétée par des engrais minéraux permet de restaurer la fertilité sans jachère longue. Deuxièmement, la rotation céréales/légumineuses (maïs/arachide ou niébé) et les cultures associées limitent l'épuisement du sol et les bioagresseurs. Troisièmement, l'agroforesterie cacaoyère — associer un verger de cacaoyers sélectionnés aux cultures vivrières — pérennise la couverture du sol, procure un revenu de rente et fixe durablement l'exploitation sur la même parcelle.

\textbf{5. Un équilibre autrefois, et aujourd'hui ? (4 pts).} On disait le système « en équilibre avec le milieu » parce que la faible densité de population laissait à chaque parcelle une jachère de 15 à 20 ans, suffisante pour régénérer la forêt et la fertilité : le prélèvement restait inférieur à la capacité de renouvellement du milieu. Cet équilibre n'est plus possible spontanément aujourd'hui, car la croissance démographique (facteur 6 en 60 ans) impose des jachères trop courtes ; il ne peut être retrouvé que par l'intensification (intrants, rotations, agroforesterie) plutôt que par l'extension spatiale.

\textbf{Points de vigilance :} dans le calcul du taux de baisse, le signe négatif doit être interprété (« baisse de 40\,\% ») ; la relation de causalité demandée à la question 3 doit enchaîner explicitement : population $\rightarrow$ jachère $\rightarrow$ fertilité $\rightarrow$ rendement.
\end{corrige}

\begin{corrige}[Exercice 11 — Dissertation guidée : l'agriculture de seconde génération]
\textbf{Corrigé-type (plan rédigé attendu).}

\textbf{Introduction (4 pts).} \emph{Amorce :} chaque matin, des camions de plantain, de maïs et de tomates en provenance de l'Ouest, du Centre et de l'Extrême-Nord approvisionnent les marchés de Douala et de Yaoundé : derrière eux se profile une paysannerie transformée. \emph{Définitions :} l'\cle{agriculture de seconde génération} est une exploitation familiale modernisée (5 à 50 ha) utilisant intrants, outillage amélioré et encadrement pour produire en vue du marché ; la \cle{sécurité alimentaire} désigne l'accès de tous, à tout moment, à une nourriture suffisante et saine. \emph{Problématique :} cette agriculture modernisée suffit-elle à nourrir durablement les populations camerounaises ? \emph{Plan :} nous montrerons d'abord qu'elle constitue un progrès décisif, puis qu'elle bute sur des limites qui conditionnent sa capacité à assurer la souveraineté alimentaire.

\textbf{Première partie — Un progrès décisif (8 pts).}
\begin{itemize}
\item \emph{Semences sélectionnées et intrants :} maïs CMS à cycle court (IRAD), manioc TMS résistant à la mosaïque, riz NERICA ; engrais NPK et urée ; rendements du maïs passés d'environ 1,5 à 3-4 t/ha.
\item \emph{Outillage amélioré :} traction animale dans le bassin cotonnier (charrue « Manga »), motorisation légère à Yagoua (motopompes, batteuses de la SEMRY, deux campagnes de riz par an, 5-6 t/ha).
\item \emph{Encadrement :} SODECOTON (crédit intrants, collecte du coton), SEMRY (plaine de Yagoua), UNVDA (plaine de la Noun), UCCAO (Bafoussam) pour le maïs et le maraîchage des Hauts-Plateaux.
\item \emph{Effets :} hausse des revenus paysans (500 000 à 2 000 000 FCFA/an pour les cotonculteurs), monétarisation des campagnes, approvisionnement massif des villes en maïs, plantain, macabo, légumes.
\end{itemize}

\textbf{Deuxième partie — Des limites réelles (6 pts).}
\begin{itemize}
\item \emph{Coût et dépendance :} intrants et matériel largement importés et chers ; endettement des paysans lié au crédit intrants ; attelage complet supérieur à 1 500 000 FCFA.
\item \emph{Infrastructures :} pistes rurales impraticables en saison des pluies, pertes post-récolte de 20 à 30\,\%, stockage insuffisant.
\item \emph{Concurrence et déficits :} riz importé (environ 600 000 t/an, ordre de grandeur selon les années), blé importé à 100\,\%, huile, poisson et lait déficitaires ; les produits importés concurrencent les productions locales.
\item \emph{Environnement :} déforestation, érosion des parcelles labourées en pente, risques liés aux pesticides. \emph{Conditions :} crédit agricole accessible, recherche et vulgarisation (IRAD), désenclavement, transformation locale.
\end{itemize}

\textbf{Conclusion (2 pts).} L'agriculture de seconde génération est une étape indispensable — elle a déjà fait du Cameroun un pays quasi autosuffisant en racines, tubercules, plantain et maïs — mais elle ne suffit pas seule : sans infrastructures, crédit et transformation, le pays reste dépendant des importations. \emph{Ouverture :} la voie vers la souveraineté alimentaire passe par l'articulation de cette paysannerie modernisée avec l'agro-industrie et une politique volontariste de substitution aux importations (« import-substitution »).

\textbf{Barème indicatif :} introduction (4 pts) ; première partie (8 pts : 4 arguments $\times$ 2 pts, exemple localisé exigé pour chacun) ; deuxième partie (6 pts) ; conclusion (2 pts). Langue et organisation intégrées à la notation.

\textbf{Points de vigilance :} une problématique doit être une question explicite ; chaque argument sans exemple localisé (région, société d'encadrement, chiffre) perd la moitié de sa valeur ; éviter de confondre agriculture de seconde génération et plantations agro-industrielles (SOCAPALM, PHP), hors sujet ici.
\end{corrige}

\newpage

\coursec{Fiche bilan}

\begin{popfigure}
\begin{tikzpicture}[mindmap, grow cyclic, every node/.style={concept, font=\small},
  root concept/.append style={concept color=popblue, font=\bfseries, text=white, minimum size=3.2cm},
  level 1 concept/.append style={level distance=4.2cm, sibling angle=60, font=\footnotesize, text=white},
  level 2 concept/.append style={level distance=2.6cm, font=\scriptsize, text=popdark}]
\node[root concept]{Agriculture traditionnelle : de grandes améliorations}
  child[concept color=popgreen]{ node {Cadre et caractéristiques}
    child { node[concept color=popgreenL]{Petites exploitations familiales} }
    child { node[concept color=popgreenL]{Main-d'œuvre familiale, faible capital} }
  }
  child[concept color=poporange]{ node {Outillage amélioré}
    child { node[concept color=poporangeL]{Houe, machette $\rightarrow$ charrue, tracteur} }
    child { node[concept color=poporangeL]{Traction animale au Nord} }
  }
  child[concept color=poppurple]{ node {Techniques culturales}
    child { node[concept color=poppurpleL]{Irrigation, drainage} }
    child { node[concept color=poppurpleL]{Engrais, pesticides, semences} }
    child { node[concept color=poppurpleL]{Jachère raccourcie} }
  }
  child[concept color=popgold]{ node {Nouvelles cultures}
    child { node[concept color=popgoldL]{Cultures de rente : cacao, café, coton} }
    child { node[concept color=popgoldL]{Maraîchage, riz, maïs améliorés} }
  }
  child[concept color=poppink]{ node {Itinérante sur brûlis et vivrière}
    child { node[concept color=poppinkL]{Défriche, brûlis, rotation des parcelles} }
    child { node[concept color=poppinkL]{Autoconsommation d'abord} }
  }
  child[concept color=popblue!70]{ node {Seconde génération}
    child { node[concept color=popblueL]{Mécanisation, intrants, crédit} }
    child { node[concept color=popblueL]{Agropastoralisme, intégration au marché} }
  };
\end{tikzpicture}
\legende{Carte mentale de la leçon : les six branches à retenir (croquis schématique, non normé).}
\end{popfigure}

\begin{bilan}
\textbf{Définitions clés}
\begin{itemize}
  \item \cle{Agriculture traditionnelle} : agriculture de subsistance, à petite échelle, reposant sur la main-d'œuvre familiale, un outillage simple et de faibles rendements.
  \item \cle{Agriculture itinérante sur brûlis} : système où l'on défriche, brûle la végétation (la cendre fertilise le sol), cultive quelques années, puis abandonne la parcelle (\emph{jachère}) pour une autre.
  \item \cle{Agriculture vivrière} : production destinée d'abord à l'autoconsommation de la famille (manioc, maïs, plantain, igname, mil).
  \item \cle{Agriculture de seconde génération} : agriculture modernisée (mécanisation, intrants, sélection végétale, crédit agricole, encadrement) visant la productivité et le marché.
\end{itemize}

\textbf{Repères à connaître par cœur}
\begin{center}
\begin{tabularx}{\linewidth}{|l|X|}\hline
\rowcolor{popblueL} \textbf{Notion} & \textbf{Repère camerounais} \\ \hline
Part de l'agriculture & Environ 15 à 20 \% du PIB ; emploie plus de la moitié de la population active (ordre de grandeur). \\ \hline
Traction animale & Surtout dans le Grand Nord (Adamaoua, Nord, Extrême-Nord) : charrues tirées par des bœufs. \\ \hline
Cultures de rente adoptées & Cacao (Centre, Sud, Littoral), café (Ouest, Nord-Ouest), coton (Nord), banane, palmier à huile. \\ \hline
Encadrement & SODECOTON (coton), SEMRY (riz, plaine de Yagoua), CDC et PHP (agro-industrie). \\ \hline
Limites du brûlis & Déforestation, érosion, baisse de fertilité quand la jachère se raccourcit. \\ \hline
\end{tabularx}
\end{center}

\textbf{Méthodes express}
\begin{itemize}
  \item \emph{Décrire une carte agricole} : localiser (région, ville repère) $\rightarrow$ décrire (cultures, zones) $\rightarrow$ expliquer (climat, sols, encadrement, marchés).
  \item \emph{Construire un croquis} : titre, contour simplifié du Cameroun, zones colorées, flèche du nord, légende numérotée.
  \item \emph{Commenter un tableau} : donner les ordres de grandeur, calculer parts en \% et évolutions, dégager une idée directrice.
\end{itemize}
\end{bilan}

\begin{aretenir}[Checklist avant le Bac]
\begin{itemize}
  \item[$\square$] Je sais définir agriculture traditionnelle, vivrière, itinérante sur brûlis et de seconde génération.
  \item[$\square$] Je cite au moins trois améliorations de l'outillage (charrue, traction animale, tracteur).
  \item[$\square$] J'explique le rôle des intrants : engrais, pesticides, semences améliorées.
  \item[$\square$] Je localise les grandes cultures de rente sur une carte du Cameroun.
  \item[$\square$] Je décris les étapes du cycle défriche--brûlis--culture--jachère.
  \item[$\square$] Je connais les limites écologiques de l'agriculture sur brûlis.
  \item[$\square$] Je cite deux organismes d'encadrement (SODECOTON, SEMRY) et leur rôle.
  \item[$\square$] Je distingue autoconsommation et commercialisation dans l'agriculture vivrière.
  \item[$\square$] Je sais construire un croquis légendé avec titre, nord et échelle indicative.
  \item[$\square$] Je relie les améliorations agricoles à l'intégration nationale (marchés, routes, villes).
\end{itemize}
\end{aretenir}

\findecours

\end{document}
